% !TeX program = xelatex
% !TeX encoding = UTF-8
% ============================================================================
%  LUẬN VĂN THẠC SĨ
%  Phát hiện và phân đoạn bệnh lá sầu riêng bằng học sâu
%  kết hợp AI khả diễn giải và gán nhãn giả
%
%  ⚠ BẮT BUỘC BIÊN DỊCH BẰNG XeLaTeX — KHÔNG DÙNG pdfLaTeX
%
%  Dòng "% !TeX program = xelatex" ở đầu file giúp Overleaf tự chọn đúng
%  trình biên dịch. Nếu Overleaf vẫn báo lỗi fontspec, hãy đặt thủ công:
%      Menu (góc trên bên trái) → Compiler → XeLaTeX → Recompile
%
%  Lý do: tài liệu dùng fontspec + polyglossia để hiển thị tiếng Việt,
%  hai gói này chỉ chạy được trên XeLaTeX hoặc LuaLaTeX.
%
%  Xem hướng dẫn đầy đủ tại README.md
% ============================================================================
% Cỡ chữ 13pt được đặt qua gói `fontsize` trong preamble.tex, KHÔNG qua tuỳ chọn
% của documentclass: gói extsizes chỉ hỗ trợ 8/9/10/11/12/14/17/20pt — truyền
% 13pt vào sẽ bị âm thầm bỏ qua và tài liệu tụt về 10pt.
\documentclass[12pt, a4paper, oneside]{report}

% ============================================================================
%  preamble.tex — toàn bộ cấu hình định dạng của luận văn
%  Biên dịch: XeLaTeX (Overleaf: Menu → Compiler → XeLaTeX)
%  Thư mục tham chiếu: biblatex + biber (Overleaf tự nhận diện)
%
%  Nếu trường có template riêng, chỉ cần thay THẾ file này — nội dung các
%  chương trong chapters/ không phụ thuộc vào định dạng.
% ============================================================================

% ---------------------------------------------------------------------------
% 1. Trang giấy và lề
% ---------------------------------------------------------------------------
% Cỡ chữ 13pt — đặt qua gói `fontsize` vì lớp report chuẩn chỉ nhận 10/11/12pt
% còn extsizes thì thiếu đúng mức 13pt. Gói này cho phép cỡ chữ tuỳ ý và điều
% chỉnh đồng bộ toàn bộ các cỡ dẫn xuất (\large, \small, chú thích, chỉ số trên...).
\usepackage[fontsize=13pt]{fontsize}

% Cho phép các font Type1 (Computer Modern dùng trong công thức toán) nhận cỡ
% chữ tuỳ ý. Thiếu gói này, mọi cỡ không nằm trong bảng dựng sẵn — trong đó có
% đúng 13pt — đều bị thay thế kèm cảnh báo "Font shape ... not available".
\usepackage{anyfontsize}

\usepackage{geometry}
\geometry{
    a4paper,
    left=3.5cm,      % lề trái rộng để đóng bìa
    right=2.0cm,
    top=3.0cm,
    bottom=3.0cm,
    footskip=1.2cm
}

% ---------------------------------------------------------------------------
% 2. Font và ngôn ngữ tiếng Việt (XeLaTeX)
% ---------------------------------------------------------------------------
\usepackage{fontspec}

% Font được nạp theo TÊN TỆP (không theo tên hệ thống) để chạy được trên mọi
% môi trường: Overleaf, TeX Live, MiKTeX và Tectonic. Tectonic không dùng
% fontconfig nên tra theo tên hệ thống ("TeX Gyre Termes") sẽ thất bại.
% Ba họ font dưới đây là bản tự do tương đương Times, Helvetica và Courier.
\setmainfont{texgyretermes}[
    Extension      = .otf,
    UprightFont    = *-regular,
    BoldFont       = *-bold,
    ItalicFont     = *-italic,
    BoldItalicFont = *-bolditalic,
]
\setsansfont{texgyreheros}[
    Extension      = .otf,
    UprightFont    = *-regular,
    BoldFont       = *-bold,
    ItalicFont     = *-italic,
    BoldItalicFont = *-bolditalic,
]
\setmonofont{texgyrecursor}[
    Extension      = .otf,
    UprightFont    = *-regular,
    BoldFont       = *-bold,
    ItalicFont     = *-italic,
    BoldItalicFont = *-bolditalic,
    Scale          = 0.92,
]

\usepackage{polyglossia}
\setdefaultlanguage{vietnamese}
\setotherlanguage{english}

% Lưới an toàn: ép lại tên tiếng Việt phòng khi gói ngôn ngữ dịch thiếu.
% Dùng \AtBeginDocument thay cho \renewcaptionname để tương thích với MỌI
% phiên bản polyglossia — lệnh \renewcaptionname chỉ có ở các bản mới.
\AtBeginDocument{%
    \renewcommand{\chaptername}{Chương}%
    \renewcommand{\figurename}{Hình}%
    \renewcommand{\tablename}{Bảng}%
    \renewcommand{\contentsname}{Mục lục}%
    \renewcommand{\listfigurename}{Danh mục hình vẽ}%
    \renewcommand{\listtablename}{Danh mục bảng biểu}%
    \renewcommand{\appendixname}{Phụ lục}%
    \renewcommand{\indexname}{Chỉ mục}%
}

% ---------------------------------------------------------------------------
% 3. Giãn dòng và đoạn văn
% ---------------------------------------------------------------------------
\usepackage{setspace}
\onehalfspacing                        % giãn dòng 1,5
\setlength{\parindent}{1cm}
\setlength{\parskip}{0.35em}

% Nới ràng buộc dàn dòng. Văn bản tiếng Việt có nhiều từ đơn âm tiết ngắn và
% hầu như không ngắt từ được, lại xen kẽ nhiều thuật ngữ tiếng Anh cùng tên
% hàm, tên tệp dài — nếu giữ mặc định sẽ còn khá nhiều dòng tràn ra lề.
% \emergencystretch cho TeX một lượng giãn dự phòng ở lượt dàn dòng cuối cùng.
\tolerance=2000
\setlength{\emergencystretch}{3em}
\hbadness=6000                         % bớt cảnh báo dòng thưa chữ không đáng kể

% ---------------------------------------------------------------------------
% 4. Toán học
% ---------------------------------------------------------------------------
\usepackage{amsmath}
\usepackage{amssymb}
\usepackage{amsfonts}

% ---------------------------------------------------------------------------
% 5. Hình ảnh và bảng biểu
% ---------------------------------------------------------------------------
\usepackage{graphicx}
\graphicspath{{figures/}}
\usepackage{float}
\usepackage{booktabs}
\usepackage{multirow}
\usepackage{longtable}
\usepackage{array}
\usepackage{tabularx}
\usepackage{caption}
\usepackage{subcaption}
\captionsetup{font=small, labelfont=bf, skip=6pt}
\captionsetup[table]{position=above}

% Cột canh giữa/phải có độ rộng cố định
\newcolumntype{C}[1]{>{\centering\arraybackslash}p{#1}}
\newcolumntype{R}[1]{>{\raggedleft\arraybackslash}p{#1}}
\newcolumntype{L}[1]{>{\raggedright\arraybackslash}p{#1}}

% ---------------------------------------------------------------------------
% 6. Màu và sơ đồ TikZ (tái sử dụng từ paper_ijai.tex)
% ---------------------------------------------------------------------------
\usepackage{xcolor}
\usepackage{tikz}
\usetikzlibrary{arrows.meta, positioning, fit, calc}

\definecolor{stageblue}{HTML}{DCEBFF}
\definecolor{stagegreen}{HTML}{DDF4E7}
\definecolor{stageorange}{HTML}{FFE8CC}
\definecolor{stagegray}{HTML}{EEF1F5}
\definecolor{stagepurple}{HTML}{EDE4FF}
\definecolor{arrowgray}{HTML}{4B5563}

\tikzset{
    pipelinebox/.style={
        draw=arrowgray, rounded corners=2pt, align=center,
        minimum width=3.1cm, minimum height=1.15cm,
        font=\small, inner sep=5pt
    },
    data/.style={pipelinebox, fill=stagegray},
    classmodel/.style={pipelinebox, fill=stageblue},
    xai/.style={pipelinebox, fill=stageorange},
    sam/.style={pipelinebox, fill=stagepurple},
    segmodel/.style={pipelinebox, fill=stagegreen},
    flow/.style={-{Latex[length=2.6mm]}, thick, draw=arrowgray},
    dashedflow/.style={-{Latex[length=2.6mm]}, thick, dashed, draw=arrowgray},
    lane/.style={draw=arrowgray!45, rounded corners=4pt, inner sep=8pt}
}

% Bảng màu Okabe–Ito, đồng bộ với màu trong toàn bộ notebook
\definecolor{okOrange}{HTML}{E69F00}
\definecolor{okSky}{HTML}{56B4E9}
\definecolor{okGreen}{HTML}{009E73}
\definecolor{okBlue}{HTML}{0072B2}
\definecolor{okVermil}{HTML}{D55E00}

% ---------------------------------------------------------------------------
% 7. Đánh số chương mục và tiêu đề
% ---------------------------------------------------------------------------
\usepackage{titlesec}
\titleformat{\chapter}[display]
    {\normalfont\Large\bfseries\centering}
    {\chaptertitlename\ \thechapter}{12pt}{\LARGE}
\titlespacing*{\chapter}{0pt}{-10pt}{28pt}

\setcounter{secnumdepth}{3}
\setcounter{tocdepth}{2}

% ---------------------------------------------------------------------------
% 8. Header / footer
% ---------------------------------------------------------------------------
\usepackage{fancyhdr}
\pagestyle{fancy}
\fancyhf{}
\fancyfoot[C]{\thepage}
\renewcommand{\headrulewidth}{0pt}
\fancypagestyle{plain}{\fancyhf{}\fancyfoot[C]{\thepage}\renewcommand{\headrulewidth}{0pt}}

% ---------------------------------------------------------------------------
% 9. Liên kết chéo và trích dẫn
% ---------------------------------------------------------------------------
% Cho phép ngắt dòng bên trong văn bản đánh máy (\texttt). Mặc định LaTeX đặt
% \hyphenchar = -1 cho font đánh máy nên các chuỗi dài như tên tệp, đường dẫn
% hay định danh mã nguồn không thể ngắt và tràn ra ngoài lề.
\usepackage[htt]{hyphenat}

\usepackage[hidelinks, unicode]{hyperref}

% xurl cho phép ngắt URL tại bất kỳ ký tự nào. Cần thiết vì một số mục trong
% danh mục tham khảo chứa địa chỉ kho mã nguồn dài. Phải nạp sau hyperref.
\usepackage{xurl}
\hypersetup{
    pdftitle={Phat hien benh la sau rieng bang hoc sau ket hop XAI va pseudo-labeling},
    pdfauthor={},
    bookmarksnumbered=true
}
% Tuỳ chọn `english` ép cleveref nạp bộ định nghĩa tiếng Anh thay vì tự dò theo
% babel/polyglossia. Cần thiết vì cleveref 0.21.x không biết ngôn ngữ
% `vietnamese`; nếu để nó tự dò thì các macro liên từ (\crefpairconjunction...)
% sẽ không được định nghĩa và mọi lệnh \cref nhiều nhãn đều lỗi.
% Toàn bộ chuỗi hiển thị được Việt hoá ngay bên dưới nên tuỳ chọn này không
% làm lộ chữ tiếng Anh nào ra bản in.
\usepackage[nameinlink, english]{cleveref}

% Bản vá tương thích cleveref 0.21.x với polyglossia 1.5x: cleveref chỉ đặt
% \cref@language khi \languagename đã tồn tại, còn polyglossia bản mới định
% nghĩa \languagename muộn hơn. Hệ quả là \cref@language chưa xác định khi
% cleveref chạy hook \AtBeginDocument, gây lỗi "Undefined control sequence".
% Khai báo dự phòng dưới đây vô hại nếu cleveref đã tự đặt được giá trị.
\makeatletter
\providecommand{\cref@language}{english}
\makeatother

% Việt hoá các liên từ khi tham chiếu nhiều nhãn cùng lúc:
%   \cref{a,b}   -> "Hình 1 và Hình 2"
%   \cref{a,b,c} -> "Hình 1, Hình 2 và Hình 3"
%   \crefrange   -> "Hình 1 đến Hình 5"
%
% Phải đặt trong \AtBeginDocument: cleveref chỉ tạo các macro liên từ ở hook
% \AtBeginDocument của chính nó (sau khi đã xác định ngôn ngữ), nên ở thời điểm
% preamble chúng chưa tồn tại và \renewcommand sẽ báo "Command undefined".
% Hook khai báo tại đây chạy sau hook của cleveref vì được đăng ký muộn hơn.
\newcommand{\vnand}{ và~}
\AtBeginDocument{%
    \renewcommand{\crefpairconjunction}{\vnand}%
    \renewcommand{\crefmiddleconjunction}{, }%
    \renewcommand{\creflastconjunction}{\vnand}%
    \renewcommand{\crefpairgroupconjunction}{\vnand}%
    \renewcommand{\crefmiddlegroupconjunction}{, }%
    \renewcommand{\creflastgroupconjunction}{\vnand}%
    \renewcommand{\crefrangeconjunction}{ đến~}%
    \renewcommand{\crefrangepreconjunction}{}%
    \renewcommand{\crefrangepostconjunction}{}%
}

\crefname{figure}{Hình}{Hình}
\crefname{table}{Bảng}{Bảng}
\crefname{chapter}{Chương}{Chương}
\crefname{section}{Mục}{Mục}
\crefname{equation}{Công thức}{Công thức}

% Dùng backend=bibtex thay vì biber: bibtex có sẵn trong Tectonic và mọi bản
% TeX Live/MiKTeX, trong khi biber là chương trình ngoài phải cài riêng.
% Tệp refs.bib chỉ chứa ký tự ASCII và các dấu phụ đã escape theo cú pháp
% LaTeX (\'e, \"a...) nên bibtex xử lý được đầy đủ.
% Khai báo language=english để biblatex dùng bộ chuỗi tiếng Anh cho phần tài
% liệu tham khảo — phù hợp vì toàn bộ tài liệu trích dẫn đều là công bố quốc tế.
% autolang=none: chặn biblatex tự đổi ngôn ngữ theo ngôn ngữ tài liệu. Không có
% tuỳ chọn này, biblatex sẽ cố dựng mục tài liệu tham khảo bằng tiếng Việt,
% nhưng không tồn tại tệp vietnamese.lbx nên hàng loạt chuỗi ("and", "pages",
% "in"...) bị bỏ chưa dịch và hiện ra thô trong bản in.
\usepackage[backend=bibtex, style=ieee, sorting=none,
            language=english, autolang=none]{biblatex}

% biblatex chốt ngôn ngữ mục tài liệu tham khảo ngay lúc nạp gói, dựa theo ngôn
% ngữ chính của tài liệu. Vì không tồn tại tệp vietnamese.lbx, cần ánh xạ tường
% minh sang english.lbx — đây là cơ chế chính thức của biblatex cho các ngôn ngữ
% chưa có bộ chuỗi riêng. Thiếu dòng này, các chuỗi "and", "pages", "in",
% "jourvol"... sẽ bị bỏ chưa dịch và hiện thô trong danh mục tham khảo.
\DeclareLanguageMapping{vietnamese}{english}

\addbibresource{refs.bib}

% ---------------------------------------------------------------------------
% 10. Ghi chú việc còn thiếu — mọi chỗ số liệu chưa xác minh phải dùng \todo
% ---------------------------------------------------------------------------
% \TODO được cố tình định nghĩa ở dạng NỘI DÒNG (không dùng todonotes) để có thể
% đặt ở bất kỳ đâu, kể cả giữa các dòng của trang bìa — nơi lệnh \todo[inline]
% của todonotes sẽ phá vỡ cấu trúc đoạn và gây lỗi "There's no line here to end".
% Nền vàng viền đỏ bảo đảm không thể bỏ sót khi rà soát bản in.
% KHÔNG bọc trong \colorbox: hộp màu là một hộp ngang không thể ngắt dòng, nên
% một ghi chú dài sẽ tràn ra ngoài lề hàng chục xăng-ti-mét. Chữ đỏ đậm ngắt
% dòng bình thường mà vẫn đủ nổi bật để không bỏ sót khi rà bản in.
\newcommand{\TODO}[1]{%
    \textcolor{red}{\textbf{[\,CẦN BỔ SUNG: #1\,]}}%
}

% ---------------------------------------------------------------------------
% 11. Macro dùng chung
% ---------------------------------------------------------------------------
\usepackage{siunitx}
\sisetup{output-decimal-marker={,}, group-separator={.}, group-minimum-digits=4}

% Tên kỹ thuật viết nhất quán trong toàn luận văn
\newcommand{\EffNet}{EfficientNet-B0}
\newcommand{\GradCAM}{Grad-CAM}
\newcommand{\GradCAMpp}{Grad-CAM++}
\newcommand{\UNetpp}{UNet++}
\newcommand{\SAM}{SAM}

% Thuật ngữ tiếng Anh giữ nguyên, in nghiêng ở lần xuất hiện đầu
\newcommand{\en}[1]{\textit{#1}}

\usepackage{enumitem}
\setlist{itemsep=2pt, topsep=4pt, parsep=0pt}

\usepackage{microtype}


\begin{document}

% ===========================================================================
%  PHẦN ĐẦU — đánh số La Mã
% ===========================================================================
\pagenumbering{roman}

% ---------------------------------------------------------------------------
%  Trang bìa
%
%  Lưu ý: KHÔNG dùng môi trường `titlepage` của LaTeX. Môi trường đó tự đặt lại
%  bộ đếm trang về 1 khi kết thúc; dùng hai lần cho bìa chính và bìa phụ sẽ tạo
%  ra hai trang cùng mang số i và ii, khiến hyperref sinh trùng mỏ neo trang
%  (cảnh báo "Object @page.i already defined") và kéo theo cảnh báo tham chiếu
%  chưa xác định. Ở đây dùng khối căn giữa thông thường rồi \clearpage.
% ---------------------------------------------------------------------------
\begingroup
\thispagestyle{empty}
\centering
\vspace*{0.5cm}

{\large\bfseries \TODO{Tên Bộ / Trường Đại học chủ quản}}\\[0.3cm]
{\large\bfseries \TODO{Tên Khoa / Viện}}\\[2.5cm]

{\Large\bfseries \TODO{Họ và tên học viên}}\\[2.5cm]

{\LARGE\bfseries
PHÁT HIỆN VÀ PHÂN ĐOẠN BỆNH LÁ SẦU RIÊNG\\[0.35cm]
BẰNG HỌC SÂU KẾT HỢP AI KHẢ DIỄN GIẢI\\[0.35cm]
VÀ GÁN NHÃN GIẢ
}\\[2.5cm]

{\large LUẬN VĂN THẠC SĨ}\\[0.4cm]
{\large \TODO{Ngành và mã ngành đào tạo}}\\[3cm]

\vfill
{\large \TODO{Địa danh} --- \the\year}
\par
\endgroup
\clearpage

% ---------------------------------------------------------------------------
%  Trang bìa phụ (có tên người hướng dẫn)
% ---------------------------------------------------------------------------
\begingroup
\thispagestyle{empty}
\centering
\vspace*{0.5cm}

{\large\bfseries \TODO{Tên Bộ / Trường Đại học chủ quản}}\\[0.3cm]
{\large\bfseries \TODO{Tên Khoa / Viện}}\\[2cm]

{\Large\bfseries \TODO{Họ và tên học viên}}\\[2cm]

{\LARGE\bfseries
PHÁT HIỆN VÀ PHÂN ĐOẠN BỆNH LÁ SẦU RIÊNG\\[0.35cm]
BẰNG HỌC SÂU KẾT HỢP AI KHẢ DIỄN GIẢI\\[0.35cm]
VÀ GÁN NHÃN GIẢ
}\\[1.5cm]

{\large LUẬN VĂN THẠC SĨ}\\[0.3cm]
{\large \TODO{Ngành và mã ngành đào tạo}}\\[2.5cm]

{\large NGƯỜI HƯỚNG DẪN KHOA HỌC}\\[0.4cm]
{\large\bfseries \TODO{Học hàm, học vị và họ tên người hướng dẫn}}\\[3cm]

\vfill
{\large \TODO{Địa danh} --- \the\year}
\par
\endgroup
\clearpage

% ---------------------------------------------------------------------------
%  Lời cam đoan
% ---------------------------------------------------------------------------
\chapter*{LỜI CAM ĐOAN}
\addcontentsline{toc}{chapter}{Lời cam đoan}

Tôi xin cam đoan luận văn này là công trình nghiên cứu của riêng tôi, được thực
hiện dưới sự hướng dẫn khoa học của \TODO{tên người hướng dẫn}. Các số liệu, kết
quả thực nghiệm trình bày trong luận văn đều được sinh ra trực tiếp từ mã nguồn và
các \en{notebook} thực nghiệm kèm theo, có thể tái tạo được và chưa từng được
công bố trong bất kỳ công trình nào khác. Mọi tham khảo từ tài liệu của tác giả
khác đều được trích dẫn đầy đủ.

Tôi xin chịu hoàn toàn trách nhiệm về nội dung luận văn của mình.

\vspace{1.5cm}
\begin{flushright}
\begin{minipage}{7cm}
\centering
\TODO{Địa danh}, ngày \quad tháng \quad năm \the\year\\[0.3cm]
\textbf{Học viên thực hiện}\\[2.5cm]
\TODO{Họ và tên học viên}
\end{minipage}
\end{flushright}

% ---------------------------------------------------------------------------
%  Lời cảm ơn
% ---------------------------------------------------------------------------
\chapter*{LỜI CẢM ƠN}
\addcontentsline{toc}{chapter}{Lời cảm ơn}

\TODO{Viết lời cảm ơn: người hướng dẫn, cơ sở đào tạo, đơn vị cung cấp dữ liệu, gia đình.}

% ---------------------------------------------------------------------------
%  Mục lục và các danh mục
% ---------------------------------------------------------------------------
\tableofcontents
\listoffigures
\addcontentsline{toc}{chapter}{Danh mục hình vẽ}
\listoftables
\addcontentsline{toc}{chapter}{Danh mục bảng biểu}

% !TeX root = ../main.tex
\chapter*{DANH MỤC TỪ VIẾT TẮT VÀ THUẬT NGỮ}
\addcontentsline{toc}{chapter}{Danh mục từ viết tắt và thuật ngữ}

\begin{longtable}{L{2.6cm} L{4.6cm} L{6.5cm}}
\toprule
\textbf{Viết tắt} & \textbf{Tiếng Anh đầy đủ} & \textbf{Giải thích} \\
\midrule
\endfirsthead
\toprule
\textbf{Viết tắt} & \textbf{Tiếng Anh đầy đủ} & \textbf{Giải thích} \\
\midrule
\endhead
\bottomrule
\endfoot

AI      & Artificial Intelligence            & Trí tuệ nhân tạo \\
CAM     & Class Activation Map               & Bản đồ kích hoạt theo lớp \\
CC      & Connected Component                & Thành phần liên thông \\
CNN     & Convolutional Neural Network       & Mạng nơ-ron tích chập \\
CRF     & Conditional Random Field           & Trường ngẫu nhiên có điều kiện \\
EDA     & Exploratory Data Analysis          & Phân tích khám phá dữ liệu \\
FN      & False Negative                     & Âm tính giả --- bỏ sót vùng bệnh \\
FP      & False Positive                     & Dương tính giả --- báo nhầm vùng lành \\
GPU     & Graphics Processing Unit           & Bộ xử lý đồ hoạ \\
Grad-CAM & Gradient-weighted Class Activation Mapping & Bản đồ kích hoạt có trọng số gradient \\
HSV     & Hue--Saturation--Value             & Không gian màu sắc độ--bão hoà--độ sáng \\
IoU     & Intersection over Union            & Tỉ số giao trên hợp \\
KDE     & Kernel Density Estimation          & Ước lượng mật độ hạt nhân \\
LR      & Learning Rate                      & Tốc độ học \\
ReLU    & Rectified Linear Unit              & Hàm kích hoạt tuyến tính chỉnh lưu \\
RGB     & Red--Green--Blue                   & Không gian màu đỏ--lục--lam \\
SAM     & Segment Anything Model             & Mô hình nền tảng phân đoạn tổng quát \\
TP      & True Positive                      & Dương tính thật \\
ViT     & Vision Transformer                 & Mạng Transformer cho thị giác máy tính \\
XAI     & Explainable Artificial Intelligence & Trí tuệ nhân tạo khả diễn giải \\
\end{longtable}

\vspace{0.6cm}
\noindent\textbf{Quy ước thuật ngữ.} Các thuật ngữ tiếng Anh đã trở thành chuẩn
trong lĩnh vực (\en{pseudo-label}, \en{fine-tuning}, \en{early stopping},
\en{ground truth}, \en{backbone}, \en{encoder}, \en{decoder}\dots) được giữ
nguyên dạng tiếng Anh và in nghiêng ở lần xuất hiện đầu tiên, kèm giải thích
tiếng Việt trong ngoặc. Ở các lần sau, thuật ngữ được dùng trực tiếp không in
nghiêng nhằm giữ mạch đọc.

\vspace{0.4cm}
\noindent\textbf{Quy ước ký hiệu phiên bản.} Luận văn xây dựng ba phiên bản
pipeline phân đoạn, ký hiệu thống nhất là \textbf{V1}, \textbf{V2}, \textbf{V3}.
Mỗi phiên bản gồm một bộ \en{pseudo-label} và một mô hình phân đoạn huấn luyện
trên bộ nhãn đó; chi tiết trong \cref{chap:phuongphap}.

\vspace{0.4cm}
\noindent\textbf{Quy ước số thập phân.} Luận văn dùng dấu phẩy làm dấu thập phân
theo chuẩn tiếng Việt (ví dụ: 0,6239). Trong các đoạn mã nguồn và tên tệp trích
dẫn nguyên trạng, dấu chấm được giữ nguyên theo đúng dữ liệu gốc.


% ===========================================================================
%  PHẦN THÂN — đánh số Ả Rập
% ===========================================================================
\pagenumbering{arabic}
\setcounter{page}{1}

% !TeX root = ../main.tex
\chapter{MỞ ĐẦU}
\label{chap:modau}

\section{Bối cảnh nghiên cứu}
\label{sec:boicanh}

Sầu riêng (\textit{Durio zibethinus}) là một trong những cây ăn quả nhiệt đới có
giá trị kinh tế cao nhất khu vực Đông Nam Á. Tại Việt Nam, cây sầu riêng được
canh tác tập trung ở Tây Nguyên, Đông Nam Bộ và đồng bằng sông Cửu Long, đóng vai
trò quan trọng trong cơ cấu thu nhập của nhiều hộ nông dân. Đặc điểm của loại cây
này là chu kỳ kiến thiết cơ bản dài --- thường phải mất từ bốn đến sáu năm kể từ
khi trồng mới cho thu hoạch ổn định --- nên mỗi tổn thất do sâu bệnh gây ra đều
kéo theo hệ quả kinh tế lâu dài, không thể khắc phục trong một vụ.

Trong các nhóm bệnh hại trên cây sầu riêng, bệnh trên lá chiếm vị trí đặc biệt vì
hai lý do. Thứ nhất, lá là cơ quan quang hợp; diện tích lá bị tổn thương giảm
trực tiếp làm suy giảm khả năng tích luỹ chất khô, ảnh hưởng đến năng suất và
chất lượng quả. Thứ hai, triệu chứng trên lá thường xuất hiện \emph{sớm hơn}
triệu chứng trên thân và quả, nên lá đóng vai trò chỉ báo cho phép can thiệp kịp
thời. Việc phát hiện chính xác loại bệnh và xác định được phạm vi tổn thương ngay
từ giai đoạn đầu vì thế có ý nghĩa quyết định đối với hiệu quả phòng trừ.

\subsection{Nguồn gốc bộ dữ liệu}
\label{subsec:nguongoc}

Bộ dữ liệu được sử dụng trong luận văn gồm \num{4437} ảnh lá sầu riêng thuộc năm
lớp: bốn lớp bệnh và một lớp lá khoẻ. Chi tiết về phân bố lớp, đặc trưng màu sắc
và cách phân chia tập dữ liệu được trình bày trong \cref{sec:dulieu}.

Ảnh được thu thập trực tiếp tại các vườn trồng sầu riêng ở khu vực Đông Nam Á.
Nhãn ở mức ảnh --- tức việc xác định mỗi ảnh thuộc lớp bệnh nào --- do cán bộ
chuyên môn nông nghiệp thực hiện và thẩm định; nhóm nghiên cứu không tự gán nhãn
dựa trên quan sát của người không có chuyên môn bệnh học thực vật. Đây là cơ sở
để xem nhãn mức ảnh của bộ dữ liệu là đáng tin cậy, và cũng là tiền đề cho toàn
bộ quy trình học giám sát yếu được đề xuất: nếu nhãn mức ảnh sai thì bản đồ chú ý
sinh ra từ nó, và do đó cả nhãn giả lẫn mô hình phân đoạn, đều mất giá trị.

Bộ dữ liệu và trọng số mô hình được chia sẻ theo yêu cầu hợp lý gửi tới tác giả
liên hệ; luận văn không công bố công khai dữ liệu thô.

Cần lưu ý ba đặc điểm về trạng thái của dữ liệu khi được đưa vào nghiên cứu này,
vì chúng ảnh hưởng trực tiếp đến phạm vi kết luận có thể rút ra:

\begin{itemize}
    \item \textbf{Ảnh đã được tiền xử lý trước.} Toàn bộ \num{4437} ảnh đều ở
          đúng kích thước $224 \times 224$ điểm ảnh với tỉ lệ khung hình bằng 1
          và độ lệch chuẩn bằng 0 --- nghĩa là việc cắt và co giãn đã được thực
          hiện ở khâu trước. Hệ quả là luận văn không thể khảo sát ảnh hưởng của
          độ phân giải đầu vào lên chất lượng phân đoạn.
    \item \textbf{Siêu dữ liệu EXIF đã bị loại bỏ.} Kiểm tra trực tiếp các tệp
          ảnh cho thấy không còn thông tin về thời điểm chụp, thiết bị chụp hay
          toạ độ định vị. Vì vậy không thể phân tích ảnh hưởng của điều kiện chụp
          --- ánh sáng, thiết bị, mùa vụ --- lên hiệu năng mô hình.
    \item \textbf{Cấu trúc tên tệp phản ánh quy trình gán nhãn.} Các tệp được đặt
          tên theo dạng \texttt{to\_\allowbreak label\_\allowbreak N.\allowbreak jpg} với chỉ số $N$ đánh liên tục
          xuyên suốt các lớp, cho thấy dữ liệu ban đầu tồn tại dưới dạng một kho
          ảnh chung chờ gán nhãn, sau đó mới được phân vào các thư mục theo lớp.
\end{itemize}

\TODO{Bổ sung nếu lấy được từ nhóm thu thập dữ liệu gốc: (1) địa phương cụ thể
và thời gian thu thập; (2) thiết bị chụp và độ phân giải ảnh gốc trước khi co về
$224\times224$; (3) số lượng chuyên gia tham gia gán nhãn và mức độ đồng thuận
giữa họ nếu có nhiều hơn một người; (4) giấy phép sử dụng dữ liệu. Bốn thông tin
này không tồn tại trong kho mã nguồn và cũng không được ghi trong công bố liên
quan của nhóm.}

Phương pháp chẩn đoán truyền thống dựa hoàn toàn vào quan sát trực tiếp của cán
bộ kỹ thuật nông nghiệp. Cách làm này bộc lộ ba hạn chế cố hữu: tốn nhiều công
lao động khi diện tích canh tác lớn; kết quả mang tính chủ quan và thiếu nhất
quán giữa những người quan sát khác nhau; và đặc biệt khó phân biệt các bệnh có
biểu hiện hình thái tương tự nhau. Những hạn chế này tạo ra nhu cầu thực tế về
một công cụ chẩn đoán tự động, khách quan và có khả năng triển khai ở quy mô lớn.

Học sâu (\en{deep learning}) đã chứng minh hiệu quả trong bài toán nhận dạng bệnh
cây trồng qua nhiều nghiên cứu, từ công trình nền tảng của Mohanty và cộng
sự~\cite{mohanty2016plantdisease} trên bộ dữ liệu PlantVillage, đến các khảo sát
tổng quan của Kamilaris và Prenafeta-Boldú~\cite{kamilaris2018survey} hay
Ferentinos~\cite{ferentinos2018plantdisease}. Tuy nhiên, phần lớn các nghiên cứu
này dừng lại ở mức \emph{phân loại ảnh}: mô hình trả lời câu hỏi ``ảnh này thuộc
bệnh nào'' nhưng không trả lời được câu hỏi ``vùng bệnh nằm ở đâu và rộng bao
nhiêu''. Đây chính là khoảng trống mà luận văn hướng tới.

\section{Phát biểu vấn đề}
\label{sec:vande}

Luận văn giải quyết ba vấn đề gắn bó chặt chẽ với nhau.

\subsection{Chi phí gán nhãn ở mức điểm ảnh}

Để huấn luyện một mô hình phân đoạn (\en{segmentation}) theo cách giám sát đầy đủ,
mỗi ảnh cần một mặt nạ nhị phân (\en{binary mask}) chỉ rõ từng điểm ảnh thuộc vùng
bệnh hay vùng lành. Công việc khoanh vùng thủ công này đòi hỏi chuyên gia bệnh
học thực vật, tốn khoảng vài phút cho mỗi ảnh, và với bộ dữ liệu \num{4437} ảnh
thì tổng chi phí trở nên khó chấp nhận trong điều kiện nghiên cứu thông thường.
Trong khi đó, nhãn ở \emph{mức ảnh} --- tức chỉ ghi nhãn ``ảnh này bị bệnh
Phomopsis'' --- lại rẻ hơn nhiều bậc và thường đã có sẵn.

Bài toán đặt ra là: \emph{liệu có thể tận dụng nhãn mức ảnh vốn rẻ để tạo ra
nhãn mức điểm ảnh mà không cần khoanh vùng thủ công hay không?} Đây là phạm trù
của học giám sát yếu (\en{weakly-supervised learning}).

\subsection{Tính khả diễn giải của mô hình}

Mạng nơ-ron tích chập vận hành như một hộp đen: người dùng nhận được nhãn dự đoán
kèm độ tin cậy nhưng không biết mô hình căn cứ vào đặc trưng nào. Trong bối cảnh
nông nghiệp, nơi một khuyến cáo phun thuốc sai có thể gây thiệt hại trực tiếp,
sự thiếu minh bạch này là rào cản thực sự đối với việc chấp nhận công nghệ. Cán
bộ kỹ thuật cần nhìn thấy mô hình ``đang nhìn vào đâu'' để đối chiếu với kinh
nghiệm chuyên môn trước khi tin vào kết luận của nó.

\subsection{Nhu cầu định vị chính xác vùng tổn thương}

Nhãn phân loại cho biết \emph{có} bệnh, nhưng quyết định canh tác lại cần biết
\emph{mức độ} bệnh. Tỉ lệ diện tích lá bị tổn thương là căn cứ để lựa chọn giữa
theo dõi tiếp, phun thuốc cục bộ hay cắt tỉa. Chỉ mô hình phân đoạn mới cung cấp
được đại lượng này.

Ba vấn đề trên gợi ý một hướng tiếp cận thống nhất: nếu bản đồ chú ý của mô hình
phân loại có thể được trực quan hoá để giải quyết vấn đề thứ hai, thì chính bản
đồ đó cũng có thể được chuyển hoá thành mặt nạ phân đoạn để giải quyết vấn đề thứ
nhất và thứ ba. Đây là ý tưởng trung tâm của luận văn.

\section{Mục tiêu và câu hỏi nghiên cứu}
\label{sec:muctieu}

\subsection{Mục tiêu tổng quát}

Xây dựng và đánh giá một quy trình học giám sát yếu hoàn chỉnh cho bài toán phát
hiện bệnh lá sầu riêng, đi từ phân loại ảnh, qua giải thích bằng bản đồ kích
hoạt, đến sinh nhãn giả và huấn luyện mô hình phân đoạn --- toàn bộ mà không sử
dụng bất kỳ nhãn thủ công nào ở mức điểm ảnh.

\subsection{Mục tiêu cụ thể}

\begin{enumerate}
    \item Huấn luyện mô hình phân loại \EffNet{} cho năm lớp bệnh lá sầu riêng và
          đánh giá đầy đủ trên tập kiểm thử độc lập.
    \item Áp dụng \GradCAM{} và \GradCAMpp{} đa tỉ lệ để trực quan hoá vùng chú ý
          của mô hình phân loại, đồng thời kiểm chứng rằng mô hình học đặc trưng
          bệnh học có ý nghĩa chứ không phải học nhiễu nền.
    \item Thiết kế quy trình chuyển hoá bản đồ nhiệt thành mặt nạ nhị phân, với
          cơ chế kiểm soát độ phủ theo từng lớp bệnh.
    \item Khảo sát mô hình nền tảng \SAM{} trong vai trò bộ tinh chỉnh biên
          (\en{label refiner}) cho nhãn giả.
    \item Huấn luyện và so sánh ba phiên bản mô hình phân đoạn tương ứng với ba
          thế hệ nhãn giả, phân tích tách bạch ảnh hưởng của từng yếu tố.
    \item Chỉ rõ giới hạn phương pháp luận của việc đánh giá phân đoạn khi không
          có nhãn chuẩn (\en{ground truth}) ở mức điểm ảnh.
\end{enumerate}

\subsection{Câu hỏi nghiên cứu}

\begin{itemize}
    \item[\textbf{CH1.}] Mô hình phân loại \EffNet{} đạt độ chính xác bao nhiêu
          trên năm lớp bệnh lá sầu riêng, và những cặp lớp nào dễ bị nhầm lẫn
          nhất?
    \item[\textbf{CH2.}] Bản đồ chú ý sinh ra từ mô hình phân loại có tập trung
          vào vùng tổn thương thực sự hay không, và mức độ tin cậy của chúng khác
          nhau thế nào giữa các lớp bệnh?
    \item[\textbf{CH3.}] Việc tinh chỉnh biên bằng \SAM{} có cải thiện chất lượng
          nhãn giả, và qua đó cải thiện mô hình phân đoạn, hay không?
    \item[\textbf{CH4.}] Trong điều kiện không có nhãn chuẩn ở mức điểm ảnh, các
          chỉ số IoU và Dice giữa những phiên bản khác nhau của quy trình có so
          sánh trực tiếp được với nhau không? Nếu không thì phép so sánh nào là
          hợp lệ?
\end{itemize}

Câu hỏi CH4 là câu hỏi mang tính phương pháp luận và, như sẽ thấy ở
\cref{chap:thucnghiem}, câu trả lời cho nó làm thay đổi hoàn toàn cách diễn giải
các con số của ba câu hỏi còn lại.

\section{Đối tượng, phạm vi và giới hạn}
\label{sec:phamvi}

\subsection{Đối tượng nghiên cứu}

Đối tượng là ảnh lá sầu riêng đơn lẻ ở độ phân giải $224 \times 224$ điểm ảnh,
thuộc năm lớp: Algal Leaf Spot (đốm tảo), Allocaridara Attack (rầy phấn tấn
công), Healthy Leaf (lá khoẻ), Leaf Blight (cháy lá) và Phomopsis Leaf Spot (đốm
lá Phomopsis).

\subsection{Phạm vi}

\begin{itemize}
    \item \textbf{Bài toán phân loại}: đa lớp, năm lớp, một nhãn trên mỗi ảnh.
    \item \textbf{Bài toán phân đoạn}: nhị phân --- phân biệt vùng bệnh với phần
          còn lại. Luận văn \emph{không} thực hiện phân đoạn đa lớp, tức không
          phân biệt loại bệnh ở mức điểm ảnh.
    \item \textbf{Dữ liệu}: chỉ sử dụng bộ dữ liệu \num{4437} ảnh nêu trên; không
          thu thập bổ sung và không sử dụng dữ liệu ngoài.
    \item \textbf{Mô hình}: các kiến trúc có kích thước vừa phải, huấn luyện được
          trên một GPU đơn.
\end{itemize}

\subsection{Giới hạn cần lưu ý trước khi đọc kết quả}
\label{subsec:gioihan-quantrong}

Có một giới hạn mang tính bản chất cần được nêu ngay từ đầu, vì nó chi phối cách
diễn giải toàn bộ kết quả phân đoạn trong luận văn:

\begin{quote}
\textbf{Bộ dữ liệu không có nhãn chuẩn ở mức điểm ảnh.} Vì vậy, mọi chỉ số IoU và
Dice được báo cáo trong \cref{chap:thucnghiem} đo mức độ mô hình phân đoạn
\emph{tái tạo lại được bộ nhãn giả} mà chính nó được huấn luyện trên đó --- chứ
\emph{không} đo độ chính xác so với vùng bệnh thật.
\end{quote}

Hệ quả trực tiếp của giới hạn này là: các phiên bản sử dụng bộ nhãn giả khác nhau
sẽ được đánh giá trên những mốc tham chiếu khác nhau, nên giá trị IoU của chúng
\emph{không so sánh trực tiếp được với nhau}. Luận văn không né tránh mà phân
tích thẳng vào vấn đề này ở \cref{sec:caveat}, đồng thời chỉ rõ phép so sánh nào
vẫn giữ được tính hợp lệ nội tại.

Ngoài ra, luận văn có ba giới hạn thứ cấp: (i) các hạn chế bắt nguồn từ trạng
thái dữ liệu đầu vào đã nêu ở \cref{subsec:nguongoc} --- ảnh đã được co về
$224 \times 224$ và siêu dữ liệu EXIF đã bị loại bỏ --- khiến không thể khảo sát
ảnh hưởng của độ phân giải cũng như điều kiện chụp; (ii) các ngưỡng độ phủ theo
lớp được chọn theo suy luận bệnh học chứ chưa được hiệu chỉnh dựa trên nhãn
chuẩn; (iii) hệ thống được đánh giá trên ảnh chụp lá đơn lẻ, chưa kiểm chứng
trên ảnh chụp toàn tán trong điều kiện đồng ruộng.

\section{Đóng góp của luận văn}
\label{sec:donggop}

\begin{enumerate}
    \item \textbf{Một quy trình học giám sát yếu bốn giai đoạn hoàn chỉnh} cho
          bệnh lá sầu riêng: phân loại $\rightarrow$ giải thích bằng \GradCAMpp{}
          $\rightarrow$ tinh chỉnh biên bằng \SAM{} $\rightarrow$ phân đoạn bằng
          \UNetpp{}, được cài đặt đầy đủ và có thể tái tạo.

    \item \textbf{Cơ chế ngưỡng phân vị theo lớp} (\en{per-class percentile
          thresholding}) cho phép kiểm soát trực tiếp độ phủ của mặt nạ theo đặc
          điểm từng bệnh, thay cho ngưỡng Otsu vốn giả định phân bố hai đỉnh và
          gây phân đoạn thừa khi bản đồ nhiệt được hợp nhất đa tỉ lệ. Cơ chế này
          đạt độ phủ mục tiêu với sai lệch dưới 0,4\% cho mọi lớp.

    \item \textbf{Phân tích định lượng về \SAM{} trong vai trò bộ tinh chỉnh
          nhãn}, bao gồm việc chỉ ra và giải thích một chế độ thất bại chưa được
          nhấn mạnh trong tài liệu: cơ chế bảo vệ dựa trên IoU chỉ chặn được
          trường hợp \SAM{} \emph{trôi} sang vùng khác, mà không chặn được trường
          hợp \SAM{} \emph{nở} mặt nạ ra bao trọn cả chiếc lá. Hiện tượng nở mặt
          nạ được đo cụ thể: độ phủ trung bình tăng từ 15,04\% lên 25,43\%.

    \item \textbf{Một phân tích tách bạch (\en{ablation}) trung thực} về khả năng
          so sánh của các chỉ số phân đoạn khi thiếu nhãn chuẩn. Luận văn chỉ ra
          rằng phiên bản có chỉ số cao nhất về mặt số học lại không phải phiên bản
          định vị bệnh tốt nhất, và giải thích tại sao --- một kết luận có giá trị
          cảnh báo cho các nghiên cứu học giám sát yếu tương tự.

    \item \textbf{Bộ tài liệu kiểm kê số liệu} (\texttt{INVENTORY.md}) truy vết
          mọi con số trong luận văn về đúng notebook và ô lệnh đã sinh ra nó, kèm
          danh sách các điểm chưa đủ bằng chứng. Cách làm này bảo đảm tính minh
          bạch và tái lập của kết quả.
\end{enumerate}

\section{Bố cục luận văn}
\label{sec:bocuc}

Phần còn lại của luận văn được tổ chức như sau.

\textbf{\Cref{chap:cosolythuyet}} trình bày cơ sở lý thuyết: mạng nơ-ron tích
chập và học chuyển giao, kiến trúc U-Net cùng biến thể \UNetpp{}, các hàm mất mát
cho dữ liệu mất cân bằng, họ phương pháp \GradCAM{}, khung học giám sát yếu, mô
hình nền tảng \SAM{}, và các kỹ thuật xử lý ảnh cổ điển được sử dụng.

\textbf{\Cref{chap:phuongphap}} mô tả phương pháp đề xuất: tổng quan quy trình
bốn giai đoạn, phân tích khám phá bộ dữ liệu, thiết kế mô hình phân loại, và ba
thế hệ nhãn giả V1--V3 cùng lý do kỹ thuật của từng cải tiến.

\textbf{\Cref{chap:thucnghiem}} trình bày kết quả thực nghiệm: hiệu năng phân
loại, phân tích định tính bản đồ chú ý, chất lượng nhãn giả qua các phiên bản,
hiện tượng nở mặt nạ của \SAM{}, kết quả phân đoạn, phân tích tách bạch và phân
tích lỗi.

\textbf{\Cref{chap:ketluan}} tổng kết đóng góp, trả lời trực tiếp bốn câu hỏi
nghiên cứu, thảo luận hạn chế và đề xuất hướng phát triển.

Bốn phụ lục cung cấp bảng siêu tham số đầy đủ, nhật ký huấn luyện chi tiết, mô tả
ứng dụng minh hoạ và hướng dẫn tái tạo toàn bộ kết quả.

% !TeX root = ../main.tex
\chapter{CƠ SỞ LÝ THUYẾT VÀ NGHIÊN CỨU LIÊN QUAN}
\label{chap:cosolythuyet}

Chương này trình bày nền tảng lý thuyết của các thành phần được sử dụng trong quy
trình đề xuất. Thứ tự trình bày bám theo trình tự của chính quy trình: từ mạng
phân loại (\cref{sec:cnn}), qua kiến trúc phân đoạn (\cref{sec:unet}) và hàm mất
mát (\cref{sec:loss}), đến các phương pháp giải thích (\cref{sec:xai}), khung học
giám sát yếu (\cref{sec:weaksup}), mô hình nền tảng \SAM{} (\cref{sec:sam}), các
kỹ thuật xử lý ảnh cổ điển (\cref{sec:xulyanh}) và cuối cùng là hệ chỉ số đánh
giá (\cref{sec:metrics}).

\section{Mạng nơ-ron tích chập và học chuyển giao}
\label{sec:cnn}

\subsection{Mạng nơ-ron tích chập}

Mạng nơ-ron tích chập (\en{Convolutional Neural Network}, CNN) là kiến trúc chủ
đạo trong thị giác máy tính hiện đại. Khác với mạng kết nối đầy đủ, CNN khai thác
hai giả định cấu trúc phù hợp với dữ liệu ảnh: \emph{tính cục bộ} --- một điểm
ảnh liên quan mật thiết nhất với các điểm lân cận --- và \emph{tính bất biến tịnh
tiến} --- một đặc trưng như cạnh hay góc mang cùng ý nghĩa dù xuất hiện ở vị trí
nào trong ảnh.

Phép tích chập hai chiều giữa ảnh đầu vào $X$ và nhân tích chập $K$ kích thước
$m \times n$ được định nghĩa:

\begin{equation}
    (X * K)(i, j) = \sum_{a=0}^{m-1} \sum_{b=0}^{n-1} X(i + a,\, j + b) \cdot K(a, b)
    \label{eq:convolution}
\end{equation}

Trọng số của nhân $K$ được chia sẻ trên toàn bộ ảnh, làm giảm mạnh số tham số so
với mạng kết nối đầy đủ và đồng thời áp đặt tính bất biến tịnh tiến vào chính kiến
trúc. Một khối tích chập điển hình gồm phép tích chập, chuẩn hoá theo lô
(\en{batch normalization}) và hàm kích hoạt phi tuyến. Chồng nhiều khối như vậy
tạo ra một phân cấp đặc trưng: các tầng nông học cạnh và kết cấu, các tầng sâu học
các bộ phận và khái niệm ngữ nghĩa. Trong bài toán của luận văn, phân cấp này
tương ứng với việc mạng học từ kết cấu bề mặt lá ở tầng nông đến hình thái đốm
bệnh ở tầng sâu.

\subsection{Học chuyển giao}

Bộ dữ liệu của luận văn có \num{3104} ảnh huấn luyện --- một con số nhỏ so với
quy mô cần thiết để huấn luyện CNN từ đầu. Huấn luyện từ trạng thái khởi tạo ngẫu
nhiên trên tập nhỏ như vậy hầu như chắc chắn dẫn tới quá khớp (\en{overfitting}):
mạng ghi nhớ tập huấn luyện thay vì học quy luật khái quát.

Học chuyển giao (\en{transfer learning}) giải quyết vấn đề này bằng cách khởi tạo
mạng từ trọng số đã được huấn luyện trên một tập dữ liệu lớn, phổ biến nhất là
ImageNet với hơn 1,2 triệu ảnh thuộc 1000 lớp~\cite{deng2009imagenet}. Cơ sở của
cách làm là các đặc trưng ở tầng nông --- cạnh, góc, kết cấu, chuyển sắc màu ---
mang tính tổng quát và hữu ích cho hầu hết bài toán thị giác, bất kể miền dữ liệu.

Có hai chiến lược chính:

\begin{itemize}
    \item \textbf{Trích xuất đặc trưng} (\en{feature extraction}): đóng băng toàn
          bộ phần trích đặc trưng, chỉ huấn luyện tầng phân loại cuối. Phù hợp khi
          dữ liệu đích rất nhỏ hoặc rất giống miền nguồn.
    \item \textbf{Tinh chỉnh toàn bộ} (\en{full fine-tuning}): cập nhật mọi trọng
          số với tốc độ học nhỏ. Phù hợp khi dữ liệu đích đủ lớn và khác biệt đáng
          kể so với miền nguồn.
\end{itemize}

Ảnh lá bệnh khác xa phân bố ảnh ImageNet về cả nội dung lẫn kết cấu, trong khi
\num{3104} ảnh là đủ để tránh quên thảm hoạ (\en{catastrophic forgetting}) nếu
dùng tốc độ học đủ nhỏ. Vì vậy luận văn chọn tinh chỉnh toàn bộ với tốc độ học
$10^{-4}$; lý do lựa chọn được biện luận chi tiết ở \cref{chap:phuongphap}.

\subsection{Kiến trúc EfficientNet}

Trước EfficientNet, việc mở rộng CNN để tăng độ chính xác thường được thực hiện
theo một chiều đơn lẻ: tăng độ sâu (số tầng), tăng độ rộng (số kênh), hoặc tăng độ
phân giải ảnh đầu vào. Tan và Le~\cite{tan2019efficientnet} chỉ ra rằng ba chiều
này không độc lập --- mạng sâu hơn cần độ phân giải cao hơn để tận dụng vùng tiếp
nhận lớn, và cần nhiều kênh hơn để biểu diễn các mẫu tinh vi hơn --- và đề xuất
\emph{mở rộng hợp thành} (\en{compound scaling}), mở rộng đồng thời ba chiều theo
một hệ số duy nhất $\phi$:

\begin{equation}
    \text{độ sâu: } d = \alpha^{\phi}, \qquad
    \text{độ rộng: } w = \beta^{\phi}, \qquad
    \text{độ phân giải: } r = \gamma^{\phi}
    \label{eq:compound_scaling}
\end{equation}

với ràng buộc $\alpha \cdot \beta^{2} \cdot \gamma^{2} \approx 2$ và
$\alpha, \beta, \gamma \geq 1$. Ràng buộc này bảo đảm chi phí tính toán tăng xấp
xỉ $2^{\phi}$ lần khi hệ số hợp thành tăng $\phi$ đơn vị --- một cách kiểm soát
ngân sách tính toán rõ ràng. Các hệ số $\alpha, \beta, \gamma$ được xác định bằng
tìm kiếm lưới trên mạng cơ sở.

Mạng cơ sở EfficientNet-B0 được tìm ra bằng tìm kiếm kiến trúc tự động
(\en{neural architecture search}), với khối xây dựng chính là MBConv --- khối
tích chập nghịch đảo có nút cổ chai (\en{inverted residual bottleneck}) kết hợp
cơ chế \en{squeeze-and-excitation} để hiệu chỉnh trọng số giữa các kênh. Nhờ tích
chập tách theo chiều sâu (\en{depthwise separable convolution}), EfficientNet-B0
chỉ có khoảng \num{5.3} triệu tham số nhưng đạt độ chính xác cao hơn nhiều mạng
lớn hơn nó vài lần, chẳng hạn ResNet-50~\cite{he2016resnet} với khoảng
\num{25} triệu tham số.

Trong luận văn, EfficientNet-B0 giữ vai trò kép, và đây là một lựa chọn thiết kế
có chủ đích chứ không phải sự trùng hợp:

\begin{enumerate}
    \item Làm \textbf{mạng phân loại} ở Giai đoạn 1. Sau khi thay tầng phân loại
          1000 lớp bằng tầng 5 lớp, mô hình có \num{4013953} tham số.
    \item Làm \textbf{bộ mã hoá} (\en{encoder}) cho mạng phân đoạn ở Giai đoạn 4,
          với phần trích đặc trưng gồm \num{4007548} tham số.
\end{enumerate}

Việc dùng chung một kiến trúc xương sống cho cả hai giai đoạn bảo đảm bản đồ chú
ý sinh ra ở Giai đoạn 2 và đặc trưng mà mạng phân đoạn khai thác ở Giai đoạn 4
đến từ cùng một không gian biểu diễn, giúp quy trình nhất quán về mặt khái niệm.

\section{Kiến trúc phân đoạn ảnh}
\label{sec:unet}

\subsection{Từ phân loại đến phân đoạn}

Phân loại ảnh gán một nhãn cho toàn ảnh; phân đoạn ngữ nghĩa
(\en{semantic segmentation}) gán một nhãn cho \emph{từng điểm ảnh}. Chuyển từ bài
toán thứ nhất sang bài toán thứ hai đặt ra một mâu thuẫn cốt lõi: mạng phân loại
liên tục giảm độ phân giải không gian để mở rộng vùng tiếp nhận và trích xuất ngữ
nghĩa, nhưng phân đoạn lại đòi hỏi đầu ra có độ phân giải bằng đúng ảnh gốc. Với
đầu vào $224 \times 224$, bản đồ đặc trưng ở tầng cuối của EfficientNet-B0 chỉ
còn $7 \times 7$ --- mất mát thông tin không gian tới 1024 lần về diện tích.

\subsection{U-Net}

U-Net~\cite{ronneberger2015unet} giải quyết mâu thuẫn trên bằng kiến trúc mã
hoá--giải mã đối xứng hình chữ U. Nhánh mã hoá giảm dần độ phân giải để thu ngữ
nghĩa; nhánh giải mã tăng dần độ phân giải để khôi phục chi tiết không gian. Đóng
góp quyết định là các \emph{kết nối tắt} (\en{skip connection}) nối trực tiếp mỗi
tầng mã hoá với tầng giải mã tương ứng:

\begin{equation}
    x^{i}_{\text{dec}} = \mathcal{F}\big(\,[\,\mathcal{U}(x^{i+1}_{\text{dec}}) \;;\; x^{i}_{\text{enc}}\,]\,\big)
    \label{eq:unet_skip}
\end{equation}

trong đó $\mathcal{U}(\cdot)$ là phép tăng mẫu, $[\,\cdot\,;\,\cdot\,]$ là phép
nối theo chiều kênh và $\mathcal{F}(\cdot)$ là khối tích chập. Nhờ kết nối tắt,
thông tin không gian độ phân giải cao ở nhánh mã hoá được đưa thẳng sang nhánh
giải mã mà không phải đi qua nút cổ chai, cho phép khôi phục biên chính xác. U-Net
ban đầu được thiết kế cho ảnh y sinh --- bài toán cũng có đặc điểm dữ liệu huấn
luyện ít, tương tự bối cảnh của luận văn.

\subsection{UNet++}

Zhou và cộng sự~\cite{zhou2018unetpp} chỉ ra một điểm yếu của U-Net: kết nối tắt
nối trực tiếp đặc trưng \emph{nông} của bộ mã hoá với đặc trưng \emph{sâu} của bộ
giải mã, tạo ra \emph{khoảng cách ngữ nghĩa} (\en{semantic gap}) giữa hai luồng
thông tin có mức trừu tượng rất khác nhau. Việc nối trực tiếp buộc bộ giải mã phải
tự dung hoà hai nguồn không đồng nhất.

\UNetpp{} thu hẹp khoảng cách này bằng các kết nối tắt \emph{lồng nhau, dày đặc}:
thay vì một đường nối thẳng, đặc trưng của bộ mã hoá đi qua một chuỗi khối tích
chập trung gian trước khi tới bộ giải mã. Node $x^{i,j}$ ($i$ là chỉ số tầng, $j$
là chỉ số dọc theo đường nối) được tính:

\begin{equation}
    x^{i,j} =
    \begin{cases}
        \mathcal{F}\big(x^{i-1,j}\big) & \text{nếu } j = 0 \\[6pt]
        \mathcal{F}\Big(\big[\,[x^{i,k}]_{k=0}^{j-1},\; \mathcal{U}(x^{i+1,j-1})\,\big]\Big) & \text{nếu } j > 0
    \end{cases}
    \label{eq:unetpp}
\end{equation}

Mỗi node nhận toàn bộ các node trước đó cùng tầng cộng với đặc trưng đã tăng mẫu
từ tầng dưới. Kết quả là đặc trưng được tinh chỉnh dần qua nhiều bước trung gian
thay vì nhảy một bước, và bộ giải mã nhận được biểu diễn ở nhiều mức trừu tượng
khác nhau.

Đặc điểm này đặc biệt phù hợp với bài toán của luận văn. Bệnh đốm lá tạo ra các
tổn thương \emph{nhiều kích thước khác nhau trong cùng một ảnh}: đốm nhỏ vài chục
điểm ảnh nằm cạnh mảng hoại tử lớn. Kiến trúc lồng nhau cho phép mô hình kết hợp
thông tin đa tỉ lệ tự nhiên hơn so với kết nối tắt đơn.

Chi phí của \UNetpp{} là khiêm tốn. Với cùng bộ mã hoá EfficientNet-B0, mô hình
\UNetpp{} có \num{6569581} tham số so với \num{6251469} tham số của U-Net thường
--- chỉ tăng \num{318112} tham số, tương đương 5,1\%. Trong đó bộ mã hoá chiếm
\num{4007548} tham số và bộ giải mã chiếm \num{2561888} tham số. Mức tăng nhỏ này
là lý do luận văn chọn \UNetpp{} mà không lo ngại về ngân sách tính toán.

\section{Hàm mất mát cho dữ liệu mất cân bằng}
\label{sec:loss}

\subsection{Bản chất của mất cân bằng trong phân đoạn tổn thương}

Trong bài toán phân đoạn nhị phân vùng bệnh, số điểm ảnh thuộc lớp nền lớn hơn
nhiều lần số điểm ảnh thuộc lớp tổn thương. Với độ phủ mục tiêu 18--22\% như
thiết kế ở \cref{chap:phuongphap}, tỉ lệ nền trên tổn thương xấp xỉ 4:1; với các
ảnh bệnh nhẹ, tỉ lệ này còn cao hơn nhiều. Mất cân bằng dẫn tới một nghiệm suy
biến nguy hiểm: mô hình dự đoán toàn bộ ảnh là nền vẫn đạt độ chính xác theo điểm
ảnh rất cao nhưng hoàn toàn vô dụng.

\subsection{Binary Cross-Entropy}

Hàm mất mát cơ bản cho phân loại nhị phân theo từng điểm ảnh:

\begin{equation}
    \mathcal{L}_{\text{BCE}} = -\frac{1}{N}\sum_{i=1}^{N}
        \Big[ y_i \log(p_i) + (1 - y_i)\log(1 - p_i) \Big]
    \label{eq:bce}
\end{equation}

trong đó $y_i \in \{0,1\}$ là nhãn và $p_i \in (0,1)$ là xác suất dự đoán của
điểm ảnh thứ $i$. Hạn chế của BCE là xử lý mọi điểm ảnh với trọng số như nhau:
khi lớp nền áp đảo, tổng đóng góp của các điểm nền lấn át tín hiệu từ vùng tổn
thương.

\subsection{Dice Loss}

Milletari và cộng sự~\cite{milletari2016vnet} đề xuất tối ưu trực tiếp hệ số Dice
--- một đại lượng đo độ chồng lấp, không phụ thuộc vào số điểm ảnh nền:

\begin{equation}
    \mathcal{L}_{\text{Dice}} = 1 - \frac{2\sum_{i} p_i y_i + \epsilon}
                                         {\sum_{i} p_i + \sum_{i} y_i + \epsilon}
    \label{eq:dice}
\end{equation}

Hằng số $\epsilon$ giữ cho biểu thức xác định khi cả dự đoán lẫn nhãn đều rỗng
--- trường hợp thực sự xảy ra trong luận văn, vì \num{683} ảnh lá khoẻ được gán
mặt nạ toàn số không (\cref{chap:phuongphap}). Ưu điểm của Dice là tối ưu trực
tiếp đại lượng dùng để đánh giá; nhược điểm là gradient kém ổn định khi vùng đích
rất nhỏ.

\subsection{Focal Loss}

Lin và cộng sự~\cite{lin2017focal} tiếp cận mất cân bằng theo hướng khác: thay vì
đổi đại lượng tối ưu, họ \emph{giảm trọng số} của các mẫu đã được phân loại tốt:

\begin{equation}
    \mathcal{L}_{\text{Focal}} = -\alpha \,(1 - p_t)^{\gamma} \log(p_t),
    \qquad
    p_t = \begin{cases} p & \text{nếu } y = 1 \\ 1 - p & \text{nếu } y = 0 \end{cases}
    \label{eq:focal}
\end{equation}

Hệ số điều biến $(1-p_t)^{\gamma}$ là mấu chốt. Với một điểm ảnh nền dễ, mô hình
dự đoán $p_t = 0{,}95$, hệ số bằng $0{,}05^{2} = 0{,}0025$ khi $\gamma = 2$ ---
đóng góp của nó bị giảm 400 lần. Ngược lại, một điểm ảnh khó với $p_t = 0{,}5$
chỉ bị giảm 4 lần. Nhờ vậy gradient tập trung vào vùng biên tổn thương --- chính
là nơi mô hình cần học nhất. Tham số $\alpha$ cân bằng thêm giữa hai lớp. Luận văn
dùng $\alpha = 0{,}25$ và $\gamma = 2{,}0$ theo giá trị khuyến nghị của tác giả
gốc.

\subsection{Hàm mất mát kết hợp FocalDice}

Hai hàm trên bổ khuyết cho nhau: Focal Loss hoạt động ở mức \emph{từng điểm ảnh},
điều tiết đóng góp theo độ khó; Dice Loss hoạt động ở mức \emph{toàn vùng}, tối
ưu độ chồng lấp tổng thể. Luận văn sử dụng tổ hợp:

\begin{equation}
    \mathcal{L}_{\text{FocalDice}} = \mathcal{L}_{\text{Focal}} + \mathcal{L}_{\text{Dice}}
    \label{eq:focaldice}
\end{equation}

Cần lưu ý một hệ quả thực nghiệm quan trọng: do thành phần Focal chủ động thu nhỏ
gradient của phần lớn điểm ảnh, giá trị mất mát tổng thể giảm chậm hơn nhiều so
với BCE + Dice. Điều này khiến \emph{tốc độ giảm của hàm mất mát không còn là chỉ
báo đáng tin cậy cho tiến bộ của mô hình} --- một quan sát dẫn tới hệ quả trực
tiếp về tiêu chí dừng sớm, được phân tích ở \cref{chap:thucnghiem}.

\subsection{Thuật toán tối ưu}

Luận văn sử dụng hai bộ tối ưu. Adam~\cite{kingma2015adam} ước lượng thích nghi
mô-men bậc một và bậc hai của gradient, cho tốc độ hội tụ nhanh mà ít cần điều
chỉnh tốc độ học. AdamW~\cite{loshchilov2019adamw} sửa một khiếm khuyết của Adam:
trong Adam, suy giảm trọng số (\en{weight decay}) được cộng vào gradient nên bị
chia cho ước lượng mô-men bậc hai, làm cường độ chính quy hoá thay đổi không kiểm
soát giữa các tham số. AdamW tách suy giảm trọng số khỏi bước cập nhật gradient,
khôi phục đúng ý nghĩa chính quy hoá.

Về lịch tốc độ học, luận văn dùng hai chiến lược khác nhau cho hai nhóm thực
nghiệm. \en{ReduceLROnPlateau} giảm tốc độ học khi chỉ số theo dõi ngừng cải
thiện --- đơn giản và không cần biết trước tổng số bước. \en{CosineAnnealingLR},
xuất phát từ công trình của Loshchilov và Hutter~\cite{loshchilov2017sgdr}, giảm
tốc độ học theo quy luật cô-sin:

\begin{equation}
    \eta_t = \eta_{\min} + \tfrac{1}{2}\big(\eta_{\max} - \eta_{\min}\big)
             \left(1 + \cos\!\Big(\frac{T_{\text{cur}}}{T_{\max}}\pi\Big)\right)
    \label{eq:cosine}
\end{equation}

Lịch cô-sin giảm chậm ở đầu để mô hình khám phá không gian nghiệm, rồi giảm nhanh
về cuối để tinh chỉnh --- và quan trọng là không bao giờ bị kẹt như
\en{ReduceLROnPlateau} khi chỉ số theo dõi dao động.

\section{Trí tuệ nhân tạo khả diễn giải}
\label{sec:xai}

\subsection{Class Activation Map}

Zhou và cộng sự~\cite{zhou2016cam} đề xuất CAM, phương pháp đầu tiên định vị được
vùng ảnh mà mạng phân loại dựa vào để ra quyết định. Với mạng kết thúc bằng gộp
trung bình toàn cục (\en{global average pooling}) rồi tới một tầng tuyến tính,
điểm số của lớp $c$ là $y^{c} = \sum_{k} w^{c}_{k} \cdot \frac{1}{Z}\sum_{i,j} A^{k}_{ij}$.
Hoán đổi thứ tự hai phép tổng cho ra bản đồ định vị:

\begin{equation}
    M^{c}(i, j) = \sum_{k} w^{c}_{k} \, A^{k}_{ij}
    \label{eq:cam}
\end{equation}

Hạn chế nghiêm trọng của CAM là nó \emph{đòi hỏi kiến trúc cụ thể}: mạng bắt buộc
phải có gộp trung bình toàn cục ngay trước tầng phân loại. Mạng nào không thoả
mãn phải sửa kiến trúc rồi huấn luyện lại.

\subsection{Grad-CAM}

Selvaraju và cộng sự~\cite{selvaraju2017gradcam} tổng quát hoá CAM bằng nhận xét
rằng trọng số $w^{c}_{k}$ trong \cref{eq:cam} có thể được thay bằng gradient trung
bình --- một đại lượng tính được ở \emph{bất kỳ} tầng tích chập nào của
\emph{bất kỳ} kiến trúc nào, không cần huấn luyện lại:

\begin{equation}
    \alpha^{c}_{k} = \frac{1}{Z}\sum_{i}\sum_{j} \frac{\partial y^{c}}{\partial A^{k}_{ij}}
    \label{eq:gradcam_alpha}
\end{equation}

\begin{equation}
    L^{c}_{\GradCAM} = \operatorname{ReLU}\!\left(\sum_{k} \alpha^{c}_{k} A^{k}\right)
    \label{eq:gradcam}
\end{equation}

Hàm ReLU ở \cref{eq:gradcam} giữ lại chỉ những vùng có ảnh hưởng \emph{dương} tới
điểm số của lớp quan tâm; vùng có ảnh hưởng âm là bằng chứng chống lại lớp đó nên
bị loại. Tính tổng quát về kiến trúc khiến \GradCAM{} trở thành phương pháp giải
thích được dùng rộng rãi nhất cho CNN.

\subsection{Grad-CAM++}

\GradCAM{} bộc lộ điểm yếu khi một ảnh chứa \emph{nhiều thể hiện} của cùng một
lớp. Do \cref{eq:gradcam_alpha} lấy trung bình gradient trên toàn bản đồ đặc
trưng, một vùng có gradient lớn sẽ chi phối giá trị trung bình và làm các vùng
khác bị san phẳng. Hệ quả là bản đồ nhiệt có xu hướng chỉ làm nổi bật thể hiện rõ
nhất và bỏ sót các thể hiện còn lại.

Đây chính xác là tình huống của bệnh đốm lá: một chiếc lá thường mang hàng chục
đốm bệnh phân tán, mỗi đốm là một thể hiện riêng của cùng lớp bệnh.

Chattopadhay và cộng sự~\cite{chattopadhay2018gradcampp} khắc phục bằng cách thay
trọng số trung bình đồng đều bằng trọng số riêng cho từng vị trí không gian, dựa
trên đạo hàm bậc hai và bậc ba:

\begin{equation}
    \alpha^{kc}_{ij} = \frac{\dfrac{\partial^{2} y^{c}}{\partial (A^{k}_{ij})^{2}}}
    {2\dfrac{\partial^{2} y^{c}}{\partial (A^{k}_{ij})^{2}}
     + \displaystyle\sum_{a}\sum_{b} A^{k}_{ab} \dfrac{\partial^{3} y^{c}}{\partial (A^{k}_{ij})^{3}}}
    \label{eq:gradcampp_alpha}
\end{equation}

\begin{equation}
    w^{c}_{k} = \sum_{i}\sum_{j} \alpha^{kc}_{ij} \cdot
                \operatorname{ReLU}\!\left(\frac{\partial y^{c}}{\partial A^{k}_{ij}}\right)
    \label{eq:gradcampp_w}
\end{equation}

Mỗi vị trí $(i,j)$ nay có trọng số riêng $\alpha^{kc}_{ij}$, nên nhiều vùng kích
hoạt cùng tồn tại trong bản đồ kết quả thay vì bị một vùng lấn át. Đó là lý do
luận văn chuyển từ \GradCAM{} sang \GradCAMpp{} khi nâng cấp nhãn giả từ V1 lên
V2.

\subsection{Hợp nhất đa tỉ lệ}

Bản đồ kích hoạt kế thừa độ phân giải của tầng được chọn, kéo theo một sự đánh
đổi trực tiếp:

\begin{itemize}
    \item Tầng \textbf{sâu} (với đầu vào $224 \times 224$, EfficientNet-B0 cho
          bản đồ $7 \times 7$): giàu ngữ nghĩa, biết chắc ``đây là vùng bệnh'',
          nhưng khi phóng lên $224 \times 224$ thì mỗi ô tương ứng một khối
          $32 \times 32$ điểm ảnh --- biên trở nên rất thô.
    \item Tầng \textbf{nông hơn} ($14 \times 14$): định vị chi tiết hơn gấp bốn
          lần về diện tích, nhưng ngữ nghĩa yếu hơn nên dễ kích hoạt nhầm trên
          kết cấu nền.
\end{itemize}

Cách dung hoà là hợp nhất tuyến tính bản đồ từ hai tầng:

\begin{equation}
    H_{\text{final}} = \lambda \, H_{\text{fine}} + (1 - \lambda) \, H_{\text{coarse}}
    \label{eq:multiscale}
\end{equation}

Luận văn dùng $\lambda = 0{,}6$, thiên về tầng tinh để mặt nạ sắc nét hơn nhưng
vẫn giữ đủ ngữ cảnh ngữ nghĩa từ tầng thô. Lựa chọn cụ thể của các tầng được trình
bày ở \cref{chap:phuongphap}.

\section{Học giám sát yếu và gán nhãn giả}
\label{sec:weaksup}

\subsection{Phân loại các mức giám sát}

Theo mức độ chi tiết của nhãn có sẵn, các bài toán phân đoạn được chia thành:

\begin{itemize}
    \item \textbf{Giám sát đầy đủ}: mỗi ảnh có mặt nạ ở mức điểm ảnh. Cho kết quả
          tốt nhất nhưng chi phí gán nhãn cao nhất.
    \item \textbf{Giám sát yếu} (\en{weakly-supervised}): chỉ có nhãn thô hơn ---
          nhãn mức ảnh, hộp bao, nét vẽ nguệch ngoạc hoặc điểm đánh dấu.
    \item \textbf{Giám sát bán phần} (\en{semi-supervised}): một phần nhỏ có mặt
          nạ đầy đủ, phần lớn không có nhãn.
    \item \textbf{Không giám sát}: hoàn toàn không có nhãn.
\end{itemize}

Luận văn thuộc nhóm thứ hai, ở dạng khó nhất: chỉ có nhãn \emph{mức ảnh}.

\subsection{Gán nhãn giả}

Gán nhãn giả (\en{pseudo-labeling}), được Lee~\cite{lee2013pseudolabel} giới thiệu
cho học bán giám sát, là kỹ thuật dùng chính dự đoán của mô hình làm nhãn huấn
luyện cho vòng học tiếp theo. Trong bối cảnh chuyển từ phân loại sang phân đoạn,
ý tưởng được vận dụng như sau: bản đồ kích hoạt của mạng phân loại được nhị phân
hoá thành mặt nạ, và mặt nạ đó đóng vai trò nhãn giả để huấn luyện mạng phân đoạn.

Chuỗi suy luận có thể tóm tắt: \emph{nhãn mức ảnh $\rightarrow$ mạng phân loại
$\rightarrow$ bản đồ kích hoạt $\rightarrow$ ngưỡng hoá $\rightarrow$ mặt nạ giả
$\rightarrow$ mạng phân đoạn}. Barbedo~\cite{barbedo2019lesions} đã chứng minh
tính khả thi của hướng tiếp cận dựa trên tổn thương này trong lĩnh vực bệnh cây
trồng.

\subsection{Hạn chế cố hữu cần ghi nhận}

Khung phương pháp trên mang một hạn chế bản chất, không thể khắc phục bằng kỹ
thuật, và luận văn coi việc nêu rõ nó là một phần của đóng góp:

\begin{quote}
Bản đồ kích hoạt trả lời câu hỏi \emph{``mô hình nhìn vào đâu để phân loại''}, chứ
không trả lời câu hỏi \emph{``vùng bệnh nằm ở đâu''}. Hai câu hỏi này tương quan
với nhau nhưng không đồng nhất.
\end{quote}

Mạng phân loại chỉ cần \emph{đủ} bằng chứng để phân biệt các lớp, nên nó có thể
tập trung vào một vài đốm điển hình nhất mà bỏ qua phần còn lại --- vẫn phân loại
đúng. Vì vậy nhãn giả sinh ra từ bản đồ kích hoạt thường chưa phủ hết vùng bệnh
thực. Hệ quả kéo theo là mọi chỉ số đánh giá tính trên nhãn giả đều đo mức tái
tạo nhãn giả chứ không đo độ chính xác bệnh học --- vấn đề được phân tích trọn vẹn
ở \cref{sec:caveat}.

\section{Mô hình nền tảng Segment Anything}
\label{sec:sam}

\subsection{Kiến trúc và khả năng}

\SAM{}~\cite{kirillov2023sam} là mô hình nền tảng (\en{foundation model}) cho bài
toán phân đoạn, được Meta AI huấn luyện trên bộ dữ liệu SA-1B gồm khoảng 11 triệu
ảnh và hơn 1 tỉ mặt nạ. Kiến trúc gồm ba thành phần:

\begin{itemize}
    \item \textbf{Bộ mã hoá ảnh}: một Vision Transformer, chạy một lần cho mỗi
          ảnh và cho ra biểu diễn nhúng dùng lại được.
    \item \textbf{Bộ mã hoá gợi ý}: mã hoá gợi ý đầu vào dưới dạng điểm, hộp bao,
          mặt nạ thô hoặc văn bản.
    \item \textbf{Bộ giải mã mặt nạ}: kết hợp hai nguồn nhúng trên để sinh mặt nạ,
          rất nhẹ nên cho phép thử nhiều gợi ý khác nhau với chi phí thấp.
\end{itemize}

Luận văn dùng biến thể ViT-B --- bản nhỏ nhất trong ba biến thể ViT-B/ViT-L/ViT-H,
với tệp trọng số khoảng 375 MB.

\subsection{Vai trò bộ tinh chỉnh nhãn}

Điểm mạnh của \SAM{} là khả năng khái quát không cần huấn luyện lại
(\en{zero-shot}) và chất lượng biên rất cao. Điều này gợi ý một cách dùng khác với
mục đích ban đầu: thay vì dùng \SAM{} để phân đoạn trực tiếp, ta dùng mặt nạ giả
thô làm gợi ý và để \SAM{} tinh chỉnh lại biên. Về nguyên tắc, cách làm này kết
hợp được ưu điểm của cả hai nguồn --- tri thức về \emph{vị trí} bệnh từ bản đồ
kích hoạt, và tri thức về \emph{biên} vật thể từ \SAM{}.

Tuy nhiên, cách dùng này chứa một rủi ro then chốt. \SAM{} được huấn luyện để phân
đoạn \emph{đơn vị ngữ nghĩa tự nhiên} của cảnh. Khi nhận gợi ý là vài điểm nằm
trên một đốm bệnh, đơn vị ngữ nghĩa mà \SAM{} nhận diện có thể không phải là đốm
bệnh mà là \emph{cả chiếc lá} chứa đốm đó. Cơ chế bảo vệ dựa trên IoU thường được
dùng để phát hiện sai lệch, nhưng như phân tích ở \cref{chap:thucnghiem}, cơ chế
này chỉ chặn được trường hợp mặt nạ \emph{trôi} sang vùng khác chứ không chặn được
trường hợp mặt nạ \emph{nở} ra bao trọn mặt nạ ban đầu. Đây là một trong những
phát hiện thực nghiệm chính của luận văn.

\section{Kỹ thuật xử lý ảnh cổ điển}
\label{sec:xulyanh}

Bản đồ kích hoạt là ảnh xám liên tục, trong khi nhãn phân đoạn phải là mặt nạ nhị
phân sạch. Khoảng cách giữa hai dạng biểu diễn được lấp bằng các kỹ thuật xử lý
ảnh cổ điển dưới đây.

\subsection{Ngưỡng hoá Otsu}

Phương pháp Otsu~\cite{otsu1979threshold} chọn ngưỡng tự động bằng cách cực đại
hoá phương sai giữa hai lớp:

\begin{equation}
    \sigma_b^{2}(t) = \omega_0(t)\,\omega_1(t)\,\big[\mu_0(t) - \mu_1(t)\big]^{2}
    \label{eq:otsu}
\end{equation}

với $\omega_0, \omega_1$ là tỉ lệ điểm ảnh và $\mu_0, \mu_1$ là giá trị trung bình
của hai lớp khi cắt tại ngưỡng $t$. Otsu hoạt động tốt khi biểu đồ mức xám có
\emph{hai đỉnh} tách bạch. Giả định này bị vi phạm khi bản đồ nhiệt được hợp nhất
đa tỉ lệ theo \cref{eq:multiscale}: phép hợp nhất làm phân bố trở nên trơn và đơn
đỉnh, khiến Otsu chọn ngưỡng quá thấp và gây phân đoạn thừa. Đây là lý do trực
tiếp khiến luận văn thay Otsu bằng ngưỡng phân vị theo lớp, trình bày ở
\cref{chap:phuongphap}.

\subsection{Phép toán hình thái}

Bốn phép toán hình thái cơ bản trên ảnh nhị phân $A$ với phần tử cấu trúc $B$:

\begin{align}
    \text{Co (\en{erosion}):} \quad & A \ominus B = \{z \mid B_z \subseteq A\} \label{eq:erosion}\\
    \text{Giãn (\en{dilation}):} \quad & A \oplus B = \{z \mid B_z \cap A \neq \emptyset\} \label{eq:dilation}\\
    \text{Mở (\en{opening}):} \quad & A \circ B = (A \ominus B) \oplus B \label{eq:opening}\\
    \text{Đóng (\en{closing}):} \quad & A \bullet B = (A \oplus B) \ominus B \label{eq:closing}
\end{align}

Trong quy trình của luận văn, phép \emph{đóng} lấp các lỗ nhỏ bên trong vùng bệnh
--- cần thiết vì đốm bệnh thường có tâm sáng hơn viền nên bị ngưỡng hoá cắt rỗng
--- còn phép \emph{mở} loại các điểm nhiễu rải rác bên ngoài vùng chính. Hình dạng
phần tử cấu trúc có ý nghĩa thực tế: phần tử \emph{elip} bám sát hình dạng tròn
hoặc bầu dục của đốm bệnh tốt hơn phần tử \emph{chữ nhật}, và đây là một trong các
cải tiến từ nhãn giả V1 sang V2.

\subsection{Lọc thành phần liên thông}

Sau ngưỡng hoá, mặt nạ thường gồm nhiều vùng liên thông rời rạc. Phân tích thành
phần liên thông (\en{connected component analysis}) gán nhãn cho từng vùng và cho
phép lọc theo diện tích hoặc theo thứ hạng. Chiến lược giữ lại $k$ vùng lớn nhất
có diện tích vượt một ngưỡng tối thiểu phù hợp với đặc thù bệnh đốm lá: cần giữ
\emph{nhiều} vùng vì bệnh biểu hiện thành nhiều đốm, nhưng phải loại các mảnh vụn
nhiễu mà phép mở hình thái không xử lý hết.

\subsection{Lọc song phương}

Lọc song phương (\en{bilateral filter}) của Tomasi và
Manduchi~\cite{tomasi1998bilateral} làm trơn ảnh nhưng bảo toàn biên, nhờ nhân
trọng số tính trên cả khoảng cách không gian lẫn chênh lệch cường độ:

\begin{equation}
    I^{\text{filt}}(p) = \frac{1}{W_p}\sum_{q \in \Omega} I(q)\,
        \underbrace{G_{\sigma_s}\big(\lVert p - q \rVert\big)}_{\text{gần về không gian}}\;
        \underbrace{G_{\sigma_r}\big(|I(p) - I(q)|\big)}_{\text{giống về cường độ}}
    \label{eq:bilateral}
\end{equation}

Hai điểm ảnh chỉ ảnh hưởng lẫn nhau khi vừa gần nhau về vị trí vừa tương đồng về
cường độ. Áp dụng lên bản đồ nhiệt, bộ lọc này khiến kích hoạt ``lan'' theo các
vùng đồng màu và dừng lại tại ranh giới màu sắc thực của tổn thương.

\subsection{Trường ngẫu nhiên có điều kiện}

Dense CRF~\cite{krahenbuhl2011densecrf} tinh chỉnh kết quả phân đoạn bằng cách
cực tiểu hoá hàm năng lượng gồm thế đơn (dựa trên dự đoán ban đầu) và thế cặp
(khuyến khích các điểm ảnh gần nhau và giống màu nhận cùng nhãn):

\begin{equation}
    E(\mathbf{x}) = \sum_{i} \psi_u(x_i) + \sum_{i < j} \psi_p(x_i, x_j)
    \label{eq:crf}
\end{equation}

Luận văn có khảo sát Dense CRF như một hướng tinh chỉnh thay thế cho \SAM{}, nhưng
không đưa vào quy trình cuối vì lý do triển khai được nêu ở \cref{chap:thucnghiem}.

\section{Chỉ số đánh giá}
\label{sec:metrics}

\subsection{Chỉ số cho bài toán phân loại}

Từ ma trận nhầm lẫn, với TP, FP, FN lần lượt là số dương tính thật, dương tính giả
và âm tính giả:

\begin{equation}
    \text{Precision} = \frac{TP}{TP + FP}, \qquad
    \text{Recall} = \frac{TP}{TP + FN}, \qquad
    F_1 = \frac{2 \cdot P \cdot R}{P + R}
    \label{eq:prf}
\end{equation}

Với bài toán chẩn đoán bệnh cây trồng, \textbf{recall quan trọng hơn precision}:
bỏ sót một cây bệnh (âm tính giả) gây thiệt hại lan rộng, trong khi cảnh báo nhầm
một cây khoẻ (dương tính giả) chỉ tốn thêm một lần kiểm tra. Nhận định này định
hướng cách diễn giải kết quả ở \cref{chap:thucnghiem}.

Do bộ dữ liệu gần cân bằng (tỉ lệ mất cân bằng 1,33), luận văn báo cáo cả
\en{macro-average} --- trung bình không trọng số qua các lớp, phản ánh hiệu năng
trên lớp ít mẫu --- và \en{weighted-average}.

\subsection{Chỉ số cho bài toán phân đoạn}

Hai chỉ số chuẩn đo độ chồng lấp giữa mặt nạ dự đoán $X$ và mặt nạ tham chiếu $Y$:

\begin{equation}
    \text{IoU} = \frac{|X \cap Y|}{|X \cup Y|} = \frac{TP}{TP + FP + FN}
    \label{eq:iou}
\end{equation}

\begin{equation}
    \text{Dice} = \frac{2|X \cap Y|}{|X| + |Y|} = \frac{2\,TP}{2\,TP + FP + FN}
    \label{eq:dicecoef}
\end{equation}

Hai chỉ số liên hệ với nhau qua một song ánh đơn điệu:

\begin{equation}
    \text{Dice} = \frac{2\,\text{IoU}}{1 + \text{IoU}}, \qquad
    \text{IoU} = \frac{\text{Dice}}{2 - \text{Dice}}
    \label{eq:iou_dice_relation}
\end{equation}

Từ \cref{eq:iou_dice_relation} suy ra hai hệ quả cần lưu ý khi đọc kết quả. Thứ
nhất, Dice luôn lớn hơn hoặc bằng IoU, nên báo cáo Dice thay vì IoU luôn cho con
số ``đẹp'' hơn dù chất lượng không đổi. Thứ hai, do quan hệ là đơn điệu, việc xếp
hạng các mô hình theo IoU hay theo Dice luôn cho cùng thứ tự --- báo cáo cả hai
không cung cấp thêm thông tin về thứ hạng, mà chỉ giúp đối chiếu với các công bố
dùng chỉ số khác nhau.

Trong bài toán phân đoạn nhị phân, $F_1$ tính trên tập điểm ảnh trùng khớp hoàn
toàn với hệ số Dice --- điều này giải thích vì sao hai cột $F_1$ và Dice trong
các bảng kết quả ở \cref{chap:thucnghiem} luôn bằng nhau.

\subsection{Ghi chú về mốc tham chiếu}

Mọi công thức trong mục này đều giả định $Y$ là nhãn chuẩn (\en{ground truth}).
Trong luận văn, do không có nhãn thủ công ở mức điểm ảnh, $Y$ là \emph{nhãn giả}.
Các công thức vẫn tính được và vẫn có ý nghĩa so sánh nội bộ, nhưng ý nghĩa diễn
giải thay đổi hoàn toàn: chúng đo mức độ mô hình tái tạo nhãn giả, chứ không đo
độ chính xác so với tổn thương thực. \Cref{sec:caveat} phân tích trọn vẹn hệ quả
của sự thay đổi này.

\section{Nghiên cứu liên quan}
\label{sec:lienquan}

\subsection{Học sâu trong phát hiện bệnh cây trồng}

Mohanty và cộng sự~\cite{mohanty2016plantdisease} là công trình mở đường, đạt độ
chính xác 99,35\% trên bộ PlantVillage với 54.306 ảnh thuộc 38 lớp bệnh--cây
trồng. Tuy nhiên các tác giả cũng ghi nhận độ chính xác sụt giảm mạnh khi kiểm
thử trên ảnh chụp trong điều kiện thực địa --- do PlantVillage được chụp trên nền
đồng nhất trong phòng thí nghiệm. Đây là bài học nền tảng về khoảng cách giữa dữ
liệu chuẩn hoá và dữ liệu thực tế.

Ferentinos~\cite{ferentinos2018plantdisease} mở rộng khảo sát trên 87.848 ảnh với
25 loài cây, so sánh nhiều kiến trúc CNN và đạt độ chính xác cao nhất 99,53\%.
Kamilaris và Prenafeta-Boldú~\cite{kamilaris2018survey} tổng hợp 40 nghiên cứu
ứng dụng học sâu trong nông nghiệp và kết luận rằng học sâu vượt trội so với các
kỹ thuật xử lý ảnh truyền thống, đồng thời chỉ ra hai điểm nghẽn: thiếu dữ liệu
có nhãn chất lượng và thiếu khả năng diễn giải.

Fuentes và cộng sự~\cite{fuentes2017tomato} tiếp cận theo hướng phát hiện đối
tượng thay vì phân loại, cho phép định vị nhiều vùng bệnh trên cùng một ảnh cà
chua. Barbedo~\cite{barbedo2019lesions} lập luận rằng nhận dạng ở mức
\emph{từng tổn thương} thay vì mức toàn lá giúp tăng đáng kể lượng mẫu huấn luyện
hiệu dụng và cải thiện khả năng khái quát --- lập luận trực tiếp ủng hộ hướng
phân đoạn mà luận văn theo đuổi.

\subsection{Vị trí của luận văn}

Đối chiếu với các công trình trên, luận văn có ba điểm khác biệt:

\begin{enumerate}
    \item \textbf{Đối tượng}: sầu riêng là cây trồng chưa được nghiên cứu nhiều
          trong lĩnh vực phát hiện bệnh bằng học sâu, so với cà chua, lúa hay
          khoai tây.
    \item \textbf{Mức đầu ra}: hầu hết công trình dừng ở phân loại hoặc phát hiện
          bằng hộp bao; luận văn hướng tới phân đoạn ở mức điểm ảnh mà không dùng
          nhãn thủ công nào ở mức đó.
    \item \textbf{Thái độ với kết quả}: thay vì chỉ báo cáo chỉ số cao nhất đạt
          được, luận văn phân tích thẳng vào việc \emph{các chỉ số ấy có nghĩa gì
          và không có nghĩa gì} khi thiếu nhãn chuẩn --- một khía cạnh thường bị
          bỏ qua trong các nghiên cứu học giám sát yếu.
\end{enumerate}

\section{Kết luận chương}

Chương này đã trình bày nền tảng lý thuyết cho toàn bộ quy trình đề xuất:
EfficientNet với cơ chế mở rộng hợp thành làm mạng phân loại và bộ mã hoá;
\UNetpp{} với kết nối tắt lồng nhau làm mạng phân đoạn; tổ hợp Focal và Dice làm
hàm mất mát cho dữ liệu mất cân bằng; \GradCAMpp{} hợp nhất đa tỉ lệ để sinh bản
đồ kích hoạt; \SAM{} trong vai trò bộ tinh chỉnh biên; cùng các kỹ thuật xử lý ảnh
cổ điển để chuyển bản đồ nhiệt thành mặt nạ nhị phân.

Hai điểm lý thuyết sẽ được viện dẫn nhiều lần ở các chương sau, cần được ghi nhận
ngay: (i) bản đồ kích hoạt trả lời câu hỏi ``mô hình nhìn vào đâu'' chứ không phải
``bệnh nằm ở đâu'' (\cref{sec:weaksup}); và (ii) mọi chỉ số IoU/Dice đều được định
nghĩa \emph{tương đối} với một mốc tham chiếu, nên khi mốc tham chiếu thay đổi thì
các con số không còn so sánh trực tiếp được với nhau (\cref{sec:metrics}).

\Cref{chap:phuongphap} trình bày cách các thành phần lý thuyết trên được lắp ghép
thành một quy trình hoàn chỉnh, cùng lý do kỹ thuật của từng lựa chọn thiết kế.

% !TeX root = ../main.tex
\chapter{PHƯƠNG PHÁP ĐỀ XUẤT}
\label{chap:phuongphap}

\section{Tổng quan quy trình}
\label{sec:tongquan}

Quy trình đề xuất gồm bốn giai đoạn nối tiếp, trong đó đầu ra của giai đoạn trước
trở thành đầu vào của giai đoạn sau. Điểm mấu chốt của thiết kế là \emph{không có
giai đoạn nào cần nhãn thủ công ở mức điểm ảnh}: toàn bộ thông tin không gian được
chiết xuất từ nhãn mức ảnh thông qua cơ chế giải thích của mạng phân loại.

\begin{figure}[H]
\centering
\resizebox{\textwidth}{!}{%
\begin{tikzpicture}[node distance=0.85cm and 0.95cm]
    \node[data] (dataset) {Bộ dữ liệu lá sầu riêng\\4.437 ảnh RGB\\5 lớp bệnh};
    \node[data, right=of dataset] (split) {Chia tập 70/10/20\\Train 3.104\\Val 443 / Test 890};
    \node[classmodel, right=of split] (classifier) {Giai đoạn 1\\EfficientNet-B0\\phân loại mức ảnh};
    \node[xai, right=of classifier] (gradcam) {Giai đoạn 2\\GradCAM++ đa tỉ lệ\\$0{,}6H_{\text{fine}}+0{,}4H_{\text{coarse}}$};
    \node[xai, right=of gradcam] (post) {Sinh mặt nạ thích ứng\\ngưỡng phân vị theo lớp\\hình thái + top-3 CC};
    \node[sam, right=of post] (samrefine) {Giai đoạn 3\\Tinh chỉnh bằng SAM\\gợi ý điểm FG/BG\\bảo vệ IoU (min = 0,15)};
    \node[segmodel, right=of samrefine] (segtrain) {Giai đoạn 4\\Phân đoạn UNet++\\bộ mã hoá EfficientNet-B0\\hàm mất mát FocalDice};

    \node[data, below=2.15cm of split] (newimage) {Ảnh lá đầu vào\\$224\times224\times3$};
    \node[classmodel, right=of newimage] (inferclass) {Dự đoán lớp bệnh\\kèm xác suất};
    \node[segmodel, right=of inferclass] (inferseg) {Định vị tổn thương\\mặt nạ nhị phân};
    \node[data, right=of inferseg] (decision) {Kết quả cuối\\lớp bệnh +\\vùng nhiễm bệnh};

    \draw[flow] (dataset) -- (split);
    \draw[flow] (split) -- (classifier);
    \draw[flow] (classifier) -- (gradcam);
    \draw[flow] (gradcam) -- (post);
    \draw[flow] (post) -- (samrefine);
    \draw[flow] (samrefine) -- (segtrain);

    \draw[dashedflow] (classifier.south) -- ++(0,-0.7) -| (inferclass.north);
    \draw[dashedflow] (segtrain.south) -- ++(0,-0.7) -| (inferseg.north);
    \draw[flow] (newimage) -- (inferclass);
    \draw[flow] (inferclass) -- (inferseg);
    \draw[flow] (inferseg) -- (decision);

    \node[lane, fit=(dataset)(split)(classifier)(gradcam)(post)(samrefine)(segtrain), inner xsep=7pt, inner ysep=12pt,
        label={[font=\bfseries\small]above:Huấn luyện và xây dựng nhãn giả}] {};
    \node[lane, fit=(newimage)(inferclass)(inferseg)(decision), inner xsep=7pt, inner ysep=12pt,
        label={[font=\bfseries\small]below:Suy luận}] {};
\end{tikzpicture}%
}
\caption[Tổng quan quy trình học giám sát yếu đề xuất]{Tổng quan quy trình phân
tích bệnh lá sầu riêng theo hướng giám sát yếu. Mô hình phân loại học nhãn mức ảnh
trước, sau đó các giải thích GradCAM++ đa tỉ lệ của nó được chuyển thành mặt nạ
phân đoạn giả thông qua ngưỡng phân vị theo lớp. SAM tinh chỉnh biên của các mặt
nạ này trước khi chúng được dùng để huấn luyện mô hình phân đoạn UNet++ cuối cùng.
Khi suy luận, hệ thống trả về đồng thời lớp bệnh và vùng tổn thương được định vị.}
\label{fig:pipeline}
\end{figure}

Bảng dưới tóm tắt vai trò và sản phẩm của từng giai đoạn.

\begin{table}[H]
\centering
\caption{Vai trò và sản phẩm của bốn giai đoạn}
\label{tab:giaidoan}
\small
\begin{tabular}{L{1.4cm} L{3.6cm} L{3.4cm} L{4.2cm}}
\toprule
\textbf{Giai đoạn} & \textbf{Nhiệm vụ} & \textbf{Đầu vào} & \textbf{Sản phẩm} \\
\midrule
1 & Phân loại bệnh & Ảnh + nhãn mức ảnh & Mô hình phân loại đã huấn luyện \\
2 & Sinh nhãn giả & Mô hình phân loại + 3.104 ảnh huấn luyện & 3.104 mặt nạ nhị phân \\
3 & Tinh chỉnh biên & Mặt nạ giả + ảnh gốc & 3.104 mặt nạ biên sắc nét hơn \\
4 & Phân đoạn & Ảnh + mặt nạ giả & Mô hình phân đoạn đầu cuối \\
\bottomrule
\end{tabular}
\end{table}

Luận văn xây dựng \emph{ba phiên bản} của quy trình, khác nhau ở cách sinh nhãn
giả ở Giai đoạn 2--3, còn Giai đoạn 4 giữ nguyên kiến trúc để bảo đảm tính so
sánh:

\begin{itemize}
    \item \textbf{V1} --- \GradCAM{} một tầng, ngưỡng cố định 0,5
          (\cref{sec:pseudov1}).
    \item \textbf{V2} --- \GradCAMpp{} đa tỉ lệ, ngưỡng phân vị theo lớp
          (\cref{sec:pseudov2}).
    \item \textbf{V3} --- nhãn V2 được \SAM{} tinh chỉnh biên
          (\cref{sec:samrefine}).
\end{itemize}

\section{Bộ dữ liệu và phân tích khám phá}
\label{sec:dulieu}

Phân tích khám phá dữ liệu (\en{Exploratory Data Analysis}, EDA) được thực hiện
trước khi huấn luyện nhằm ba mục đích: xác nhận tính toàn vẹn của dữ liệu, phát
hiện sớm các đặc điểm có thể gây khó khăn cho mô hình, và định hướng lựa chọn chỉ
số đánh giá.

\subsection{Thành phần và phân bố lớp}

Bộ dữ liệu gồm \num{4437} ảnh chia thành năm lớp như trong \cref{tab:phanbolop}.

\begin{table}[H]
\centering
\caption{Phân bố số lượng ảnh theo lớp}
\label{tab:phanbolop}
\begin{tabular}{c l r r}
\toprule
\textbf{Nhãn} & \textbf{Lớp} & \textbf{Số ảnh} & \textbf{Tỉ lệ} \\
\midrule
0 & Algal Leaf Spot (đốm tảo)          & 733 & 16,5\% \\
1 & Allocaridara Attack (rầy phấn)     & 913 & 20,6\% \\
2 & Healthy Leaf (lá khoẻ)             & 976 & 22,0\% \\
3 & Leaf Blight (cháy lá)              & 937 & 21,1\% \\
4 & Phomopsis Leaf Spot (đốm Phomopsis)& 878 & 19,8\% \\
\midrule
\multicolumn{2}{l}{\textbf{Tổng}} & \textbf{4.437} & \textbf{100\%} \\
\bottomrule
\end{tabular}
\end{table}

Nhãn số được gán theo thứ tự bảng chữ cái của tên thư mục, thống nhất giữa mọi
tập con và mọi lần chạy.

\begin{figure}[H]
\centering
\includegraphics[width=\textwidth]{figures/eda_class_distribution.png}
\caption[Phân bố lớp của bộ dữ liệu]{Phân bố số lượng ảnh theo lớp (trái) và tỉ lệ
phần trăm (phải). Đường nét đứt biểu thị giá trị trung bình 887 ảnh mỗi lớp.}
\label{fig:class_distribution}
\end{figure}

Tỉ lệ mất cân bằng --- thương của lớp nhiều nhất trên lớp ít nhất --- bằng
$976/733 = 1{,}33$. Đây là mức rất thấp; theo thông lệ, chỉ khi tỉ lệ vượt khoảng
1:10 mới cần các biện pháp can thiệp như lấy mẫu quá mức hay trọng số theo lớp.
Do đó luận văn \emph{không} áp dụng biện pháp nào và huấn luyện trực tiếp trên
phân bố gốc. Cân bằng gần hoàn hảo cũng khiến \en{macro-F1} trở thành chỉ số tổng
hợp phù hợp, vì không lớp nào đủ áp đảo để bóp méo giá trị trung bình.

\subsection{Phân chia tập dữ liệu}

Bộ dữ liệu được chia theo tỉ lệ 70/10/20 với phân tầng theo lớp.

\begin{table}[H]
\centering
\caption{Phân chia tập dữ liệu cho bài toán phân loại}
\label{tab:split}
\begin{tabular}{l r r L{6.0cm}}
\toprule
\textbf{Tập} & \textbf{Số ảnh} & \textbf{Tỉ lệ} & \textbf{Vai trò} \\
\midrule
Huấn luyện   & 3.104 & 70,0\% & Cập nhật trọng số \\
Kiểm định    & 443   & 10,0\% & Dừng sớm, chọn điểm kiểm tra, điều chỉnh tốc độ học \\
Kiểm thử     & 890   & 20,1\% & Đánh giá cuối, dùng đúng một lần \\
\midrule
\textbf{Tổng} & \textbf{4.437} & \textbf{100\%} & \\
\bottomrule
\end{tabular}
\end{table}

\begin{figure}[H]
\centering
\includegraphics[width=\textwidth]{figures/eda_dataset_splits.png}
\caption[Phân chia tập dữ liệu]{Số lượng ảnh trong ba tập con (trái) và phân bố
lớp trong từng tập (phải). Phân bố lớp đồng đều trên cả ba tập, không có hiện
tượng lệch khi chia dữ liệu.}
\label{fig:dataset_splits}
\end{figure}

Lựa chọn tỉ lệ này phản ánh một cân nhắc cụ thể. Tập kiểm định chỉ chiếm 10\%
(443 ảnh, trung bình khoảng 89 ảnh mỗi lớp) --- vừa đủ để ước lượng hiệu năng phục
vụ dừng sớm --- trong khi phần tiết kiệm được dồn cho tập kiểm thử 20\% (890 ảnh).
Tập kiểm thử lớn cho ước lượng hiệu năng cuối cùng có phương sai thấp hơn, điều
quan trọng khi cần báo cáo một con số duy nhất làm kết quả chính thức.

Phân bố theo lớp trong tập huấn luyện --- bộ ảnh sẽ được dùng để sinh nhãn giả ở
các giai đoạn sau --- như sau: Algal Leaf Spot 513 ảnh, Allocaridara Attack 639,
Healthy Leaf 683, Leaf Blight 655 và Phomopsis Leaf Spot 614.

\subsection{Đặc điểm hình thái của các lớp}

\begin{figure}[H]
\centering
\includegraphics[width=\textwidth]{figures/eda_samples_per_class.png}
\caption[Ảnh mẫu của năm lớp]{Hai ảnh mẫu cho mỗi lớp bệnh. Sự khác biệt hình thái
giữa Algal Leaf Spot và Phomopsis Leaf Spot là nhỏ nhất trong năm lớp.}
\label{fig:samples_per_class}
\end{figure}

\begin{table}[H]
\centering
\caption{Đặc điểm hình thái quan sát được của từng lớp}
\label{tab:hinhthai}
\small
\begin{tabular}{L{3.4cm} L{9.6cm}}
\toprule
\textbf{Lớp} & \textbf{Biểu hiện trên lá} \\
\midrule
Algal Leaf Spot & Đốm tròn hoặc bầu dục màu vàng--nâu--xám, rải rác trên bề mặt lá, viền tương đối rõ. \\
Allocaridara Attack & Vết cắn, đường kẻ và lỗ thủng không đều do côn trùng; kết cấu lá bị phá vỡ mạnh. \\
Healthy Leaf & Lá xanh đồng màu, bề mặt nhẵn, không có tổn thương hay đổi màu. \\
Leaf Blight & Mảng hoại tử lớn màu nâu--đen, thường lan từ rìa lá vào trong, ranh giới không rõ. \\
Phomopsis Leaf Spot & Đốm nhỏ hơn Algal, màu nâu tối viền vàng nhạt, thường tụ thành cụm. \\
\bottomrule
\end{tabular}
\end{table}

Quan sát quan trọng nhất từ \cref{fig:samples_per_class} là cặp \emph{Algal Leaf
Spot} và \emph{Phomopsis Leaf Spot} có hình thái gần nhau: cả hai đều tạo đốm nhỏ
trên nền lá xanh, khác biệt chủ yếu ở kích thước và mật độ đốm. Đây là cặp lớp dự
kiến gây nhầm lẫn nhiều nhất, và dự đoán này sẽ được kiểm chứng bằng ma trận nhầm
lẫn ở \cref{chap:thucnghiem}.

\subsection{Phân tích màu sắc}

\begin{figure}[H]
\centering
\includegraphics[width=\textwidth]{figures/eda_rgb_histogram.png}
\caption[Biểu đồ phân bố kênh màu RGB]{Phân bố ba kênh màu R, G, B cho từng lớp.
Healthy Leaf là lớp duy nhất có kênh Green dịch rõ về phía giá trị cao.}
\label{fig:rgb_histogram}
\end{figure}

\begin{figure}[H]
\centering
\includegraphics[width=\textwidth]{figures/eda_mean_rgb.png}
\caption[Giá trị trung bình các kênh màu theo lớp]{Giá trị trung bình của ba kênh
màu theo từng lớp, tính trên 30 ảnh mỗi lớp.}
\label{fig:mean_rgb}
\end{figure}

\begin{table}[H]
\centering
\caption{Giá trị trung bình các kênh màu theo lớp (mẫu 30 ảnh mỗi lớp)}
\label{tab:meanrgb}
\begin{tabular}{l r r r r}
\toprule
\textbf{Lớp} & \textbf{R} & \textbf{G} & \textbf{B} & \textbf{G $-$ R} \\
\midrule
Algal Leaf Spot     & 95,8  & 125,8 & 106,6 & 30,0 \\
Allocaridara Attack & 94,4  & 134,5 & 98,9  & 40,1 \\
Healthy Leaf        & 96,4  & 133,0 & 92,1  & 36,6 \\
Leaf Blight         & 103,5 & 137,8 & 107,2 & 34,3 \\
Phomopsis Leaf Spot & 97,2  & 138,6 & 88,5  & 41,4 \\
\bottomrule
\end{tabular}
\end{table}

Phân tích màu cho hai nhận định có ý nghĩa thực tiễn.

Thứ nhất, \emph{Algal Leaf Spot có hiệu $G - R$ thấp nhất} (30,0) --- dấu hiệu của
việc mô lá chuyển sang nâu đỏ, làm giảm ưu thế của kênh lục. Đây là đặc trưng màu
sắc phân biệt được lớp này khỏi các lớp còn lại.

Thứ hai, và quan trọng hơn về mặt phương pháp, \emph{các giá trị trung bình toàn
ảnh không đủ để phân tách năm lớp}: khoảng biến thiên của kênh lục giữa các lớp
chỉ từ 125,8 đến 138,6, quá hẹp so với độ lệch trong nội bộ từng lớp. Điều này bác
bỏ khả năng xây dựng một bộ phân loại đơn giản dựa trên đặc trưng màu thủ công, và
củng cố lựa chọn dùng CNN để học đặc trưng \emph{cục bộ} --- vốn là nơi thông tin
bệnh học thực sự nằm.

\subsection{Độ sáng và dung lượng tệp}

\begin{figure}[H]
\centering
\includegraphics[width=\textwidth]{figures/eda_filesize_brightness.png}
\caption[Dung lượng tệp và độ sáng]{Phân bố dung lượng tệp (trái) và phân bố độ
sáng trung bình theo lớp (phải).}
\label{fig:filesize_brightness}
\end{figure}

\begin{figure}[H]
\centering
\includegraphics[width=\textwidth]{figures/eda_brightness_kde.png}
\caption[Mật độ phân bố độ sáng theo lớp]{Ước lượng mật độ hạt nhân của độ sáng
trung bình cho từng lớp.}
\label{fig:brightness_kde}
\end{figure}

Dung lượng tệp tập trung quanh giá trị trung bình 11,36 KB với độ lệch chuẩn chỉ
0,80 KB (nhỏ nhất 8,91 KB, lớn nhất 15,07 KB). Độ đồng đều cao này xác nhận toàn
bộ ảnh đã đi qua cùng một quy trình nén và chuẩn hoá --- nhất quán với nhận định
về trạng thái tiền xử lý của dữ liệu ở \cref{subsec:nguongoc}.

Về hình học, kiểm tra trên mẫu 100 ảnh cho thấy chiều rộng, chiều cao và tỉ lệ
khung hình đều có độ lệch chuẩn bằng 0 --- mọi ảnh đúng $224 \times 224$ và vuông.
Hệ quả trực tiếp: bước thay đổi kích thước trong quy trình tiền xử lý trở thành
phép đồng nhất, và không có biến dạng hình học nào được đưa thêm vào dữ liệu.

\subsection{Kiểm chứng luồng nạp dữ liệu}

Trước khi huấn luyện, luồng nạp dữ liệu được kiểm chứng bằng cách lấy một lô mẫu
và đối chiếu các đại lượng thống kê. Kết quả: kích thước lô $[32, 3, 224, 224]$
đúng như thiết kế; giá trị điểm ảnh sau chuẩn hoá nằm trong khoảng
$[-2{,}118;\ 2{,}640]$; nhãn nằm trong $[0, 4]$; độ lệch chuẩn theo từng kênh xấp
xỉ 1,0 (\num{1.0298}, \num{1.0201}, \num{1.0387}).

Giá trị trung bình theo kênh sau chuẩn hoá là $(-0{,}618;\ 0{,}001;\ -0{,}320)$
thay vì $(0, 0, 0)$. Đây \emph{không} phải lỗi mà là hệ quả có thể dự đoán được:
phép chuẩn hoá dùng giá trị trung bình và độ lệch chuẩn của ImageNet, trong khi
ảnh lá cây có kênh lục trội hơn và kênh đỏ, lam yếu hơn so với phân bố ImageNet.
Độ lệch chuẩn theo kênh xấp xỉ 1,0 cho thấy phép chuẩn hoá vẫn phát huy đúng tác
dụng chính là đưa dữ liệu về thang đo phù hợp cho tối ưu bằng gradient.

\section{Giai đoạn 1 --- Mô hình phân loại}
\label{sec:giaidoan1}

\subsection{Kiến trúc}

Mô hình phân loại dùng EfficientNet-B0 tiền huấn luyện trên ImageNet, thay tầng
phân loại 1000 lớp gốc bằng cấu trúc:

\begin{equation}
    \text{Dropout}(p = 0{,}2) \;\rightarrow\; \text{Linear}(1280 \rightarrow 5)
    \label{eq:cls_head}
\end{equation}

Toàn bộ mô hình có \num{4013953} tham số và được tinh chỉnh hoàn toàn --- không
tầng nào bị đóng băng --- theo biện luận ở \cref{sec:cnn}.

\subsection{Tăng cường dữ liệu}

Tăng cường dữ liệu (\en{data augmentation}) được áp dụng \emph{chỉ trên tập huấn
luyện}: lật ngang, lật dọc, xoay và biến đổi màu (\en{color jitter}). Tập kiểm
định và kiểm thử không được tăng cường, vì nếu tăng cường thì mỗi lần đánh giá sẽ
cho kết quả khác nhau, làm nhiễu chính tín hiệu dùng để dừng sớm.

Lật dọc là phép biến đổi hợp lệ trong bài toán này --- khác với ảnh chữ viết hay
ảnh phong cảnh --- vì một chiếc lá không có hướng ``đúng'' cố định khi chụp.

\subsection{Cấu hình huấn luyện}

\begin{table}[H]
\centering
\caption{Cấu hình huấn luyện mô hình phân loại}
\label{tab:cauhinh_cls}
\begin{tabular}{L{4.6cm} L{8.4cm}}
\toprule
\textbf{Thành phần} & \textbf{Thiết lập} \\
\midrule
Kiến trúc          & EfficientNet-B0, tiền huấn luyện ImageNet \\
Tầng phân loại     & Dropout(0,2) + Linear(1280 $\rightarrow$ 5) \\
Số tham số         & 4.013.953 (toàn bộ khả huấn luyện) \\
Hàm mất mát        & Cross-Entropy \\
Bộ tối ưu          & Adam, tốc độ học $1 \times 10^{-4}$ \\
Lịch tốc độ học    & ReduceLROnPlateau theo \texttt{val\_acc}, patience = 2, hệ số 0,5 \\
Kích thước lô      & 16 \\
Số epoch tối đa    & 50 \\
Dừng sớm           & patience = 5 theo \texttt{val\_acc} \\
Tăng cường         & Lật ngang/dọc, xoay, biến đổi màu \\
\bottomrule
\end{tabular}
\end{table}

Hai lựa chọn cần biện luận:

\textbf{Tốc độ học $10^{-4}$.} Giá trị này nhỏ hơn một bậc so với mặc định
$10^{-3}$ của Adam. Với tinh chỉnh toàn bộ, tốc độ học lớn khiến các bước cập nhật
đầu tiên --- khi tầng phân loại mới còn khởi tạo ngẫu nhiên và sinh gradient lớn
--- phá huỷ chính các trọng số tiền huấn luyện mà ta muốn tận dụng.

\textbf{Kích thước lô 16.} Lô nhỏ tạo nhiễu gradient, mà nhiễu này có tác dụng
chính quy hoá nhẹ --- hữu ích với tập huấn luyện chỉ hơn ba nghìn ảnh.

\subsection{Chiến lược dừng sớm}

Cần phân biệt rõ hai cơ chế cùng theo dõi \texttt{val\_acc} nhưng phục vụ hai mục
đích khác nhau: \en{ReduceLROnPlateau} với patience = 2 \emph{giảm tốc độ học} khi
hiệu năng chững lại, còn dừng sớm với patience = 5 \emph{kết thúc huấn luyện}. Việc
đặt patience của bộ lập lịch nhỏ hơn của dừng sớm là có chủ đích: mô hình được cho
ít nhất hai cơ hội giảm tốc độ học để thoát vùng phẳng trước khi bị dừng hẳn.

Điểm kiểm tra chỉ được ghi đè khi \texttt{val\_acc} vượt kỷ lục cũ, nên tệp
\texttt{best\_\allowbreak model.\allowbreak pth} luôn ứng với mô hình tốt nhất từng đạt được. Đây chính là
tệp được nạp lại ở tất cả các giai đoạn sau.

\section{Giai đoạn 2 --- Sinh nhãn giả phiên bản V1}
\label{sec:pseudov1}

\subsection{Trích xuất bản đồ nhiệt}

Nhãn giả V1 dùng \GradCAM{} một tầng theo \cref{eq:gradcam}, với tầng đích là
\texttt{base.\allowbreak features.\allowbreak 8.\allowbreak 0} --- tầng tích chập cuối cùng của EfficientNet-B0, cho
bản đồ đặc trưng $7 \times 7$. Đây là lựa chọn tiêu chuẩn: tầng sâu nhất giữ nhiều
ngữ nghĩa nhất trong khi vẫn còn thông tin không gian.

Quy trình được áp dụng cho \emph{toàn bộ} \num{3104} ảnh huấn luyện. Phép biến đổi
đầu vào giống hệt tập kiểm định và kiểm thử, không tăng cường --- yêu cầu bắt buộc
để mặt nạ sinh ra ổn định giữa các lần chạy.

\subsection{Nhị phân hoá và hậu xử lý hình thái}

Bản đồ nhiệt đã chuẩn hoá về $[0, 1]$ được nhị phân hoá bằng ngưỡng cố định:

\begin{equation}
    M(i,j) = \begin{cases}
        1 & \text{nếu } H(i,j) \geq 0{,}5 \\
        0 & \text{nếu ngược lại}
    \end{cases}
    \label{eq:threshold_v1}
\end{equation}

Sau đó áp dụng lần lượt phép \emph{đóng} rồi phép \emph{mở} theo
\cref{eq:closing,eq:opening} với phần tử cấu trúc \emph{chữ nhật} $5 \times 5$.
Kích thước này tương ứng khoảng 0,05\% diện tích ảnh --- đủ để làm mịn biên mà
không xoá mất các đốm bệnh nhỏ.

\begin{figure}[H]
\centering
\includegraphics[width=\textwidth]{figures/pl1_morphology_demo.png}
\caption[Minh hoạ hậu xử lý hình thái]{Tác động của hậu xử lý hình thái. Từ trái
sang phải: ảnh gốc, bản đồ nhiệt Grad-CAM, mặt nạ chỉ qua ngưỡng hoá, và mặt nạ
sau khi áp dụng phép đóng rồi phép mở với phần tử cấu trúc $5 \times 5$.}
\label{fig:morphology_demo}
\end{figure}

Toàn bộ quá trình sinh \num{3104} mặt nạ mất 14 phút 00 giây trên GPU, tương đương
tốc độ 3,69 ảnh mỗi giây.

\subsection{Ba hạn chế của phiên bản V1}
\label{subsec:hanche_v1}

Phân tích thống kê mặt nạ V1 bộc lộ ba vấn đề, mỗi vấn đề dẫn tới một cải tiến cụ
thể ở phiên bản V2.

\textbf{Thứ nhất --- lớp lá khoẻ sinh mặt nạ vô nghĩa.} Độ phủ trung bình của lớp
Healthy Leaf đạt 26,4\%, cao gần gấp đôi các lớp bệnh (12,3--15,1\%), với giá trị
lớn nhất lên tới 59,1\%. Điều này hoàn toàn logic mà lại rất tai hại: lá khoẻ
\emph{không có} tổn thương nào để mạng tập trung, nên bản đồ chú ý khuếch tán khắp
phiến lá; ngưỡng hoá một bản đồ khuếch tán tất yếu cho vùng rộng. Nếu đưa
\num{683} mặt nạ này vào huấn luyện phân đoạn, mô hình sẽ học rằng lá khoẻ cũng có
26\% diện tích ``bệnh''.

\textbf{Thứ hai --- ngưỡng cố định không thích ứng.} Ngưỡng 0,5 áp dụng cho mọi
ảnh, bất kể phân bố bản đồ nhiệt của từng ảnh. Ảnh bệnh nhẹ bị phân đoạn thiếu,
ảnh có nền nổi bật bị phân đoạn thừa. Hệ quả là độ lệch chuẩn của độ phủ rất lớn
(0,039--0,076 tuỳ lớp).

\textbf{Thứ ba --- biên quá thô.} Bản đồ $7 \times 7$ phóng lên $224 \times 224$
khiến mỗi ô tương ứng một khối $32 \times 32$ điểm ảnh. Các đốm bệnh rời rạc bị
gộp thành một khối lớn duy nhất, làm mất thông tin về số lượng và vị trí từng đốm
--- đúng như phân tích lý thuyết ở \cref{sec:xai}.

\section{Giai đoạn 2' --- Sinh nhãn giả phiên bản V2}
\label{sec:pseudov2}

Phiên bản V2 khắc phục cả ba hạn chế trên bằng bốn thay đổi kỹ thuật, tóm tắt ở
\cref{tab:v1_v2}.

\begin{table}[H]
\centering
\caption{So sánh thiết kế nhãn giả V1 và V2}
\label{tab:v1_v2}
\small
\begin{tabular}{L{2.6cm} L{3.5cm} L{3.5cm} L{3.6cm}}
\toprule
\textbf{Thành phần} & \textbf{V1} & \textbf{V2} & \textbf{Lý do thay đổi} \\
\midrule
Phương pháp CAM & Grad-CAM & Grad-CAM++ & Tránh hiệu ứng san phẳng khi có nhiều đốm bệnh \\
Tầng đích & \texttt{features.8} ($7\times7$) & \texttt{features.6} ($14\times14$) + \texttt{features.8} & Tăng chi tiết không gian gấp bốn lần về diện tích \\
Ngưỡng hoá & Cố định 0,5 & Phân vị theo lớp & Kiểm soát trực tiếp độ phủ \\
Phần tử cấu trúc & Chữ nhật $5\times5$ & Elip $5\times5$ & Bám hình dạng tròn của đốm bệnh \\
Lọc nhiễu & Không & Top-3 CC $\geq$ 400 px & Loại mảnh vụn mà phép mở bỏ sót \\
Lá khoẻ & Xử lý như lớp bệnh & Mặt nạ toàn số 0 & Loại 683 mặt nạ dương tính giả \\
\bottomrule
\end{tabular}
\end{table}

\subsection{Grad-CAM++ đa tỉ lệ}

Hai tầng đích được chọn theo \cref{eq:multiscale}:

\begin{itemize}
    \item \textbf{Tầng tinh}: \texttt{base.\allowbreak features.\allowbreak 6.\allowbreak 3.\allowbreak block.\allowbreak 3.\allowbreak 0}, độ phân giải
          $14 \times 14$ --- nắm bắt các đốm nhỏ.
    \item \textbf{Tầng thô}: \texttt{base.\allowbreak features.\allowbreak 8.\allowbreak 0}, độ phân giải
          $7 \times 7$ --- cung cấp ngữ cảnh ngữ nghĩa.
\end{itemize}

Bản đồ nhiệt hợp nhất:

\begin{equation}
    H_{\text{final}} = 0{,}6 \cdot H_{\text{fine}} + 0{,}4 \cdot H_{\text{coarse}}
    \label{eq:fusion_v2}
\end{equation}

Trọng số thiên về tầng tinh nhằm ưu tiên độ sắc nét của biên, nhưng vẫn giữ 40\%
đóng góp từ tầng thô để tránh kích hoạt nhầm trên kết cấu nền. Trọng số của từng
tầng được tính theo \cref{eq:gradcampp_alpha,eq:gradcampp_w}.

\subsection{Ngưỡng phân vị theo lớp}

Đây là cải tiến trung tâm của V2. Thay vì cố định giá trị ngưỡng, ta cố định
\emph{tỉ lệ diện tích} được giữ lại. Với tỉ lệ mục tiêu $p$, ngưỡng được tính là
phân vị thứ $(1 - p) \times 100$ của chính bản đồ nhiệt đó:

\begin{equation}
    \tau = \operatorname{percentile}\big(H,\; (1 - p) \times 100\big),
    \qquad
    M(i,j) = \mathbf{1}\big[H(i,j) \geq \tau\big]
    \label{eq:percentile}
\end{equation}

Cách làm này đảo ngược quan hệ nhân quả so với V1: ở V1 ngưỡng là tham số và độ
phủ là kết quả không kiểm soát được; ở V2 độ phủ là tham số còn ngưỡng được suy ra
cho từng ảnh. Nhờ đó ngưỡng tự thích ứng với phân bố riêng của từng bản đồ nhiệt.

\begin{table}[H]
\centering
\caption{Tỉ lệ độ phủ mục tiêu theo lớp}
\label{tab:keeppct}
\begin{tabular}{l c L{7.2cm}}
\toprule
\textbf{Lớp} & \textbf{Giữ lại} & \textbf{Căn cứ} \\
\midrule
Algal Leaf Spot     & 18\% & Đốm nhỏ rải rác, cần độ chính xác cao \\
Allocaridara Attack & 20\% & Tổn thương côn trùng trên vùng rộng hơn \\
Healthy Leaf        & 0\%  & Không có mô tổn thương --- ép về không \\
Leaf Blight         & 22\% & Mảng hoại tử lan rộng nhất \\
Phomopsis Leaf Spot & 18\% & Đốm nhỏ, tương tự Algal \\
\bottomrule
\end{tabular}
\end{table}

Các giá trị 18--22\% được chọn dựa trên đặc điểm bệnh học quan sát từ
\cref{tab:hinhthai}, theo thứ tự: bệnh đốm chiếm diện tích nhỏ nhất, bệnh cháy lá
lan rộng nhất. Cần nói rõ đây là \emph{tham số suy luận, chưa được hiệu chỉnh dựa
trên nhãn chuẩn} --- một hạn chế đã nêu ở \cref{subsec:gioihan-quantrong}.

Lý do không dùng ngưỡng Otsu đã được phân tích ở \cref{sec:xulyanh}: phép hợp nhất
đa tỉ lệ theo \cref{eq:fusion_v2} làm bản đồ nhiệt trở nên trơn và đơn đỉnh, vi
phạm giả định hai đỉnh mà Otsu dựa vào.

\subsection{Ép về không cho lớp lá khoẻ}

Với lớp Healthy Leaf, tỉ lệ giữ lại được đặt bằng 0, tức mặt nạ là ma trận không:

\begin{equation}
    M_{\text{healthy}}(i,j) = 0 \quad \forall (i,j)
    \label{eq:hardzero}
\end{equation}

Đây là một \emph{tri thức tiên nghiệm} được đưa thẳng vào quy trình: lá khoẻ theo
định nghĩa không có tổn thương, nên mọi kích hoạt trên ảnh lá khoẻ đều là dương
tính giả xét theo mục tiêu phân đoạn bệnh. Biện pháp này loại bỏ \num{683} mặt nạ
nhiễu và cung cấp cho mô hình phân đoạn các mẫu âm sạch --- điều mà bản thân bản đồ
kích hoạt không bao giờ tự cho được.

\subsection{Hậu xử lý}

Hai bước nối tiếp sau ngưỡng hoá:

\begin{enumerate}
    \item \textbf{Hình thái với phần tử elip $5 \times 5$}: đóng rồi mở. Phần tử
          elip bám sát hình dạng tròn hoặc bầu dục của đốm bệnh tốt hơn phần tử
          chữ nhật dùng ở V1.
    \item \textbf{Lọc top-3 thành phần liên thông}, chỉ giữ các vùng có diện tích
          $\geq 400$ điểm ảnh.
\end{enumerate}

Hai tham số của bước lọc phản ánh đặc thù bài toán. Giữ \emph{ba} vùng thay vì một
vì bệnh đốm lá biểu hiện thành nhiều tổn thương tách biệt --- giữ một vùng sẽ làm
mất phần lớn thông tin. Ngưỡng \emph{400 điểm ảnh}, tương đương khoảng 0,8\% diện
tích ảnh, đủ lớn để loại các mảnh vụn cỡ trung mà phép mở hình thái không xử lý
hết, nhưng vẫn nhỏ hơn một đốm bệnh điển hình.

\begin{figure}[H]
\centering
\includegraphics[width=\textwidth]{figures/pl2_v1_vs_v2.png}
\caption[So sánh mặt nạ V1 và V2]{So sánh chất lượng mặt nạ trên cùng một ảnh. Từ
trái sang phải: ảnh gốc, bản đồ nhiệt GradCAM++ đa tỉ lệ, mặt nạ V1 (ngưỡng cố
định 0,5, phần tử chữ nhật) và mặt nạ V2 (ngưỡng phân vị, phần tử elip, lọc top-3
thành phần liên thông).}
\label{fig:v1_vs_v2}
\end{figure}

Toàn bộ quá trình sinh nhãn V2 mất 4 phút 08 giây (12,51 ảnh mỗi giây) --- nhanh
hơn V1 khoảng 3,4 lần dù tính hai bản đồ kích hoạt thay vì một, nhờ cách cài đặt
hiệu quả hơn cho phép truyền ngược.

\subsection{Đánh đổi của thiết kế V2}
\label{subsec:tradeoff_v2}

Ngưỡng phân vị bảo đảm độ phủ bám sát mục tiêu, nhưng cái giá phải trả cần được
nêu thẳng: \emph{độ phủ trở nên gần như đồng nhất trong cùng một lớp}. Một ảnh chỉ
có một đốm bệnh nhỏ và một ảnh bị bệnh phủ kín lá đều nhận mặt nạ chiếm xấp xỉ
cùng tỉ lệ diện tích.

Nói cách khác, mặt nạ V2 \emph{không mã hoá được mức độ nặng nhẹ của bệnh}. Mô
hình phân đoạn huấn luyện trên nhãn này học cách ``phủ đúng diện tích'' hơn là học
cách ``phát hiện vùng bệnh thực sự''. Đây là đánh đổi được chấp nhận có ý thức để
loại bỏ hiện tượng phân đoạn thừa của Otsu, và hệ quả của nó sẽ hiện rõ trong kết
quả thực nghiệm ở \cref{chap:thucnghiem}.

\section{Giai đoạn 3 --- Tinh chỉnh biên bằng SAM}
\label{sec:samrefine}

\subsection{Động lực}

Mặt nạ V2 định vị đúng vùng bệnh nhưng biên vẫn thô, vì bản đồ nhiệt gốc chỉ có độ
phân giải $14 \times 14$ và $7 \times 7$. \SAM{} được kỳ vọng bổ khuyết đúng điểm
này: giữ nguyên thông tin \emph{vị trí} từ V2 và thay thế thông tin \emph{biên}
bằng biên chính xác ở mức điểm ảnh mà \SAM{} sinh ra.

\subsection{Sinh gợi ý}

Với mỗi mặt nạ V2 khác rỗng, các gợi ý được trích xuất như sau:

\begin{itemize}
    \item \textbf{Điểm tiền cảnh} (5 điểm): trọng tâm của mặt nạ, cộng 4 điểm lấy
          ngẫu nhiên trong vùng mặt nạ. Việc bổ sung điểm ngẫu nhiên là cần thiết
          --- nếu chỉ dùng trọng tâm thì với mặt nạ hình vành khuyên hoặc gồm
          nhiều cụm rời, trọng tâm có thể rơi ra ngoài vùng thực.
    \item \textbf{Điểm nền} (3 điểm): lấy ngẫu nhiên trong phần bù của mặt nạ đã
          giãn nở 30 điểm ảnh. Vùng đệm 30 điểm ảnh bảo đảm điểm nền không rơi sát
          biên tổn thương --- nơi nhãn vốn không chắc chắn.
\end{itemize}

Mặt nạ rỗng --- toàn bộ \num{683} ảnh lá khoẻ --- được bỏ qua và giữ nguyên trạng
thái rỗng.

\subsection{Cơ chế bảo vệ dựa trên IoU}

Để phòng trường hợp \SAM{} chọn nhầm vùng, một cơ chế bảo vệ được cài đặt:

\begin{equation}
    M_{\text{refined}} =
    \begin{cases}
        M_{\SAM} & \text{nếu } \operatorname{IoU}(M_{\SAM},\, M_{\text{V2}}) \geq 0{,}15 \\[4pt]
        M_{\text{V2}} & \text{nếu ngược lại}
    \end{cases}
    \label{eq:iou_guard}
\end{equation}

Ý tưởng là nếu mặt nạ \SAM{} sinh ra chồng lấp quá ít với mặt nạ V2 thì nhiều khả
năng \SAM{} đã chuyển sang một đối tượng khác, khi đó ta quay về dùng mặt nạ V2.

Cần lưu ý ngay rằng cơ chế này \emph{chỉ chặn được một trong hai kiểu sai lệch}.
Phân tích đầy đủ về khiếm khuyết của nó, cùng bằng chứng định lượng, được trình
bày ở \cref{chap:thucnghiem} --- đây là một trong những phát hiện chính của luận
văn.

\subsection{Hậu xử lý}

Mặt nạ sau tinh chỉnh được làm sạch bằng hai bước: loại các thành phần liên thông
có diện tích nhỏ hơn 5\% tổng diện tích mặt nạ hoặc nhỏ hơn 30 điểm ảnh, sau đó
áp dụng phép đóng hình thái $5 \times 5$ để lấp các lỗ nhỏ.

\section{Giai đoạn 4 --- Mô hình phân đoạn}
\label{sec:giaidoan4}

\subsection{Kiến trúc}

Phiên bản V1 dùng EfficientNet-UNet tự cài đặt: bộ mã hoá EfficientNet-B0 trích
đặc trưng tại bốn mức, bộ giải mã U-Net với kết nối tắt trực tiếp, tổng cộng
\num{7276925} tham số (bộ mã hoá \num{4007548}, bộ giải mã \num{3269312}).

Hai phiên bản V2 và V3 dùng \UNetpp{} với bộ mã hoá EfficientNet-B0
(\num{6569581} tham số) theo biện luận ở \cref{sec:unet}. Đầu ra là một kênh
\en{logit}, hàm \en{sigmoid} được đưa vào trong hàm mất mát thay vì đặt ở đầu ra
mô hình --- cách làm chuẩn để bảo đảm ổn định số học.

\subsection{Tăng cường dữ liệu cho phân đoạn}

Khác với phân loại, tăng cường cho phân đoạn phải biến đổi \emph{đồng bộ} cả ảnh
lẫn mặt nạ. Phiên bản V1 chỉ dùng lật ngang và lật dọc. Hai phiên bản V2 và V3 dùng
bộ tăng cường mạnh hơn nhiều qua thư viện
Albumentations~\cite{buslaev2020albumentations}, liệt kê ở \cref{tab:augment}.

\begin{table}[H]
\centering
\caption{Bộ tăng cường dữ liệu cho phân đoạn (V2 và V3)}
\label{tab:augment}
\small
\begin{tabular}{L{4.2cm} c L{7.2cm}}
\toprule
\textbf{Phép biến đổi} & \textbf{Xác suất} & \textbf{Mục đích} \\
\midrule
RandomResizedCrop (0,7--1,0) & 1,00 & Bất biến với tỉ lệ và vị trí tổn thương \\
HorizontalFlip / VerticalFlip & 0,50 & Bất biến với hướng đặt lá \\
RandomRotate90 & 0,50 & Bất biến với góc quay vuông \\
ShiftScaleRotate & 0,50 & Bất biến với xoay nhỏ và tịnh tiến \\
ElasticTransform & 0,30 & Mô phỏng lá cong vênh trong thực tế \\
GridDistortion & 0,20 & Mô phỏng biến dạng phối cảnh \\
ColorJitter & 0,50 & Bất biến với điều kiện ánh sáng \\
GaussNoise & 0,20 & Bền vững với nhiễu cảm biến \\
\bottomrule
\end{tabular}
\end{table}

Hai phép biến đổi đàn hồi và biến dạng lưới đáng chú ý vì chúng làm biến dạng
\emph{cả mặt nạ}. Với nhãn chuẩn, điều này chấp nhận được; với nhãn giả vốn đã
không chắc chắn về biên, việc biến dạng thêm có thể làm nhiễu tín hiệu huấn luyện.
Đây là một điểm cần cân nhắc, được thảo luận lại ở \cref{chap:thucnghiem}.

\subsection{Cấu hình của ba phiên bản}

\begin{table}[H]
\centering
\caption{Cấu hình huấn luyện ba phiên bản mô hình phân đoạn}
\label{tab:cauhinh_seg}
\small
\begin{tabular}{L{3.2cm} L{3.2cm} L{3.2cm} L{3.2cm}}
\toprule
\textbf{Thành phần} & \textbf{V1} & \textbf{V2} & \textbf{V3} \\
\midrule
Kiến trúc & EfficientNet-UNet & UNet++ EB0 & UNet++ EB0 \\
Số tham số & 7.276.925 & 6.569.581 & 6.569.581 \\
Nhãn huấn luyện & Pseudo V1 & Pseudo V2 & SAM tinh chỉnh \\
Hàm mất mát & BCE + Dice & FocalDice & FocalDice \\
Bộ tối ưu & Adam $10^{-4}$ & AdamW $10^{-4}$, wd $10^{-4}$ & AdamW $10^{-4}$, wd $10^{-4}$ \\
Lịch tốc độ học & ReduceLROnPlateau & CosineAnnealing $T_{\max}=50$ & CosineAnnealing $T_{\max}=50$ \\
Kích thước lô & 8 & 4 & 8 \\
Epoch tối đa & 30 & 50 & 50 \\
Dừng sớm & patience 7 & patience 5 & patience 10 \\
Tiêu chí dừng & \texttt{val\_loss} & \texttt{val\_loss} & \texttt{val\_loss} \\
Cắt gradient & Không & \texttt{max\_norm} = 1,0 & \texttt{max\_norm} = 1,0 \\
Tăng cường & Chỉ lật & Albumentations & Albumentations \\
\bottomrule
\end{tabular}
\end{table}

Thiết kế thực nghiệm được tổ chức sao cho \textbf{V2 và V3 chỉ khác nhau ở một
biến duy nhất là bộ nhãn huấn luyện}. Kiến trúc, hàm mất mát, bộ tối ưu, lịch tốc
độ học, bộ tăng cường và cách chia tập đều giữ nguyên. Nhờ vậy, chênh lệch hiệu
năng giữa hai phiên bản có thể quy cho tác động của việc tinh chỉnh biên bằng
\SAM{} --- đây là phép so sánh có kiểm soát duy nhất trong luận văn.

Hai khác biệt còn lại giữa V2 và V3 --- kích thước lô 4 so với 8 và patience 5 so
với 10 --- xuất phát từ ràng buộc kỹ thuật khi chạy thực nghiệm chứ không phải lựa
chọn thiết kế, và ảnh hưởng của chúng được đánh giá ở \cref{chap:thucnghiem}.

\subsection{Chia tập cho bài toán phân đoạn}
\label{subsec:split_seg}

Toàn bộ \num{3104} ảnh huấn luyện kèm nhãn giả được chia lại theo tỉ lệ 70/15/15,
cho \num{2172} ảnh huấn luyện, 465 ảnh kiểm định và 467 ảnh kiểm thử, với hạt
giống ngẫu nhiên cố định bằng 42.

Cần lưu ý một chi tiết kỹ thuật có hệ quả trực tiếp tới việc diễn giải kết quả:

\begin{quote}
Phiên bản V1 chia tập bằng \texttt{torch.\allowbreak random\_\allowbreak split} với bộ sinh số ngẫu nhiên
của PyTorch, trong khi V2 và V3 chia bằng
\texttt{numpy.\allowbreak random.\allowbreak default\_rng(42).\allowbreak
permutation}. Dù cùng hạt giống 42, hai bộ
sinh số ngẫu nhiên này tạo ra \emph{hai hoán vị khác nhau}, nên \textbf{tập kiểm
thử của V1 không trùng với tập kiểm thử của V2 và V3}.
\end{quote}

Đây là một trong nhiều lý do khiến chỉ số của V1 không so sánh trực tiếp được với
V2 và V3, bên cạnh lý do quan trọng hơn về khác biệt mốc tham chiếu, được phân
tích trọn vẹn ở \cref{sec:caveat}. Ngược lại, V2 và V3 dùng chung đúng một hoán vị
nên tập kiểm thử của chúng trùng khớp hoàn toàn.

\section{Kết luận chương}

Chương này đã trình bày chi tiết quy trình bốn giai đoạn cùng ba phiên bản nhãn
giả. Các quyết định thiết kế chính và căn cứ của chúng được tóm tắt như sau:

\begin{itemize}
    \item Dùng chung EfficientNet-B0 cho cả phân loại và bộ mã hoá phân đoạn, bảo
          đảm bản đồ kích hoạt và đặc trưng phân đoạn cùng không gian biểu diễn.
    \item Chuyển từ \GradCAM{} một tầng sang \GradCAMpp{} đa tỉ lệ để xử lý đặc
          thù nhiều đốm bệnh trên một lá.
    \item Thay ngưỡng cố định bằng ngưỡng phân vị theo lớp, đảo ngược quan hệ giữa
          tham số và kết quả để kiểm soát trực tiếp độ phủ.
    \item Ép mặt nạ lớp lá khoẻ về không, đưa tri thức tiên nghiệm vào quy trình
          để loại \num{683} mặt nạ dương tính giả.
    \item Giữ nguyên mọi yếu tố giữa V2 và V3 ngoại trừ bộ nhãn, tạo ra phép so
          sánh có kiểm soát cho tác động của \SAM{}.
\end{itemize}

Hai điểm đã được nêu trước và sẽ được kiểm chứng bằng số liệu ở chương sau: đánh
đổi giữa kiểm soát độ phủ và khả năng mã hoá mức độ bệnh
(\cref{subsec:tradeoff_v2}), và khiếm khuyết của cơ chế bảo vệ dựa trên IoU
(\cref{eq:iou_guard}).

% !TeX root = ../main.tex
\chapter{THỰC NGHIỆM VÀ KẾT QUẢ}
\label{chap:thucnghiem}

\section{Thiết lập thực nghiệm}
\label{sec:thietlap}

\subsection{Môi trường phần cứng và phần mềm}

\begin{table}[H]
\centering
\caption{Môi trường thực nghiệm}
\label{tab:moitruong}
\begin{tabular}{L{4.4cm} L{8.6cm}}
\toprule
\textbf{Thành phần} & \textbf{Cấu hình} \\
\midrule
Bộ xử lý đồ hoạ & NVIDIA Quadro M1000M, 2,15 GB VRAM \\
Hệ điều hành    & Windows 10 Pro \\
Python          & 3.13.12 \\
PyTorch         & 2.6.0+cu124~\cite{paszke2019pytorch} \\
torchvision     & 0.21.0+cu124 \\
CUDA            & 12.4 \\
Thư viện phân đoạn & segmentation-models-pytorch~\cite{iakubovskii2019smp} \\
Thư viện tăng cường & Albumentations~\cite{buslaev2020albumentations} \\
Thư viện xử lý ảnh & OpenCV~\cite{bradski2000opencv} \\
\bottomrule
\end{tabular}
\end{table}

Dung lượng bộ nhớ đồ hoạ chỉ 2,15 GB là ràng buộc chi phối nhiều quyết định thực
nghiệm và cần được nêu rõ để người đọc hiểu đúng bối cảnh. Ràng buộc này giải
thích trực tiếp việc chọn kiến trúc quy mô vừa (EfficientNet-B0 thay vì các biến
thể lớn hơn), kích thước lô nhỏ, và --- như phân tích ở \cref{subsec:confound}
--- cả một số khác biệt cấu hình giữa các phiên bản vốn không nằm trong ý đồ
thiết kế ban đầu.

\subsection{Chỉ số đánh giá và quy ước báo cáo}

Bài toán phân loại được đánh giá bằng \en{accuracy}, \en{precision}, \en{recall}
và $F_1$ theo \cref{eq:prf}, báo cáo cả ở mức tổng hợp lẫn mức từng lớp. Bài toán
phân đoạn dùng IoU và Dice theo \cref{eq:iou,eq:dicecoef}, kèm \en{precision} và
\en{recall} tính trên tập điểm ảnh.

Mọi chỉ số phân đoạn được tính \emph{theo từng ảnh rồi lấy trung bình} trên toàn
tập, chứ không gộp toàn bộ điểm ảnh của cả tập rồi tính một lần. Cách thứ nhất
cho mỗi ảnh trọng số như nhau, phù hợp khi diện tích tổn thương biến thiên mạnh
giữa các ảnh; cách thứ hai sẽ khiến các ảnh có vùng bệnh lớn chi phối kết quả.

Toàn bộ chỉ số báo cáo trong chương này được lấy trực tiếp từ ô kết quả của các
\en{notebook} thực nghiệm hoặc từ tệp kết quả đã lưu, và có thể truy vết về đúng
nguồn sinh ra chúng.

\section{Kết quả phân loại}
\label{sec:ketqua_cls}

\subsection{Diễn biến huấn luyện}

Mô hình được huấn luyện tối đa 50 epoch với cơ chế dừng sớm. Thực tế quá trình
dừng ở epoch 16 khi \texttt{val\_acc} không cải thiện trong 5 epoch liên tiếp.

\begin{figure}[H]
\centering
\includegraphics[width=\textwidth]{figures/cls_training_curves.png}
\caption[Đường cong huấn luyện mô hình phân loại]{Diễn biến hàm mất mát (trái) và
độ chính xác (phải) qua 16 epoch. Đường nét đứt đánh dấu độ chính xác kiểm định
tốt nhất 98,19\% đạt tại epoch 11.}
\label{fig:cls_curves}
\end{figure}

\begin{table}[H]
\centering
\caption{Diễn biến huấn luyện mô hình phân loại qua 16 epoch}
\label{tab:cls_history}
\small
\begin{tabular}{c r r r r l}
\toprule
\textbf{Epoch} & \textbf{train\_loss} & \textbf{train\_acc} & \textbf{val\_loss} & \textbf{val\_acc} & \textbf{Ghi chú} \\
\midrule
1  & 1,0963 & 0,5892 & 0,4712 & 0,8623 & \\
2  & 0,4369 & 0,8547 & 0,2316 & 0,9187 & \\
3  & 0,2740 & 0,9056 & 0,1820 & 0,9368 & \\
4  & 0,2293 & 0,9249 & 0,1222 & 0,9616 & \\
5  & 0,1906 & 0,9327 & 0,1175 & 0,9639 & \\
6  & 0,1602 & 0,9478 & 0,1017 & 0,9684 & \\
7  & 0,1487 & 0,9507 & 0,1356 & 0,9503 & lần đầu không cải thiện \\
8  & 0,1220 & 0,9597 & 0,0926 & 0,9729 & \\
9  & 0,1130 & 0,9613 & 0,0887 & 0,9729 & \\
10 & 0,0989 & 0,9662 & 0,0890 & 0,9752 & \\
\textbf{11} & 0,0977 & 0,9671 & 0,0811 & \textbf{0,9819} & \textbf{điểm kiểm tra tốt nhất} \\
12 & 0,0820 & 0,9729 & 0,0910 & 0,9774 & \\
13 & 0,0791 & 0,9771 & 0,1023 & 0,9684 & \\
14 & 0,0747 & 0,9758 & 0,1108 & 0,9729 & \\
15 & 0,0559 & 0,9800 & 0,1050 & 0,9616 & \\
16 & 0,0476 & 0,9858 & 0,0831 & 0,9684 & dừng sớm kích hoạt \\
\bottomrule
\end{tabular}
\end{table}

Diễn biến chia thành ba pha rõ rệt. \emph{Pha hội tụ nhanh} (epoch 1--6): hàm mất
mát kiểm định giảm từ 0,4712 xuống 0,1017, độ chính xác tăng vọt từ 86,23\% lên
96,84\% --- hiệu quả điển hình của học chuyển giao, mô hình chỉ cần thích ứng đặc
trưng sẵn có với miền dữ liệu mới. \emph{Pha tinh chỉnh} (epoch 7--11): cải thiện
chậm dần, đạt đỉnh 98,19\% tại epoch 11. \emph{Pha quá khớp} (epoch 12--16): dấu
hiệu rất rõ --- \texttt{train\_acc} tiếp tục tăng tới 98,58\% trong khi
\texttt{val\_acc} dao động rồi giảm về 96,84\%, đồng thời \texttt{val\_loss} tăng
từ 0,0811 lên 0,1108 tại epoch 14. Dừng sớm đã can thiệp đúng thời điểm.

Đáng chú ý là ở epoch 16, \texttt{train\_acc} (98,58\%) đã vượt \texttt{val\_acc}
(96,84\%) --- khoảng cách 1,74 điểm phần trăm. Khoảng cách này còn nhỏ, cho thấy
mức quá khớp ở giai đoạn cuối là nhẹ và không phá huỷ khả năng khái quát.

\subsection{Hiệu năng trên tập kiểm thử}

\begin{table}[H]
\centering
\caption{Hiệu năng phân loại trên tập kiểm thử (890 ảnh)}
\label{tab:cls_test}
\begin{tabular}{l c c}
\toprule
\textbf{Chỉ số} & \textbf{Kiểm định (tốt nhất)} & \textbf{Kiểm thử} \\
\midrule
Accuracy            & \textbf{98,19\%} & \textbf{97,19\%} \\
Precision (weighted) & ---             & 97,22\% \\
Recall (weighted)    & ---             & 97,19\% \\
$F_1$ (weighted)     & ---             & 97,19\% \\
\bottomrule
\end{tabular}
\end{table}

Mô hình phân loại đúng 865 trên 890 ảnh kiểm thử, tức chỉ sai 25 ảnh.

Khoảng cách giữa kiểm định (98,19\%) và kiểm thử (97,19\%) là 1,00 điểm phần trăm.
Chênh lệch nhỏ và \emph{theo đúng chiều dự kiến} --- hiệu năng trên tập kiểm thử
thấp hơn --- là dấu hiệu lành mạnh: tập kiểm định đã bị dùng gián tiếp để chọn
điểm kiểm tra và điều chỉnh tốc độ học nên không còn khách quan, còn tập kiểm thử
chỉ được dùng đúng một lần. Nếu khoảng cách này lớn (chẳng hạn trên 5 điểm phần
trăm), đó sẽ là dấu hiệu mô hình đã bị điều chỉnh quá mức theo tập kiểm định.

\subsection{Ma trận nhầm lẫn}

\begin{figure}[H]
\centering
\includegraphics[width=0.8\textwidth]{figures/cls_confusion_matrix.png}
\caption[Ma trận nhầm lẫn trên tập kiểm thử]{Ma trận nhầm lẫn của mô hình
EfficientNet-B0 trên 890 ảnh kiểm thử. Phần lớn nhầm lẫn xảy ra giữa Algal Leaf
Spot và Phomopsis Leaf Spot.}
\label{fig:confusion}
\end{figure}

\begin{table}[H]
\centering
\caption{Ma trận nhầm lẫn trên tập kiểm thử (hàng: nhãn thật, cột: dự đoán)}
\label{tab:confusion}
\small
\begin{tabular}{l r r r r r | r}
\toprule
 & \textbf{Algal} & \textbf{Alloc.} & \textbf{Healthy} & \textbf{Blight} & \textbf{Phom.} & \textbf{Tổng} \\
\midrule
Algal Leaf Spot     & \textbf{142} & 1 & 1 & 0 & 3 & 147 \\
Allocaridara Attack & 2 & \textbf{180} & 0 & 0 & 1 & 183 \\
Healthy Leaf        & 1 & 0 & \textbf{195} & 0 & 0 & 196 \\
Leaf Blight         & 1 & 3 & 1 & \textbf{182} & 1 & 188 \\
Phomopsis Leaf Spot & 4 & 4 & 2 & 0 & \textbf{166} & 176 \\
\bottomrule
\end{tabular}
\end{table}

Ma trận nhầm lẫn xác nhận dự đoán đưa ra từ phân tích khám phá ở
\cref{sec:dulieu}: cặp \emph{Algal Leaf Spot} và \emph{Phomopsis Leaf Spot} nhầm
lẫn qua lại nhiều nhất --- 2 ảnh Algal bị nhận thành Phomopsis và 3 ảnh Phomopsis
bị nhận thành Algal, tổng cộng 5 trên 24 lỗi. Điều này phù hợp với quan sát hình
thái ở \cref{tab:hinhthai}: cả hai bệnh đều tạo đốm nhỏ trên nền lá xanh, khác
biệt chủ yếu ở kích thước và mật độ đốm.

Hai quan sát khác đáng lưu ý. \emph{Leaf Blight có precision tuyệt đối 100\%}:
cột Leaf Blight không chứa bất kỳ giá trị khác không nào nằm ngoài đường chéo ---
không một ảnh thuộc lớp khác từng bị nhận nhầm thành Leaf Blight. Mảng hoại tử lớn
màu nâu đen là dấu hiệu đặc thù mà không lớp nào khác có. Tuy vậy chính lớp này
lại có \emph{recall thấp nhất} (96,28\%): 7 ảnh Leaf Blight bị phân sang các lớp
khác, phân tán khá đều (2 sang Algal, 1 sang Allocaridara, 3 sang Healthy, 1 sang
Phomopsis). Cách phân bố lỗi này gợi ý rằng các ảnh Leaf Blight ở giai đoạn bệnh
sớm --- khi mảng hoại tử chưa lan rộng --- có biểu hiện gần với bệnh đốm hoặc thậm
chí gần lá khoẻ.

Đáng chú ý về mặt ứng dụng là hai chiều nhầm lẫn liên quan tới lớp Healthy Leaf.
Theo chiều \emph{báo động giả}, 5 ảnh lá khoẻ bị nhận thành lớp bệnh (1 sang Algal,
4 sang Phomopsis). Theo chiều \emph{bỏ sót}, 6 ảnh bệnh bị nhận thành lá khoẻ (2
từ Algal, 3 từ Leaf Blight, 1 từ Phomopsis).

Theo tiêu chí ưu tiên \en{recall} đã nêu ở \cref{sec:metrics}, chiều thứ nhất là
loại lỗi ít nghiêm trọng --- cảnh báo nhầm một cây khoẻ chỉ tốn thêm một lần kiểm
tra. Chiều thứ hai mới đáng lo, vì bỏ sót cây bệnh đồng nghĩa mất cơ hội can thiệp
sớm. Với 6 ảnh trên tổng số 694 ảnh bệnh trong tập kiểm thử, tỉ lệ bỏ sót bệnh là
0,86\% --- mức chấp nhận được cho ứng dụng thực tế.

\subsection{Hiệu năng theo từng lớp}

\begin{figure}[H]
\centering
\includegraphics[width=\textwidth]{figures/cls_per_class_metrics.png}
\caption[Chỉ số theo từng lớp]{Precision, recall và $F_1$ cho từng lớp trên tập
kiểm thử. Cả năm lớp đều đạt $F_1$ trên 95\%.}
\label{fig:per_class}
\end{figure}

\begin{table}[H]
\centering
\caption{Hiệu năng phân loại theo từng lớp trên tập kiểm thử}
\label{tab:cls_perclass}
\begin{tabular}{l c c c r}
\toprule
\textbf{Lớp} & \textbf{Precision} & \textbf{Recall} & \textbf{$F_1$} & \textbf{Số mẫu} \\
\midrule
Algal Leaf Spot     & 94,67\%  & 96,60\% & 95,62\% & 147 \\
Allocaridara Attack & 95,74\%  & 98,36\% & \textbf{97,04\%} & 183 \\
Healthy Leaf        & \textbf{97,99\%}  & \textbf{99,49\%} & \textbf{98,73\%} & 196 \\
Leaf Blight         & \textbf{100,00\%} & 96,81\% & 98,38\% & 188 \\
Phomopsis Leaf Spot & 97,08\%  & 94,32\% & 95,68\% & 176 \\
\midrule
\textbf{Trung bình có trọng số} & \textbf{97,22\%} & \textbf{97,19\%} & \textbf{97,19\%} & \textbf{890} \\
\bottomrule
\end{tabular}
\end{table}

Ba nhận xét từ \cref{tab:cls_perclass}:

\textbf{Không lớp nào bị bỏ rơi.} $F_1$ thấp nhất là 95,97\% (Algal Leaf Spot),
cao nhất 98,63\% (Allocaridara Attack) --- biên độ chỉ 2,66 điểm phần trăm. Kết
hợp với tỉ lệ mất cân bằng thấp (1,33), điều này cho thấy mô hình không thiên vị
lớp đa số.

\textbf{Recall vượt precision ở bốn trên năm lớp.} Ngoại lệ duy nhất là
Allocaridara Attack và Leaf Blight. Xu hướng chung này thuận lợi cho bài toán chẩn
đoán bệnh, theo tiêu chí đã nêu ở \cref{sec:metrics}.

\textbf{Algal Leaf Spot là lớp yếu nhất} với precision 94,70\% --- thấp nhất trong
năm lớp. Điều này nhất quán với ma trận nhầm lẫn và với phân tích màu ở
\cref{tab:meanrgb}, nơi lớp này có hiệu $G - R$ thấp nhất và phân bố kênh màu
chồng lấn nhiều nhất với các lớp còn lại.

\section{Phân tích bản đồ kích hoạt}
\label{sec:ketqua_xai}

Bản đồ kích hoạt được sinh cho toàn bộ \num{3104} ảnh huấn luyện, dùng điểm kiểm
tra tốt nhất từ \cref{sec:ketqua_cls}.

\begin{figure}[H]
\centering
\includegraphics[height=0.80\textheight, keepaspectratio]{figures/xai_gradcam_best.png}
\caption[Bản đồ kích hoạt Grad-CAM theo lớp]{Bản đồ kích hoạt của ảnh có độ tin
cậy cao nhất trong mỗi lớp. Từ trái sang phải: ảnh gốc, bản đồ nhiệt, và ảnh chồng
lớp.}
\label{fig:gradcam}
\end{figure}

\begin{table}[H]
\centering
\caption{Độ chính xác và độ tin cậy trên tập huấn luyện theo lớp}
\label{tab:gradcam_stats}
\begin{tabular}{l r c c c}
\toprule
\textbf{Lớp} & \textbf{Số ảnh} & \textbf{Accuracy} & \textbf{Tin cậy TB} & \textbf{Tin cậy nhỏ nhất} \\
\midrule
Algal Leaf Spot     & 513 & 99,03\% & 0,971 & \textbf{0,165} \\
Allocaridara Attack & 639 & 99,22\% & 0,981 & 0,220 \\
Healthy Leaf        & 683 & 99,56\% & 0,979 & 0,204 \\
Leaf Blight         & 655 & \textbf{99,85\%} & 0,988 & 0,354 \\
Phomopsis Leaf Spot & 614 & 99,84\% & \textbf{0,991} & 0,415 \\
\bottomrule
\end{tabular}
\end{table}

Độ chính xác trên tập huấn luyện đạt 99,03--99,85\%, cao hơn mức 97,19\% trên tập
kiểm thử. Chênh lệch này là bình thường và có hai nguyên nhân cộng hưởng: đây là
dữ liệu mô hình đã thấy khi huấn luyện, và khi suy luận thì không áp dụng tăng
cường dữ liệu nên bài toán dễ hơn.

Về mặt chất lượng bản đồ kích hoạt, quan sát \cref{fig:gradcam} cho thấy bốn kiểu
hành vi khác nhau:

\begin{itemize}
    \item \textbf{Algal và Phomopsis}: kích hoạt tập trung vào các đốm bệnh, chứng
          tỏ mô hình phân biệt hai lớp dựa trên đặc trưng đốm chứ không dựa vào
          màu tổng thể của lá.
    \item \textbf{Allocaridara}: kích hoạt bám vào vết cắn và lỗ thủng --- đây là
          lớp có bản đồ kích hoạt tập trung nhất, phù hợp với việc lớp này đạt
          $F_1$ cao nhất.
    \item \textbf{Leaf Blight}: kích hoạt trải rộng theo mảng hoại tử, không có
          tâm rõ ràng.
    \item \textbf{Healthy Leaf}: kích hoạt \emph{khuếch tán khắp phiến lá}, không
          có vùng trọng tâm.
\end{itemize}

Hành vi cuối cùng có ý nghĩa then chốt và cần được nhấn mạnh. Với lá khoẻ, mô hình
không có tổn thương nào để tập trung vào; nó nhận diện lớp này qua \emph{sự vắng
mặt} của dấu hiệu bệnh và qua màu xanh đồng đều trên toàn lá. Bản đồ kích hoạt vì
thế trải đều --- một hành vi hoàn toàn hợp lý xét theo nhiệm vụ phân loại, nhưng
lại vô nghĩa và gây hại nếu đem nhị phân hoá thành mặt nạ tổn thương. Đây chính là
bằng chứng thực nghiệm cho luận điểm lý thuyết ở \cref{sec:weaksup}: bản đồ kích
hoạt trả lời câu hỏi ``mô hình nhìn vào đâu'', không phải ``bệnh nằm ở đâu''.

Về độ tin cậy, giá trị trung bình đều trên 0,97 nhưng giá trị nhỏ nhất chênh lệch
đáng kể giữa các lớp: Algal Leaf Spot xuống tới 0,165 trong khi Phomopsis Leaf Spot
chỉ xuống 0,415. Điều này cho thấy Algal có một số trường hợp biên mà mô hình thực
sự không chắc chắn; bản đồ kích hoạt của chúng kém tin cậy hơn khi dùng làm nhãn
giả.

\section{Chất lượng nhãn giả}
\label{sec:ketqua_pseudo}

\subsection{Nhãn giả V1}

\begin{figure}[H]
\centering
\includegraphics[height=0.80\textheight, keepaspectratio]{figures/pl1_per_class.png}
\caption[Nhãn giả V1 theo lớp]{Mặt nạ giả V1 cho ảnh có độ tin cậy cao nhất mỗi
lớp. Từ trái sang phải: ảnh gốc, bản đồ nhiệt chồng lớp, mặt nạ giả chồng lớp.}
\label{fig:pl1}
\end{figure}

\begin{table}[H]
\centering
\caption{Thống kê độ phủ mặt nạ giả V1 theo lớp}
\label{tab:cov_v1}
\begin{tabular}{l c c c c}
\toprule
\textbf{Lớp} & \textbf{Trung bình} & \textbf{Độ lệch chuẩn} & \textbf{Nhỏ nhất} & \textbf{Lớn nhất} \\
\midrule
Algal Leaf Spot     & 15,1\% & 0,056 & 5,7\%  & 42,6\% \\
Allocaridara Attack & 13,7\% & 0,039 & 5,0\%  & 26,5\% \\
Healthy Leaf        & \textbf{26,4\%} & 0,076 & 11,1\% & \textbf{59,1\%} \\
Leaf Blight         & 12,3\% & 0,047 & 5,5\%  & 43,6\% \\
Phomopsis Leaf Spot & 14,5\% & 0,048 & 5,4\%  & 37,5\% \\
\bottomrule
\end{tabular}
\end{table}

Số liệu ở \cref{tab:cov_v1} định lượng chính xác vấn đề đã dự đoán ở
\cref{subsec:hanche_v1}. Lớp Healthy Leaf --- lớp \emph{không có} tổn thương ---
lại có độ phủ trung bình cao nhất (26,4\%), gần gấp đôi lớp Leaf Blight (12,3\%)
vốn là bệnh lan rộng nhất. Ảnh cực đoan nhất có tới 59,1\% diện tích bị đánh dấu
là ``bệnh'' trên một chiếc lá hoàn toàn khoẻ mạnh.

Độ lệch chuẩn cũng lớn (0,039--0,076), phản ánh việc ngưỡng cố định 0,5 không
thích ứng được với phân bố riêng của từng bản đồ nhiệt.

Toàn bộ \num{3104} mặt nạ được sinh trong 14 phút 00 giây. Độ chính xác dự đoán
trên tập này đạt 99,52\% với độ tin cậy trung bình 0,9822 --- nghĩa là chỉ 15 ảnh
có nhãn giả sinh ra từ dự đoán sai.

\subsection{Nhãn giả V2}

\begin{figure}[H]
\centering
\includegraphics[height=0.80\textheight, keepaspectratio]{figures/pl2_per_class.png}
\caption[Nhãn giả V2 theo lớp]{Mặt nạ giả V2 sinh bằng GradCAM++ đa tỉ lệ kết hợp
ngưỡng phân vị theo lớp. Mặt nạ của lớp Healthy Leaf là ma trận không.}
\label{fig:pl2}
\end{figure}

\begin{table}[H]
\centering
\caption{Độ phủ mục tiêu và độ phủ thực tế của nhãn giả V2}
\label{tab:cov_v2}
\begin{tabular}{l c c c c}
\toprule
\textbf{Lớp} & \textbf{Mục tiêu} & \textbf{Thực tế} & \textbf{Sai lệch} & \textbf{Độ lệch chuẩn} \\
\midrule
Algal Leaf Spot     & 18\% & 17,7\% & $-0{,}3$ & 0,008 \\
Allocaridara Attack & 20\% & 19,7\% & $-0{,}3$ & 0,007 \\
Healthy Leaf        & 0\%  & 0,0\%  & $0{,}0$  & 0,000 \\
Leaf Blight         & 22\% & 21,7\% & $-0{,}3$ & 0,007 \\
Phomopsis Leaf Spot & 18\% & 17,6\% & $-0{,}4$ & 0,009 \\
\bottomrule
\end{tabular}
\end{table}

Cơ chế ngưỡng phân vị hoạt động đúng như thiết kế: sai lệch so với mục tiêu không
vượt quá 0,4 điểm phần trăm ở bất kỳ lớp nào.

\begin{figure}[H]
\centering
\includegraphics[width=\textwidth]{figures/pl2_coverage.png}
\caption[Phân bố độ phủ V1 và V2]{Phân bố độ phủ của toàn bộ mặt nạ (trái) và độ
phủ trung bình theo lớp của V2 (phải). Cụm tại giá trị 0 trong biểu đồ bên trái
tương ứng 683 mặt nạ rỗng của lớp Healthy Leaf.}
\label{fig:pl2_coverage}
\end{figure}

\begin{table}[H]
\centering
\caption{So sánh tổng hợp nhãn giả V1 và V2}
\label{tab:v1v2_stats}
\begin{tabular}{l c c}
\toprule
\textbf{Chỉ số} & \textbf{V1} & \textbf{V2} \\
\midrule
Tổng số mặt nạ            & 3.104  & 3.104 \\
Độ chính xác dự đoán      & 99,52\% & 99,52\% \\
Độ tin cậy trung bình     & 0,9822 & 0,9822 \\
Độ phủ trung bình         & 0,1660 & 0,1504 \\
Mặt nạ gần rỗng ($<0{,}05$) & 1      & \textbf{683} \\
Mặt nạ quá lớn ($>0{,}60$)  & 0      & 0 \\
Thời gian sinh nhãn       & 14 ph 00 s & \textbf{4 ph 08 s} \\
\bottomrule
\end{tabular}
\end{table}

Hai chỉ số đầu giống hệt nhau giữa hai phiên bản, điều này là tất yếu vì cả hai
dùng chung một mô hình phân loại và cùng bộ ảnh --- chỉ khác cách chuyển bản đồ
kích hoạt thành mặt nạ.

Con số \textbf{683 mặt nạ gần rỗng} ở V2 cần được đọc đúng: đây \emph{không} phải
lỗi mà là kết quả có chủ đích của việc ép về không cho lớp Healthy Leaf theo
\cref{eq:hardzero}. Giá trị 683 trùng khớp chính xác với số ảnh lá khoẻ trong tập
huấn luyện, xác nhận cơ chế hoạt động đúng và không ảnh hưởng tới lớp nào khác.

Việc V2 nhanh hơn V1 tới 3,4 lần (4 phút 08 giây so với 14 phút 00 giây) dù phải
tính \emph{hai} bản đồ kích hoạt thay vì một là kết quả của cách cài đặt hiệu quả
hơn cho bước truyền ngược, không phải do giảm khối lượng tính toán.

\subsection{Đánh đổi đã được xác nhận}

Số liệu ở \cref{tab:cov_v2} xác nhận đánh đổi đã dự báo ở
\cref{subsec:tradeoff_v2}: độ lệch chuẩn của độ phủ giảm mạnh từ khoảng
0,039--0,076 ở V1 xuống chỉ còn 0,007--0,009 ở V2. Nói cách khác, \emph{mọi ảnh
trong cùng một lớp đều nhận mặt nạ có diện tích gần như bằng nhau}.

Đây vừa là thành công vừa là hạn chế. Thành công vì độ phủ được kiểm soát chặt chẽ
đúng như mục tiêu thiết kế. Hạn chế vì mặt nạ không còn mã hoá mức độ nặng nhẹ của
bệnh: một chiếc lá chỉ có một đốm nhỏ và một chiếc lá bị bệnh gần kín đều được gán
mặt nạ chiếm xấp xỉ 18\% hoặc 22\% diện tích. Mô hình phân đoạn huấn luyện trên
nhãn này học cách ``phủ đúng tỉ lệ diện tích'' thay vì học cách ``phát hiện đúng
vùng bệnh''.

\section{Phân tích tinh chỉnh bằng SAM}
\label{sec:ketqua_sam}

\subsection{Hiện tượng nở mặt nạ}

\begin{figure}[H]
\centering
\includegraphics[width=0.72\textwidth]{figures/sam_refinement.png}
\caption[So sánh nhãn V2 và nhãn SAM tinh chỉnh]{Bốn mẫu so sánh ảnh gốc, mặt nạ
giả V2 và mặt nạ sau khi SAM tinh chỉnh. SAM cho biên sắc nét hơn rõ rệt nhưng
đồng thời mở rộng diện tích mặt nạ đáng kể.}
\label{fig:sam_refine}
\end{figure}

Đo lường trực tiếp trên toàn bộ \num{3104} mặt nạ sau tinh chỉnh cho kết quả ở
\cref{tab:sam_expansion}.

\begin{table}[H]
\centering
\caption{Độ phủ trước và sau khi tinh chỉnh bằng SAM}
\label{tab:sam_expansion}
\begin{tabular}{l r r r r}
\toprule
\textbf{Lớp} & \textbf{Độ phủ V2} & \textbf{Độ phủ SAM} & \textbf{Độ lệch chuẩn} & \textbf{Mức tăng} \\
\midrule
Algal Leaf Spot     & 17,7\% & 34,63\% & 0,1210 & \textbf{+95,7\%} \\
Allocaridara Attack & 19,7\% & 31,98\% & 0,1112 & +62,3\% \\
Healthy Leaf        & 0,0\%  & 0,00\%  & 0,0000 & --- \\
Leaf Blight         & 21,7\% & 30,71\% & 0,1176 & +41,5\% \\
Phomopsis Leaf Spot & 17,6\% & 33,56\% & 0,1207 & +90,7\% \\
\midrule
\textbf{Toàn bộ}    & \textbf{15,04\%} & \textbf{25,43\%} & 0,1708 & \textbf{+69,1\%} \\
\bottomrule
\end{tabular}
\end{table}

Kết quả trái ngược với kỳ vọng ban đầu. \SAM{} được đưa vào quy trình với vai trò
\emph{tinh chỉnh biên} --- tức giữ nguyên vị trí và diện tích, chỉ làm biên chính
xác hơn. Thực tế nó \emph{mở rộng} diện tích mặt nạ từ 41,5\% đến 95,7\% tuỳ lớp,
trung bình toàn bộ tăng 69,1\%.

Hai chỉ số bổ sung làm rõ bản chất hiện tượng:

\begin{itemize}
    \item \textbf{200 trên \num{3104} mặt nạ (6,4\%) có độ phủ vượt 50\%} --- tức
          hơn một nửa diện tích ảnh bị đánh dấu là tổn thương, điều không thể xảy
          ra với bệnh đốm lá.
    \item \textbf{Độ lệch chuẩn tăng vọt} từ 0,007--0,009 (V2) lên 0,111--0,121
          (SAM), tức tăng khoảng 15 lần. Nhãn SAM biến thiên mạnh giữa các ảnh
          thay vì đồng nhất như V2.
    \item \textbf{683 mặt nạ rỗng được giữ nguyên rỗng} --- cơ chế ép về không cho
          lớp lá khoẻ không bị \SAM{} phá vỡ, đúng như thiết kế.
\end{itemize}

Nguyên nhân đã được dự báo về mặt lý thuyết ở \cref{sec:sam}: \SAM{} được huấn
luyện để phân đoạn \emph{đơn vị ngữ nghĩa tự nhiên} của cảnh. Khi nhận gợi ý là
vài điểm nằm trên một đốm bệnh, đơn vị ngữ nghĩa mà \SAM{} nhận diện không phải
đốm bệnh mà là \emph{cả chiếc lá} chứa đốm đó --- quan sát trực tiếp trên
\cref{fig:sam_refine} xác nhận điều này.

\subsection{Vì sao cơ chế bảo vệ IoU không phát hiện được}
\label{subsec:iouguard}

Cơ chế bảo vệ theo \cref{eq:iou_guard} lẽ ra phải phát hiện các mặt nạ sai lệch.
Nó đã không phát hiện được trường hợp nào của hiện tượng nở. Phân tích dưới đây
cho thấy điều này \emph{không phải sự cố ngẫu nhiên mà là hệ quả tất yếu} của công
thức được chọn.

\begin{quote}
\textbf{Mệnh đề.} Giả sử \SAM{} sinh ra mặt nạ $M_{\SAM}$ \emph{chứa trọn} mặt nạ
gốc $M_{\text{V2}}$, tức $M_{\text{V2}} \subseteq M_{\SAM}$ --- đúng với trường hợp
nở ra bao cả chiếc lá. Khi đó cơ chế bảo vệ với ngưỡng $\theta$ chỉ kích hoạt khi
độ phủ của \SAM{} vượt quá $1/\theta$ lần độ phủ của V2.
\end{quote}

\textbf{Chứng minh.} Khi $M_{\text{V2}} \subseteq M_{\SAM}$ thì
$|M_{\SAM} \cap M_{\text{V2}}| = |M_{\text{V2}}|$ và
$|M_{\SAM} \cup M_{\text{V2}}| = |M_{\SAM}|$. Do đó

\begin{equation}
    \operatorname{IoU}(M_{\SAM}, M_{\text{V2}})
    = \frac{|M_{\text{V2}}|}{|M_{\SAM}|}
    = \frac{c_{\text{V2}}}{c_{\SAM}}
    \label{eq:iou_containment}
\end{equation}

với $c$ ký hiệu độ phủ. Cơ chế bảo vệ kích hoạt khi
$c_{\text{V2}}/c_{\SAM} < \theta$, tương đương $c_{\SAM} > c_{\text{V2}}/\theta$.
$\square$

Áp dụng vào số liệu thực tế với $\theta = 0{,}15$ cho kết quả ở
\cref{tab:iouguard_analysis}.

\begin{table}[H]
\centering
\caption{Ngưỡng độ phủ cần thiết để cơ chế bảo vệ IoU kích hoạt}
\label{tab:iouguard_analysis}
\begin{tabular}{l c c c}
\toprule
\textbf{Lớp} & \textbf{Độ phủ V2} & \textbf{Độ phủ SAM cần vượt} & \textbf{Khả thi?} \\
\midrule
Algal Leaf Spot     & 17,7\% & 118,0\% & \textbf{Không} \\
Allocaridara Attack & 19,7\% & 131,3\% & \textbf{Không} \\
Leaf Blight         & 21,7\% & 144,7\% & \textbf{Không} \\
Phomopsis Leaf Spot & 17,6\% & 117,3\% & \textbf{Không} \\
\bottomrule
\end{tabular}
\end{table}

Kết luận rất rõ ràng: với mọi lớp bệnh, ngưỡng độ phủ cần thiết để kích hoạt cơ chế
bảo vệ đều \emph{vượt quá 100\% diện tích ảnh} --- một điều bất khả thi về mặt vật
lý. Nói cách khác:

\begin{quote}
Với cấu hình đã dùng, cơ chế bảo vệ dựa trên IoU \textbf{không bao giờ có thể}
phát hiện hiện tượng nở mặt nạ ở bất kỳ lớp bệnh nào. Nó chỉ chặn được trường hợp
\SAM{} \emph{trôi} sang một vùng hoàn toàn khác, chứ hoàn toàn bất lực trước trường
hợp \SAM{} \emph{nở} ra bao trùm mặt nạ gốc.
\end{quote}

Đây là một khiếm khuyết thiết kế mang tính hệ thống, không phải lỗi cài đặt. Nguồn
gốc của nó là IoU đối xứng --- nó phạt cả hai chiều sai lệch như nhau và do đó
không phân biệt được ``lệch vị trí'' với ``đúng vị trí nhưng quá rộng''.

\subsection{Giải pháp đề xuất}

Khiếm khuyết trên có thể khắc phục bằng cách bổ sung một ràng buộc \emph{bất đối
xứng} về tỉ lệ độ phủ, hoạt động song song với cơ chế IoU hiện có:

\begin{equation}
    M_{\text{refined}} =
    \begin{cases}
        M_{\text{V2}} & \text{nếu } \operatorname{IoU} < \theta \quad \text{(chống trôi)} \\[4pt]
        M_{\text{V2}} & \text{nếu } c_{\SAM} > \kappa \cdot c_{\text{V2}} \quad \text{(chống nở)} \\[4pt]
        M_{\SAM} & \text{trong các trường hợp còn lại}
    \end{cases}
    \label{eq:iou_guard_fixed}
\end{equation}

Với $\kappa = 2{,}0$, ràng buộc thứ hai sẽ kích hoạt cho phần lớn các trường hợp
nở đã quan sát --- chẳng hạn lớp Algal Leaf Spot có tỉ lệ nở trung bình 1,96 lần,
nằm sát ngưỡng, còn các ảnh cực đoan với độ phủ vượt 50\% đều bị chặn. Việc kiểm
chứng thực nghiệm cho giải pháp này nằm ngoài phạm vi luận văn và được đề xuất là
hướng phát triển ở \cref{chap:ketluan}.

\section{Kết quả phân đoạn}
\label{sec:ketqua_seg}

\subsection{Tổng hợp ba phiên bản}

\begin{table}[H]
\centering
\caption{Hiệu năng phân đoạn trên tập kiểm thử}
\label{tab:seg_results}
\small
\begin{tabular}{l c c c c c}
\toprule
\textbf{Phiên bản} & \textbf{IoU} & \textbf{Dice} & \textbf{Precision} & \textbf{Recall} & \textbf{Loss} \\
\midrule
V1 --- EfficientNet-UNet, nhãn V1$^{\dagger}$ & 0,7671 & 0,8623 & 0,8562 & 0,8877 & 0,3065 \\
V2 --- UNet++, nhãn V2                        & 0,5135 & 0,6301 & 0,6169 & 0,7419 & 0,5993 \\
V3 --- UNet++, nhãn SAM                       & 0,6239 & 0,7053 & 0,7090 & 0,8752 & 0,4704 \\
\bottomrule
\multicolumn{6}{l}{\small $^{\dagger}$ V1 dùng kiến trúc, phân bố nhãn và tập kiểm thử khác V2/V3 --- xem \cref{sec:caveat}.}
\end{tabular}
\end{table}

\begin{table}[H]
\centering
\caption{Diễn biến huấn luyện ba phiên bản phân đoạn}
\label{tab:seg_training}
\begin{tabular}{l C{1.7cm} C{2.5cm} C{2.5cm} C{2.5cm}}
\toprule
\textbf{Phiên bản} & \textbf{Số epoch} & \textbf{val\_loss nhỏ nhất} & \textbf{val\_IoU lớn nhất} & \textbf{Điểm kiểm tra chọn} \\
\midrule
V1 & 30 & 0,2959 (E23) & 0,7715 (E25) & E23 \\
V2 & 25 & 0,5957 (E20) & 0,5438 (E18) & E20 \\
V3 & 22 & 0,4478 (E12) & 0,6981 (E17) & E12 \\
\bottomrule
\end{tabular}
\end{table}

\begin{figure}[H]
\centering
\includegraphics[width=\textwidth]{figures/segv1_curves.png}
\caption[Đường cong huấn luyện phiên bản V1]{Diễn biến hàm mất mát, IoU và Dice
của phiên bản V1 qua 30 epoch.}
\label{fig:segv1_curves}
\end{figure}

\begin{figure}[H]
\centering
\includegraphics[width=\textwidth]{figures/segv2_curves.png}
\caption[Đường cong huấn luyện phiên bản V2]{Diễn biến huấn luyện của phiên bản V2
qua 25 epoch, kèm lịch tốc độ học cô-sin và biểu đồ so sánh với V1.}
\label{fig:segv2_curves}
\end{figure}

\begin{figure}[H]
\centering
\includegraphics[width=\textwidth]{figures/segv3_curves.png}
\caption[Đường cong huấn luyện phiên bản V3]{Diễn biến huấn luyện của phiên bản V3
qua 22 epoch, kèm biểu đồ so sánh ba phiên bản.}
\label{fig:segv3_curves}
\end{figure}

\subsection{Kết quả định tính}

\begin{figure}[H]
\centering
\includegraphics[width=\textwidth]{figures/segv1_predictions.png}
\caption[Dự đoán của phiên bản V1]{Bốn mẫu dự đoán của V1 trên tập kiểm thử: ảnh
đầu vào, nhãn giả tham chiếu, và mặt nạ dự đoán.}
\label{fig:segv1_pred}
\end{figure}

\begin{figure}[H]
\centering
\includegraphics[width=\textwidth]{figures/segv2_predictions.png}
\caption[Dự đoán của phiên bản V2]{Bốn mẫu dự đoán của V2: ảnh đầu vào, nhãn giả
V2, bản đồ xác suất và mặt nạ dự đoán kèm chỉ số IoU/Dice từng ảnh.}
\label{fig:segv2_pred}
\end{figure}

\begin{figure}[H]
\centering
\includegraphics[width=\textwidth]{figures/segv3_predictions.png}
\caption[Dự đoán của phiên bản V3]{Bốn mẫu dự đoán của V3 với nhãn tham chiếu là
mặt nạ đã được SAM tinh chỉnh. Biên dự đoán sắc nét hơn so với V2 nhưng diện tích
lớn hơn đáng kể.}
\label{fig:segv3_pred}
\end{figure}

\section{Phân tích tách bạch}
\label{sec:ablation}

\subsection{Giới hạn về khả năng so sánh giữa các phiên bản}
\label{sec:caveat}

Đọc \cref{tab:seg_results} một cách ngây thơ sẽ dẫn tới kết luận: V1 tốt nhất
(IoU 0,7671), V3 đứng thứ hai (0,6239), V2 kém nhất (0,5135) --- và do đó các
``cải tiến'' ở V2, V3 thực chất làm hỏng mô hình. \emph{Kết luận này sai}, và việc
chỉ ra vì sao nó sai là một trong những đóng góp chính của luận văn.

Vấn đề nằm ở chỗ mỗi phiên bản được đánh giá trên \emph{mốc tham chiếu của chính
nó}. Như đã nêu ở \cref{sec:metrics}, IoU và Dice luôn được định nghĩa tương đối
với một tập nhãn tham chiếu $Y$. Ở đây:

\begin{itemize}
    \item V1 được đo so với nhãn giả V1;
    \item V2 được đo so với nhãn giả V2;
    \item V3 được đo so với nhãn đã qua tinh chỉnh SAM.
\end{itemize}

Ba mốc tham chiếu này khác nhau về cả hình dạng, độ phủ lẫn độ biến thiên. So sánh
trực tiếp các con số IoU giữa chúng chẳng khác nào so sánh điểm số của ba học sinh
làm ba đề thi khác nhau với độ khó khác nhau.

\Cref{tab:confound} liệt kê đầy đủ các biến khác nhau giữa các cặp phiên bản.

\begin{table}[H]
\centering
\caption{Các biến khác biệt giữa các cặp phiên bản}
\label{tab:confound}
\small
\begin{tabular}{L{5.0cm} c c}
\toprule
\textbf{Biến} & \textbf{V1 so với V2} & \textbf{V2 so với V3} \\
\midrule
Phân bố nhãn tham chiếu   & Khác & Khác (chủ đích) \\
Kiến trúc mô hình         & Khác & Giống \\
Tập kiểm thử              & \textbf{Khác} & Giống \\
Hàm mất mát               & Khác & Giống \\
Bộ tối ưu và lịch tốc độ học & Khác & Giống \\
Bộ tăng cường dữ liệu     & Khác & Giống \\
Kích thước lô             & Khác & Khác (ngoài ý muốn) \\
Patience dừng sớm         & Khác & Khác (ngoài ý muốn) \\
\midrule
\textbf{Số biến khác nhau} & \textbf{8} & \textbf{3} \\
\bottomrule
\end{tabular}
\end{table}

Cặp V1--V2 khác nhau ở \emph{tám} biến cùng lúc. Với thiết kế thực nghiệm như vậy,
không thể quy chênh lệch hiệu năng cho bất kỳ nguyên nhân đơn lẻ nào. Đặc biệt,
việc hai phiên bản dùng \emph{hai tập kiểm thử khác nhau} --- do khác bộ sinh số
ngẫu nhiên như đã phân tích ở \cref{subsec:split_seg} --- khiến ngay cả phép so
sánh thô cũng mất cơ sở.

Vậy tại sao V1 lại đạt IoU cao nhất? Câu trả lời nằm ở \emph{độ khó của mục tiêu}
chứ không ở chất lượng mô hình:

\begin{itemize}
    \item Nhãn V1 là các khối tròn lớn, trơn, hình dạng đơn giản và dễ dự đoán ---
          hệ quả trực tiếp của việc phóng bản đồ $7 \times 7$ lên $224 \times 224$.
          Tái tạo lại một khối trơn như vậy là bài toán dễ.
    \item Nhãn V3 là các vùng do \SAM{} sinh ra: biên sắc, hình dạng phức tạp, bám
          theo cấu trúc thực của ảnh. Tái tạo chính xác các biên này khó hơn nhiều.
\end{itemize}

Nói cách khác, \textbf{IoU cao của V1 phản ánh việc mục tiêu của nó dễ, không phản
ánh việc nó định vị bệnh tốt hơn.} Trong điều kiện không có nhãn chuẩn ở mức điểm
ảnh, không thể thiết lập được hiệu năng định vị bệnh tuyệt đối của bất kỳ phiên
bản nào.

\subsection{Phép so sánh hợp lệ: V2 so với V3}

Cặp V2--V3 chỉ khác nhau ở ba biến, trong đó biến chính --- bộ nhãn huấn luyện ---
là biến ta muốn khảo sát. Kiến trúc, hàm mất mát, bộ tối ưu, bộ tăng cường và
\emph{tập kiểm thử} đều giống hệt nhau. Đây là phép so sánh có kiểm soát tốt nhất
trong luận văn.

\begin{table}[H]
\centering
\caption{So sánh V2 và V3 --- cùng kiến trúc, cùng tập kiểm thử}
\label{tab:v2v3}
\begin{tabular}{L{3.1cm} C{1.9cm} C{1.9cm} C{2.3cm} C{2.3cm}}
\toprule
\textbf{Chỉ số} & \textbf{V2} & \textbf{V3} & \textbf{Chênh tuyệt đối} & \textbf{Chênh tương đối} \\
\midrule
IoU       & 0,5135 & 0,6239 & $+0{,}1104$ & \textbf{+21,5\%} \\
Dice      & 0,6301 & 0,7053 & $+0{,}0752$ & +11,9\% \\
Precision & 0,6169 & 0,7090 & $+0{,}0921$ & +14,9\% \\
Recall    & 0,7419 & 0,8752 & $+0{,}1333$ & +18,0\% \\
\midrule
val\_loss nhỏ nhất & 0,5957 (E20) & 0,4478 (E12) & $-0{,}1479$ & $-24{,}8\%$ \\
\bottomrule
\end{tabular}
\end{table}

Cải thiện xuất hiện đồng đều trên cả bốn chỉ số, không có chỉ số nào đi ngược
chiều --- dấu hiệu của một tác động thực chứ không phải dao động ngẫu nhiên.

Bằng chứng bổ trợ mạnh mẽ đến từ tốc độ hội tụ: V3 đạt \texttt{val\_loss} 0,4478 ở
epoch 12, trong khi V2 phải tới epoch 20 mới đạt 0,5957 --- cao hơn đáng kể. Cùng
kiến trúc và cùng siêu tham số, mô hình học được nhanh hơn và sâu hơn trên nhãn
\SAM{}. Điều này cho thấy nhãn \SAM{} cung cấp \emph{tín hiệu huấn luyện sạch hơn}:
biên rõ ràng và nhất quán giúp mô hình tìm được quy luật dễ hơn.

\subsection{Ba yếu tố gây nhiễu cần trừ hao}
\label{subsec:confound}

Việc kết luận toàn bộ mức cải thiện +21,5\% là do \SAM{} sẽ là kết luận vội vàng.
Ba yếu tố sau có thể đóng góp một phần, và mục này đánh giá chiều tác động của
từng yếu tố.

\textbf{Yếu tố thứ nhất --- dừng sớm quá sớm ở V2.} Phiên bản V2 dùng patience = 5
còn V3 dùng patience = 10. Quan trọng hơn, dữ liệu cho thấy V2 \emph{vẫn đang cải
thiện} khi bị dừng: \texttt{val\_IoU} tăng từ 0,4960 tại epoch 20 lên 0,5339 tại
epoch 25 --- tức tăng 7,6\% trong 5 epoch cuối cùng trước khi dừng. Cơ chế dừng sớm
theo dõi \texttt{val\_loss} vốn đã chững lại, nên không nhận ra IoU vẫn đang lên.
Nếu V2 được cấp patience = 10 như V3, nhiều khả năng nó đã đạt hiệu năng cao hơn.
\emph{Chiều tác động: yếu tố này làm phóng đại khoảng cách V2--V3.}

\textbf{Yếu tố thứ hai --- kích thước mục tiêu.} Nhãn V3 có độ phủ trung bình
25,43\% so với 15,04\% của V2, tức lớn gấp 1,69 lần. Đây không phải chi tiết vô
hại, vì IoU có thiên lệch hệ thống theo kích thước đối tượng. Với một dải sai số
biên có bề rộng $w$ cố định, phần diện tích sai lệch tỉ lệ với chu vi nhân $w$,
trong khi mẫu số của IoU tỉ lệ với diện tích. Với các vùng có hình dạng tương tự
nhau, chu vi tỉ lệ với căn bậc hai của diện tích, nên sai số biên tương đối giảm
theo $1/\sqrt{A}$. Diện tích tăng 1,69 lần tương ứng sai số biên tương đối giảm
khoảng $1 - 1/\sqrt{1{,}69} \approx 23\%$. \emph{Chiều tác động: yếu tố này cũng
làm phóng đại lợi thế của V3.}

\textbf{Yếu tố thứ ba --- kích thước lô.} V2 dùng lô 4, V3 dùng lô 8. Với bộ nhớ
đồ hoạ chỉ 2,15 GB, đây là ràng buộc kỹ thuật chứ không phải lựa chọn thiết kế.
Lô lớn hơn cho ước lượng gradient ổn định hơn, có thể có lợi cho V3. \emph{Chiều
tác động: nhiều khả năng cũng có lợi cho V3, nhưng mức độ khó ước lượng.}

Cả ba yếu tố đều tác động \emph{cùng một chiều} là làm phóng đại lợi thế của V3.
Kết luận trung thực vì thế phải là:

\begin{quote}
Tinh chỉnh biên bằng \SAM{} \textbf{cải thiện chất lượng tín hiệu huấn luyện} ---
bằng chứng thuyết phục nhất là mô hình hội tụ nhanh hơn và đạt hàm mất mát kiểm
định thấp hơn hẳn với cùng kiến trúc và siêu tham số. Tuy nhiên, con số +21,5\%
IoU là \textbf{cận trên} của tác động này chứ không phải giá trị thực, vì cả ba
yếu tố gây nhiễu đã xác định đều tác động cùng chiều làm tăng khoảng cách.
\end{quote}

\subsection{Nghịch lý của SAM}

Một điểm thú vị về mặt phương pháp luận: \SAM{} vừa \emph{làm hỏng} nhãn theo tiêu
chí bệnh học --- mở rộng mặt nạ tới 69,1\%, bao cả những vùng lá khoẻ mạnh ---
lại vừa \emph{cải thiện} kết quả huấn luyện trên mọi chỉ số.

Hai sự thật này không mâu thuẫn. Chúng đo hai thứ khác nhau. Nhãn \SAM{} sai về
\emph{phạm vi} (bao quá rộng) nhưng đúng về \emph{cấu trúc} (biên bám theo ranh
giới thực trong ảnh, nhất quán giữa các ảnh). Mô hình học được từ tính nhất quán
cấu trúc đó, kể cả khi phạm vi bị phóng đại.

Hệ quả thực tiễn quan trọng: một mô hình phân đoạn huấn luyện trên nhãn V3 sẽ có
xu hướng dự đoán vùng bệnh rộng hơn thực tế. Với ứng dụng ước lượng \emph{tỉ lệ
diện tích lá bị bệnh} để ra quyết định canh tác, sai lệch hệ thống này phải được
hiệu chỉnh trước khi sử dụng.

\section{Phân tích lỗi và hạn chế}
\label{sec:phantichloi}

\subsection{Recall luôn vượt precision}

\begin{table}[H]
\centering
\caption{Quan hệ giữa recall và precision ở ba phiên bản}
\label{tab:recall_precision}
\begin{tabular}{l c c c}
\toprule
\textbf{Phiên bản} & \textbf{Precision} & \textbf{Recall} & \textbf{Chênh lệch} \\
\midrule
V1 & 0,8562 & 0,8877 & $+0{,}0315$ \\
V2 & 0,6169 & 0,7419 & $+0{,}1250$ \\
V3 & 0,7090 & 0,8752 & $+0{,}1662$ \\
\bottomrule
\end{tabular}
\end{table}

Cả ba phiên bản đều có recall vượt precision, và khoảng cách tăng dần từ V1 đến
V3. Điều này nghĩa là mô hình \emph{phân đoạn thừa} một cách hệ thống: nó đánh dấu
nhiều điểm ảnh là bệnh hơn so với nhãn tham chiếu.

Nguyên nhân là mô hình kế thừa đúng đặc tính của nhãn mà nó được huấn luyện. Nhãn
giả vốn đã có xu hướng bao rộng hơn tổn thương thực --- do bản đồ kích hoạt bị làm
mờ khi phóng từ độ phân giải thấp, và với V3 thì cộng thêm hiện tượng nở của
\SAM{}. Mô hình học đúng phân bố đó.

Xét theo yêu cầu ứng dụng nêu ở \cref{sec:metrics}, xu hướng này \emph{có lợi}:
trong chẩn đoán bệnh cây trồng, bỏ sót vùng bệnh nguy hiểm hơn cảnh báo thừa.

\subsection{Tiêu chí dừng sớm không phù hợp}
\label{subsec:earlystop}

Đây là hạn chế kỹ thuật nghiêm trọng nhất được phát hiện trong quá trình thực
nghiệm.

\begin{table}[H]
\centering
\caption{Diễn biến các epoch then chốt của phiên bản V3}
\label{tab:v3_epochs}
\begin{tabular}{c c c c l}
\toprule
\textbf{Epoch} & \textbf{val\_loss} & \textbf{val\_IoU} & \textbf{val\_Dice} & \textbf{Ghi chú} \\
\midrule
9  & 0,4514 & 0,6536 & 0,7348 & \\
11 & 0,4495 & 0,6711 & 0,7518 & \\
\textbf{12} & \textbf{0,4478} & 0,6479 & 0,7294 & \textbf{được chọn} (val\_loss nhỏ nhất) \\
16 & 0,4508 & 0,6917 & 0,7716 & \\
\textbf{17} & 0,4537 & \textbf{0,6981} & \textbf{0,7792} & \textbf{bị bỏ lỡ} (val\_IoU lớn nhất) \\
\bottomrule
\end{tabular}
\end{table}

Điểm kiểm tra được chọn theo tiêu chí \texttt{val\_loss} nhỏ nhất, rơi vào epoch
12 với val\_IoU = 0,6479. Nhưng epoch 17 có val\_IoU = 0,6981 --- \textbf{cao hơn
0,0502} --- lại bị bỏ qua vì val\_loss của nó cao hơn epoch 12 một lượng không
đáng kể (0,4537 so với 0,4478, chênh 1,3\%).

Nói cách khác, mô hình \emph{đã từng tốt hơn} trong quá trình huấn luyện nhưng
trạng thái đó không được lưu lại. Toàn bộ kết quả V3 báo cáo ở
\cref{tab:seg_results} vì thế là \emph{ước lượng thấp} so với khả năng thực của
phương pháp.

Nguyên nhân gốc đã được dự báo về mặt lý thuyết ở \cref{sec:loss}: hàm mất mát
FocalDice chủ động thu nhỏ gradient của các điểm ảnh dễ, khiến giá trị mất mát
tổng thể trở thành chỉ báo kém nhạy cho chất lượng phân đoạn. Hàm mất mát và IoU
không đồng pha với nhau.

Hiện tượng lặp lại ở V2 với mức độ còn nghiêm trọng hơn. Điểm kiểm tra được chọn
là epoch 20 với val\_IoU = 0,4960, trong khi epoch 18 đạt 0,5438. Đáng nói hơn,
epoch 20 là một \emph{điểm trũng cục bộ} của val\_IoU: các epoch lân cận đều cao
hơn hẳn (E18: 0,5438; E19: 0,5435; E21: 0,5379; E22: 0,5404). Nói cách khác, cơ
chế chọn theo val\_loss đã chọn đúng vào epoch có IoU kém nhất trong cả vùng.

\begin{table}[H]
\centering
\caption{Tổn thất do chọn điểm kiểm tra theo \texttt{val\_loss} ở ba phiên bản}
\label{tab:earlystop_loss}
\begin{tabular}{l l c c c c}
\toprule
\textbf{Phiên bản} & \textbf{Hàm mất mát} & \textbf{Epoch chọn} & \textbf{val\_IoU} & \textbf{val\_IoU tốt nhất} & \textbf{Tổn thất} \\
\midrule
V1 & BCE + Dice & E23 & 0,7699 & 0,7715 (E25) & $0{,}0016$ \\
V2 & FocalDice  & E20 & 0,4960 & 0,5438 (E18) & $\mathbf{0{,}0478}$ \\
V3 & FocalDice  & E12 & 0,6479 & 0,6981 (E17) & $\mathbf{0{,}0502}$ \\
\bottomrule
\end{tabular}
\end{table}

\Cref{tab:earlystop_loss} cho thấy một quy luật rất rõ và đây là bằng chứng thực
nghiệm trực tiếp cho dự báo lý thuyết ở \cref{sec:loss}. Phiên bản V1 dùng
BCE + Dice có tổn thất gần như bằng không (0,0016) --- hàm mất mát và IoU đồng pha
với nhau, chọn theo tiêu chí nào cũng cho kết quả tương đương. Hai phiên bản V2 và
V3 dùng FocalDice có tổn thất lớn gấp \emph{ba mươi lần} (0,0478 và 0,0502).

Cơ chế đằng sau đã được nêu ở \cref{sec:loss}: thành phần Focal chủ động thu nhỏ
gradient --- và qua đó thu nhỏ cả đóng góp vào giá trị hàm mất mát --- của những
điểm ảnh đã được phân loại tốt. Phần lớn diện tích ảnh là nền dễ, nên giá trị hàm
mất mát bị chi phối bởi một nhóm nhỏ điểm ảnh khó ở vùng biên. Trong khi đó IoU
tính trên toàn bộ vùng. Hai đại lượng vì thế đo hai khía cạnh khác nhau và không
còn đồng biến.

Khuyến nghị rút ra rất rõ ràng và có thể áp dụng ngay cho các nghiên cứu tương tự:
\emph{với các bài toán phân đoạn dùng hàm mất mát họ Focal, tiêu chí chọn điểm
kiểm tra và dừng sớm phải là chính chỉ số đánh giá mục tiêu} --- ở đây là val\_IoU
--- \emph{chứ không phải giá trị hàm mất mát}. Với hàm mất mát không có thành phần
điều biến theo độ khó như BCE + Dice, việc chọn theo hàm mất mát vẫn chấp nhận
được.

\subsection{Lỗi tăng cường dữ liệu ở phiên bản V1}

Trong cài đặt của V1, thuộc tính bật tăng cường được đặt trên đối tượng tập dữ
liệu \emph{dùng chung} cho cả ba tập con. Hệ quả là tập kiểm định và kiểm thử cũng
bị áp dụng tăng cường, dù chỉ gồm hai phép lật ngang và lật dọc.

Ảnh hưởng tới kết quả là nhỏ, vì phép lật bảo toàn diện tích và hình dạng của cả
ảnh lẫn mặt nạ, không làm thay đổi bản chất bài toán. Tuy nhiên đây vẫn là một sai
sót phương pháp cần ghi nhận: chỉ số của V1 mang thêm một thành phần dao động ngẫu
nhiên mà V2 và V3 không có. Đây là lý do thứ ba, bên cạnh khác mốc tham chiếu và
khác tập kiểm thử, khiến V1 không nên được dùng làm mốc so sánh.

\subsection{Dense CRF không khả dụng}

Như đã nêu ở \cref{sec:xulyanh}, Dense CRF là một hướng tinh chỉnh nhãn thay thế
cho \SAM{}, có ưu điểm là không cần mô hình nền tảng nặng và bám theo biên màu sắc
thực của ảnh.

Hướng này đã được thử nghiệm nhưng không đưa được vào quy trình: thư viện
\texttt{pydensecrf} không cài đặt được trên môi trường Python 3.13 dùng trong
nghiên cứu, do thư viện này chưa cung cấp bản dựng sẵn cho phiên bản Python đó.
Việc so sánh \SAM{} với Dense CRF trong vai trò bộ tinh chỉnh nhãn vì vậy vẫn là
một câu hỏi bỏ ngỏ, được đề xuất là hướng phát triển ở \cref{chap:ketluan}.

\subsection{Hạn chế của phép biến đổi đàn hồi}

Như đã lưu ý ở \cref{tab:augment}, bộ tăng cường của V2 và V3 gồm
\en{ElasticTransform} với xác suất 0,3 và \en{GridDistortion} với xác suất 0,2.
Hai phép này làm biến dạng \emph{đồng thời} cả ảnh lẫn mặt nạ.

Với nhãn chuẩn, đây là cách tăng cường hợp lệ. Với nhãn giả vốn đã không chắc chắn
về biên, việc biến dạng thêm có thể làm nhiễu tín hiệu huấn luyện thay vì tăng khả
năng khái quát. Giả thuyết này chưa được kiểm chứng bằng thực nghiệm tách bạch
trong luận văn, và được ghi nhận là một hướng khảo sát.

\section{Kết luận chương}

Chương này đã trình bày và phân tích toàn bộ kết quả thực nghiệm. Các phát hiện
chính:

\begin{enumerate}
    \item \textbf{Phân loại đạt 97,19\% trên tập kiểm thử} với $F_1$ 97,19\%, cả
          năm lớp đều vượt $F_1$ 95\%. Cặp nhầm lẫn chính là Algal Leaf Spot và
          Phomopsis Leaf Spot, đúng như dự đoán từ phân tích khám phá.

    \item \textbf{Bản đồ kích hoạt của lớp lá khoẻ khuếch tán khắp phiến lá}, xác
          nhận bằng thực nghiệm rằng bản đồ kích hoạt phản ánh cơ sở ra quyết định
          của mô hình chứ không phải vị trí tổn thương --- cơ sở cho biện pháp ép
          mặt nạ lớp này về không.

    \item \textbf{Ngưỡng phân vị đạt độ phủ mục tiêu với sai lệch dưới 0,4 điểm
          phần trăm}, đồng thời làm độ lệch chuẩn của độ phủ giảm khoảng năm lần
          --- xác nhận cả thành công lẫn đánh đổi đã dự báo.

    \item \textbf{SAM mở rộng độ phủ mặt nạ 69,1\%} thay vì chỉ tinh chỉnh biên,
          với 200 mặt nạ vượt 50\% diện tích ảnh.

    \item \textbf{Cơ chế bảo vệ dựa trên IoU không thể phát hiện hiện tượng nở} ---
          đã chứng minh bằng lập luận toán học rằng ngưỡng cần thiết để kích hoạt
          vượt quá 100\% diện tích ảnh với mọi lớp bệnh. Một ràng buộc bất đối xứng
          về tỉ lệ độ phủ được đề xuất để khắc phục.

    \item \textbf{Tinh chỉnh bằng SAM cải thiện V2 lên V3 trên mọi chỉ số}
          (IoU +21,5\%), nhưng ba yếu tố gây nhiễu đều tác động cùng chiều nên con
          số này là cận trên chứ không phải giá trị thực.

    \item \textbf{Chỉ số của V1 không so sánh được với V2 và V3} do khác biệt ở
          tám biến, bao gồm cả tập kiểm thử. IoU cao của V1 phản ánh mục tiêu dễ,
          không phản ánh chất lượng định vị bệnh.

    \item \textbf{Tiêu chí dừng sớm theo val\_loss là không phù hợp} với hàm mất
          mát họ Focal --- V3 bỏ lỡ điểm kiểm tra tốt hơn 0,0502 IoU.
\end{enumerate}

\Cref{chap:ketluan} tổng hợp các kết quả này thành câu trả lời trực tiếp cho bốn
câu hỏi nghiên cứu đặt ra ở \cref{sec:muctieu}, đồng thời đề xuất hướng phát triển.

% !TeX root = ../main.tex
\chapter{KẾT LUẬN VÀ HƯỚNG PHÁT TRIỂN}
\label{chap:ketluan}

\section{Tổng kết}

Luận văn đã xây dựng và đánh giá một quy trình học giám sát yếu hoàn chỉnh cho bài
toán phát hiện và phân đoạn bệnh lá sầu riêng, đi từ phân loại ảnh qua giải thích
bằng bản đồ kích hoạt, sinh nhãn giả, tinh chỉnh biên bằng mô hình nền tảng, đến
huấn luyện mô hình phân đoạn --- toàn bộ mà không sử dụng bất kỳ nhãn thủ công nào
ở mức điểm ảnh.

Kết quả chính được tóm tắt ở \cref{tab:tongket}.

\begin{table}[H]
\centering
\caption{Tổng hợp kết quả chính của luận văn}
\label{tab:tongket}
\small
\begin{tabular}{L{5.2cm} L{8.0cm}}
\toprule
\textbf{Hạng mục} & \textbf{Kết quả} \\
\midrule
Phân loại năm lớp bệnh & Accuracy 97,19\%, $F_1$ 97,19\% trên 890 ảnh kiểm thử \\
Độ chính xác kiểm định tốt nhất & 98,19\% tại epoch 11 \\
Sinh nhãn giả & 3.104 mặt nạ, không dùng nhãn thủ công nào \\
Kiểm soát độ phủ (V2) & Sai lệch dưới 0,4 điểm phần trăm mọi lớp \\
Loại nhãn nhiễu lớp lá khoẻ & 683 mặt nạ dương tính giả bị loại bỏ \\
Hiệu quả tinh chỉnh bằng SAM & IoU tăng 21,5\% (cận trên), hội tụ sớm hơn 8 epoch \\
Phát hiện về cơ chế bảo vệ IoU & Chứng minh không thể phát hiện hiện tượng nở mặt nạ \\
Phát hiện về tiêu chí dừng sớm & Mất 0,048--0,050 IoU khi dùng FocalDice \\
\bottomrule
\end{tabular}
\end{table}

\section{Trả lời các câu hỏi nghiên cứu}

\subsection{CH1 --- Hiệu năng phân loại và các cặp lớp dễ nhầm}

\emph{Câu hỏi: Mô hình \EffNet{} đạt độ chính xác bao nhiêu trên năm lớp bệnh lá
sầu riêng, và những cặp lớp nào dễ bị nhầm lẫn nhất?}

Mô hình đạt \textbf{97,19\% độ chính xác} trên tập kiểm thử 890 ảnh, với $F_1$ có
trọng số 97,19\%. Độ chính xác kiểm định tốt nhất là 98,19\% tại epoch 11; khoảng
cách 1,00 điểm phần trăm giữa hai giá trị cho thấy mô hình khái quát tốt. Cả năm
lớp đều đạt $F_1$ trên 95\%, biên độ giữa lớp cao nhất và thấp nhất chỉ 2,66 điểm
phần trăm.

Cặp nhầm lẫn chính là \textbf{Algal Leaf Spot và Phomopsis Leaf Spot}, chiếm 5
trên tổng số 24 lỗi. Kết quả này đã được dự báo từ giai đoạn phân tích khám phá dữ
liệu, dựa trên hai bằng chứng độc lập: quan sát hình thái cho thấy cả hai bệnh đều
tạo đốm nhỏ trên nền lá xanh, và phân tích màu cho thấy hai lớp có phân bố kênh
màu chồng lấn nhiều nhất.

Đáng chú ý về mặt ứng dụng, tỉ lệ bỏ sót bệnh --- ảnh bệnh bị phân loại thành lá
khoẻ --- chỉ là 6 trên 694 ảnh, tức 0,86\%.

\subsection{CH2 --- Chất lượng của bản đồ chú ý}

\emph{Câu hỏi: Bản đồ chú ý sinh ra từ mô hình phân loại có tập trung vào vùng tổn
thương thực sự hay không, và mức độ tin cậy của chúng khác nhau thế nào giữa các
lớp bệnh?}

Câu trả lời phân hoá rõ rệt theo lớp, và chính sự phân hoá này là kết quả có giá
trị nhất.

Với \textbf{bốn lớp bệnh}, bản đồ chú ý tập trung đúng vào vùng tổn thương: kích
hoạt bám vào các đốm với Algal và Phomopsis, bám vào vết cắn với Allocaridara, và
trải theo mảng hoại tử với Leaf Blight. Độ tin cậy trung bình trên tập huấn luyện
đạt 0,971--0,991.

Với \textbf{lớp lá khoẻ}, bản đồ chú ý \emph{khuếch tán khắp phiến lá}, không có
vùng trọng tâm. Hệ quả định lượng là mặt nạ sinh ra từ nhãn giả V1 cho lớp này có
độ phủ trung bình 26,4\% --- cao gần gấp đôi lớp bệnh lan rộng nhất --- với trường
hợp cực đoan lên tới 59,1\% trên một chiếc lá hoàn toàn khoẻ mạnh.

Đây là bằng chứng thực nghiệm trực tiếp cho luận điểm lý thuyết trung tâm của luận
văn: \textbf{bản đồ chú ý trả lời câu hỏi ``mô hình nhìn vào đâu để phân loại'',
không phải câu hỏi ``vùng bệnh nằm ở đâu''}. Với lá khoẻ, mô hình nhận diện qua sự
\emph{vắng mặt} của dấu hiệu bệnh --- một hành vi hoàn toàn hợp lý cho nhiệm vụ
phân loại nhưng vô nghĩa khi chuyển thành mặt nạ tổn thương.

Nhận thức này dẫn tới biện pháp ép mặt nạ lớp lá khoẻ về không, loại bỏ 683 mặt nạ
nhiễu và cung cấp cho mô hình phân đoạn các mẫu âm sạch.

\subsection{CH3 --- Hiệu quả của việc tinh chỉnh biên bằng SAM}

\emph{Câu hỏi: Việc tinh chỉnh biên bằng \SAM{} có cải thiện chất lượng nhãn giả,
và qua đó cải thiện mô hình phân đoạn, hay không?}

Câu trả lời có hai vế đối lập nhau và cả hai đều đúng.

\textbf{Xét theo tiêu chí bệnh học, \SAM{} làm nhãn xấu đi.} Thay vì chỉ tinh chỉnh
biên như kỳ vọng, \SAM{} mở rộng độ phủ mặt nạ trung bình từ 15,04\% lên 25,43\%
--- tăng 69,1\%, với mức tăng theo lớp từ 41,5\% đến 95,7\%. Có 200 trên 3.104 mặt
nạ vượt quá 50\% diện tích ảnh, điều bất khả thi với bệnh đốm lá. Nguyên nhân là
\SAM{} được huấn luyện để phân đoạn đơn vị ngữ nghĩa tự nhiên, mà đơn vị đó với
gợi ý điểm đặt trên đốm bệnh thường là \emph{cả chiếc lá}.

\textbf{Xét theo hiệu quả huấn luyện, \SAM{} cải thiện rõ rệt.} Với kiến trúc, hàm
mất mát và tập kiểm thử giống hệt nhau, phiên bản V3 vượt V2 trên toàn bộ bốn chỉ
số. Bằng chứng thuyết phục nhất không phải con số IoU mà là tốc độ hội tụ: V3 đạt
val\_loss 0,4478 tại epoch 12, trong khi V2 phải tới epoch 20 mới đạt 0,5957.

Hai vế này không mâu thuẫn vì chúng đo hai thứ khác nhau. Nhãn \SAM{} \emph{sai về
phạm vi} nhưng \emph{đúng về cấu trúc} --- biên bám theo ranh giới thực trong ảnh
và nhất quán giữa các ảnh. Mô hình học được từ tính nhất quán cấu trúc đó.

Về mức độ cải thiện, luận văn không kết luận toàn bộ 21,5\% là do \SAM{}. Ba yếu
tố gây nhiễu --- V2 bị dừng sớm khi còn đang cải thiện, thiên lệch của IoU theo
kích thước mục tiêu, và khác biệt kích thước lô --- đều tác động cùng chiều làm
phóng đại lợi thế của V3. Con số 21,5\% vì vậy là \textbf{cận trên} chứ không phải
giá trị thực.

\subsection{CH4 --- Khả năng so sánh của các chỉ số khi thiếu nhãn chuẩn}

\emph{Câu hỏi: Trong điều kiện không có nhãn chuẩn ở mức điểm ảnh, các chỉ số IoU
và Dice giữa những phiên bản khác nhau có so sánh trực tiếp được với nhau không?
Nếu không thì phép so sánh nào là hợp lệ?}

\textbf{Câu trả lời là không}, và đây là kết luận phương pháp luận quan trọng nhất
của luận văn.

Mỗi phiên bản được đánh giá trên mốc tham chiếu của chính nó: V1 so với nhãn V1,
V2 so với nhãn V2, V3 so với nhãn đã qua tinh chỉnh. Ba mốc tham chiếu khác nhau về
hình dạng, độ phủ lẫn độ biến thiên. So sánh trực tiếp các con số IoU giữa chúng
tương đương với việc so điểm của ba học sinh làm ba đề thi có độ khó khác nhau.

Minh hoạ rõ nhất là phiên bản V1 với IoU 0,7671 --- cao nhất trong ba phiên bản.
Con số này \emph{không} có nghĩa V1 định vị bệnh tốt nhất. Nó phản ánh việc mục
tiêu của V1 dễ hơn: nhãn V1 là các khối tròn lớn và trơn, hệ quả của việc phóng
bản đồ $7 \times 7$ lên $224 \times 224$, nên dễ tái tạo hơn nhiều so với các biên
sắc và phức tạp của nhãn \SAM{}.

Ngoài khác biệt về mốc tham chiếu, luận văn còn xác định thêm hai lý do kỹ thuật
khiến V1 không thể so sánh với V2 và V3: hai nhóm dùng \emph{hai tập kiểm thử khác
nhau} do khác bộ sinh số ngẫu nhiên dù cùng hạt giống, và cài đặt của V1 vô tình
áp dụng tăng cường dữ liệu lên cả tập kiểm định và kiểm thử. Tổng cộng cặp V1--V2
khác nhau ở \textbf{tám biến} cùng lúc.

\textbf{Phép so sánh hợp lệ duy nhất là V2 với V3}, vì hai phiên bản này dùng chung
kiến trúc, hàm mất mát, bộ tối ưu, bộ tăng cường và cùng một tập kiểm thử. Ngay cả
với cặp này, ba yếu tố gây nhiễu còn lại vẫn cần được trừ hao khi diễn giải.

Kết luận tổng quát: \textbf{trong khung học giám sát yếu không có nhãn chuẩn, chỉ
số IoU và Dice chỉ đo được mức độ mô hình tái tạo lại bộ nhãn giả của chính nó,
không đo được độ chính xác định vị bệnh tuyệt đối.} Mọi so sánh giữa các phiên bản
dùng bộ nhãn khác nhau đều cần được đặt trong ngữ cảnh này.

\section{Đóng góp của luận văn}

\subsection{Đóng góp về phương pháp}

\textbf{Cơ chế ngưỡng phân vị theo lớp.} Thay vì cố định giá trị ngưỡng và để độ
phủ trở thành kết quả không kiểm soát được, luận văn cố định tỉ lệ độ phủ và suy
ngược ra ngưỡng cho từng ảnh. Cơ chế này đạt độ phủ mục tiêu với sai lệch dưới 0,4
điểm phần trăm ở mọi lớp, đồng thời tránh được hiện tượng phân đoạn thừa của ngưỡng
Otsu khi bản đồ nhiệt hợp nhất đa tỉ lệ trở nên trơn và đơn đỉnh.

\textbf{Đưa tri thức tiên nghiệm vào quy trình sinh nhãn.} Việc ép mặt nạ lớp lá
khoẻ về không là một can thiệp đơn giản nhưng hiệu quả, loại bỏ 683 mặt nạ dương
tính giả mà bản thân bản đồ kích hoạt không bao giờ tự loại được.

\subsection{Đóng góp về phát hiện thực nghiệm}

\textbf{Chứng minh cơ chế bảo vệ dựa trên IoU không thể phát hiện hiện tượng nở
mặt nạ.} Khi mặt nạ mới chứa trọn mặt nạ gốc, IoU giữa chúng bằng đúng tỉ số hai
độ phủ. Do đó cơ chế bảo vệ với ngưỡng $\theta$ chỉ kích hoạt khi độ phủ mới vượt
$1/\theta$ lần độ phủ gốc. Với $\theta = 0{,}15$ và độ phủ thực tế 17,6--21,7\%,
ngưỡng cần thiết là 117--145\% diện tích ảnh --- bất khả thi về mặt vật lý. Đây là
khiếm khuyết thiết kế mang tính hệ thống, không phải lỗi cài đặt, và bắt nguồn từ
tính đối xứng của IoU: nó không phân biệt được ``lệch vị trí'' với ``đúng vị trí
nhưng quá rộng''. Luận văn đề xuất bổ sung một ràng buộc bất đối xứng về tỉ lệ độ
phủ để khắc phục.

\textbf{Xác định hệ quả của việc chọn điểm kiểm tra theo hàm mất mát khi dùng
FocalDice.} Với phiên bản V1 dùng BCE + Dice, tổn thất do chọn theo hàm mất mát
thay vì theo IoU chỉ là 0,0016 --- không đáng kể. Với hai phiên bản dùng FocalDice,
tổn thất lớn gấp ba mươi lần: 0,0478 và 0,0502. Nguyên nhân là thành phần Focal
thu nhỏ đóng góp của các điểm ảnh dễ, khiến giá trị hàm mất mát bị chi phối bởi
một nhóm nhỏ điểm ảnh biên trong khi IoU tính trên toàn vùng. Khuyến nghị rút ra
có thể áp dụng trực tiếp cho các nghiên cứu tương tự.

\subsection{Đóng góp về phương pháp luận}

\textbf{Khung phân tích tách bạch cho học giám sát yếu.} Luận văn đề xuất và áp
dụng cách tiếp cận đếm số biến khác nhau giữa các cặp cấu hình trước khi diễn giải
chênh lệch hiệu năng. Cách làm này cho thấy cặp V1--V2 khác nhau ở tám biến nên
không thể diễn giải, còn cặp V2--V3 khác ba biến nên diễn giải được với điều kiện
trừ hao. Đồng thời luận văn cũng đánh giá \emph{chiều tác động} của từng yếu tố
gây nhiễu để xác định con số quan sát được là cận trên hay cận dưới.

\textbf{Tài liệu kiểm kê số liệu.} Toàn bộ con số trong luận văn được truy vết về
đúng notebook và ô lệnh sinh ra nó, kèm danh sách các điểm chưa đủ bằng chứng.
Cách làm này bảo đảm tính minh bạch và cho phép kiểm chứng độc lập.

\section{Hạn chế}

\subsection{Hạn chế cốt lõi --- thiếu nhãn chuẩn}

Hạn chế bao trùm và quyết định phạm vi mọi kết luận là bộ dữ liệu không có nhãn
thủ công ở mức điểm ảnh. Hệ quả là:

\begin{itemize}
    \item Không thể xác lập hiệu năng định vị bệnh \emph{tuyệt đối} của bất kỳ
          phiên bản nào.
    \item Không thể hiệu chỉnh các tỉ lệ độ phủ mục tiêu 18--22\%, vốn hiện được
          chọn theo suy luận bệnh học.
    \item Không thể kiểm chứng liệu hiện tượng \SAM{} mở rộng mặt nạ có thực sự
          làm nhãn xa rời tổn thương thật hay không --- kết luận hiện tại dựa trên
          quan sát định tính và lập luận về diện tích hợp lý.
\end{itemize}

\subsection{Hạn chế về thiết kế thực nghiệm}

Các phiên bản không được thiết kế như một thí nghiệm có kiểm soát ngay từ đầu.
Cặp V1--V2 khác nhau ở tám biến; ngay cả cặp V2--V3 vẫn còn hai khác biệt ngoài ý
muốn về kích thước lô và patience. Ràng buộc bộ nhớ đồ hoạ 2,15 GB là nguyên nhân
trực tiếp của một số khác biệt này.

Ngoài ra, mỗi cấu hình chỉ được chạy \emph{một lần} với một hạt giống ngẫu nhiên
duy nhất. Không có ước lượng phương sai, nên không thể khẳng định chênh lệch quan
sát được vượt ngưỡng dao động ngẫu nhiên giữa các lần chạy.

\subsection{Hạn chế về dữ liệu}

Toàn bộ ảnh đã được co về $224 \times 224$ và bị xoá siêu dữ liệu EXIF trước khi
đến với nghiên cứu, nên không thể khảo sát ảnh hưởng của độ phân giải đầu vào hay
điều kiện chụp. Hệ thống cũng chỉ được đánh giá trên ảnh chụp lá đơn lẻ, chưa kiểm
chứng trên ảnh chụp toàn tán trong điều kiện đồng ruộng --- nơi có nền phức tạp,
nhiều lá chồng lấn và điều kiện ánh sáng biến thiên.

\subsection{Hạn chế về phạm vi bài toán}

Bài toán phân đoạn được đặt ở dạng nhị phân --- bệnh so với nền --- nên mô hình
phân đoạn không phân biệt được loại bệnh ở mức điểm ảnh. Với ảnh có nhiều loại
bệnh cùng lúc, hệ thống hiện chỉ cho một mặt nạ chung.

\section{Hướng phát triển}

\subsection{Ưu tiên cao}

\textbf{Xây dựng tập kiểm thử có nhãn chuẩn.} Đây là việc cần làm trước tiên vì nó
gỡ bỏ hạn chế cốt lõi. Chỉ cần một tập nhỏ 50--100 ảnh được chuyên gia khoanh vùng
thủ công là đủ để: đo hiệu năng định vị bệnh tuyệt đối, so sánh công bằng ba phiên
bản trên cùng một mốc tham chiếu, và hiệu chỉnh các tỉ lệ độ phủ mục tiêu. Chi phí
cho quy mô này hoàn toàn khả thi.

\textbf{Sửa cơ chế bảo vệ và chạy lại giai đoạn tinh chỉnh.} Bổ sung ràng buộc bất
đối xứng về tỉ lệ độ phủ như đã đề xuất, với hệ số $\kappa = 2{,}0$, rồi sinh lại
nhãn và huấn luyện lại. Đây là thí nghiệm rẻ nhưng cho phép kiểm chứng trực tiếp
liệu việc ngăn hiện tượng nở có cải thiện chất lượng nhãn hay không.

\textbf{Chuyển tiêu chí chọn điểm kiểm tra sang val\_IoU.} Chỉ là thay đổi vài
dòng mã nhưng theo số liệu đã đo sẽ thu hồi được khoảng 0,05 IoU cho phiên bản V3.

\subsection{Ưu tiên trung bình}

\textbf{Thiết kế lại thí nghiệm có kiểm soát chặt.} Cố định toàn bộ siêu tham số
giữa các phiên bản và chỉ thay đúng một biến mỗi lần, đồng thời chạy nhiều hạt
giống ngẫu nhiên để có ước lượng phương sai.

\textbf{So sánh \SAM{} với Dense CRF.} Dense CRF là hướng tinh chỉnh nhãn thay thế,
không cần mô hình nền tảng nặng và bám theo biên màu sắc thực. Hướng này đã được
thử nhưng không triển khai được do thư viện không tương thích với phiên bản Python
sử dụng; việc so sánh vẫn là câu hỏi bỏ ngỏ.

\textbf{Khảo sát cường độ tăng cường dữ liệu.} Kiểm chứng giả thuyết rằng các phép
biến dạng đàn hồi làm nhiễu tín hiệu huấn luyện khi nhãn vốn đã không chắc chắn về
biên.

\textbf{Thử \SAM{}2.} Phiên bản kế tiếp của \SAM{} có chất lượng biên tốt hơn và
có thể giảm mức độ nở mặt nạ.

\subsection{Hướng mở rộng dài hạn}

\textbf{Phân đoạn đa lớp.} Mở rộng từ nhị phân sang phân đoạn theo từng loại bệnh,
cho phép xử lý ảnh có nhiều bệnh đồng thời.

\textbf{Ước lượng mức độ bệnh.} Khắc phục hạn chế đã nêu rằng nhãn giả V2 không mã
hoá được mức độ nặng nhẹ, hướng tới đầu ra là tỉ lệ diện tích lá bị bệnh có hiệu
chỉnh --- đại lượng trực tiếp phục vụ quyết định canh tác.

\textbf{Kiểm chứng trong điều kiện đồng ruộng.} Đánh giá trên ảnh chụp toàn tán với
nền phức tạp và điều kiện ánh sáng biến thiên, đối chiếu với bài học từ bộ dữ liệu
PlantVillage về khoảng cách giữa dữ liệu chuẩn hoá và dữ liệu thực
tế~\cite{mohanty2016plantdisease}.

\textbf{Triển khai trên thiết bị biên.} Lượng tử hoá và cắt tỉa mô hình để chạy
trực tiếp trên điện thoại, phục vụ chẩn đoán tại vườn không cần kết nối mạng.

\section{Lời kết}

Luận văn cho thấy một quy trình học giám sát yếu hoàn toàn có thể chuyển nhãn mức
ảnh --- vốn rẻ và sẵn có --- thành mặt nạ phân đoạn ở mức điểm ảnh mà không cần bất
kỳ công khoanh vùng thủ công nào. Mô hình phân loại đạt 97,19\% độ chính xác, và
toàn bộ 3.104 mặt nạ giả được sinh tự động.

Nhưng đóng góp mà luận văn coi trọng nhất không nằm ở các con số hiệu năng. Đó là
việc chỉ ra rằng \emph{trong khung học giám sát yếu, chính các con số ấy cần được
đọc một cách thận trọng}. Phiên bản có IoU cao nhất không phải phiên bản định vị
bệnh tốt nhất. Một mô hình nền tảng có thể vừa làm hỏng nhãn theo tiêu chí bệnh học
vừa cải thiện kết quả huấn luyện. Và một cơ chế bảo vệ tưởng như hợp lý lại có thể
bất lực về mặt toán học trước đúng vấn đề mà nó được thiết kế để ngăn chặn.

Những phát hiện này chỉ lộ ra khi đối chiếu kỹ số liệu thực tế với kỳ vọng thiết
kế, thay vì dừng lại ở con số tốt nhất đạt được. Đó cũng là tinh thần mà luận văn
mong muốn đóng góp cho các nghiên cứu học giám sát yếu tiếp theo trong lĩnh vực
nông nghiệp.


% ===========================================================================
%  TÀI LIỆU THAM KHẢO
% ===========================================================================
% \sloppy nới lỏng ràng buộc giãn chữ trong đoạn: cần thiết cho danh mục tham
% khảo vì tên tạp chí và tên hội nghị dài không phải lúc nào cũng ngắt vừa dòng.
\begingroup
\sloppy
\printbibliography[heading=bibintoc, title={TÀI LIỆU THAM KHẢO}]
\endgroup

% ===========================================================================
%  PHỤ LỤC
% ===========================================================================
\appendix
% !TeX root = ../main.tex
\chapter{BẢNG SIÊU THAM SỐ ĐẦY ĐỦ}
\label{app:A}

Phụ lục này tập hợp toàn bộ siêu tham số của các thực nghiệm trong luận văn, phục
vụ mục đích tái tạo kết quả. Mọi giá trị được trích trực tiếp từ mã nguồn thực thi.

\section{Mô hình phân loại}

\begin{table}[H]
\centering
\caption{Siêu tham số mô hình phân loại EfficientNet-B0}
\label{tab:app_cls}
\begin{tabular}{L{5.4cm} L{7.6cm}}
\toprule
\textbf{Tham số} & \textbf{Giá trị} \\
\midrule
\multicolumn{2}{l}{\textit{Kiến trúc}} \\
Mạng xương sống       & EfficientNet-B0 \\
Trọng số khởi tạo     & ImageNet (IMAGENET1K\_V1) \\
Tầng phân loại        & Dropout($p = 0{,}2$) + Linear($1280 \rightarrow 5$) \\
Tổng tham số          & 4.013.953 \\
Tham số khả huấn luyện & 4.013.953 (100\%) \\
Kích thước đầu vào    & $224 \times 224 \times 3$ \\
\midrule
\multicolumn{2}{l}{\textit{Huấn luyện}} \\
Hàm mất mát           & CrossEntropyLoss \\
Bộ tối ưu             & Adam \\
Tốc độ học ban đầu    & $1 \times 10^{-4}$ \\
Lịch tốc độ học       & ReduceLROnPlateau \\
\quad chế độ          & \texttt{max} (theo \texttt{val\_acc}) \\
\quad patience        & 2 \\
\quad hệ số giảm      & 0,5 \\
Kích thước lô         & 16 \\
Số epoch tối đa       & 50 \\
Số epoch thực tế      & 16 \\
Dừng sớm --- chỉ số   & \texttt{val\_acc} \\
Dừng sớm --- patience & 5 \\
\midrule
\multicolumn{2}{l}{\textit{Dữ liệu}} \\
Chia tập              & 70 / 10 / 20 \\
Số ảnh                & 3.104 / 443 / 890 \\
Chuẩn hoá             & Trung bình và độ lệch chuẩn ImageNet \\
Tăng cường (chỉ train) & Lật ngang, lật dọc, xoay, biến đổi màu \\
\bottomrule
\end{tabular}
\end{table}

\section{Sinh nhãn giả}

\begin{table}[H]
\centering
\caption{Siêu tham số sinh nhãn giả V1 và V2}
\label{tab:app_pseudo}
\begin{tabular}{L{4.4cm} L{4.2cm} L{4.4cm}}
\toprule
\textbf{Tham số} & \textbf{V1} & \textbf{V2} \\
\midrule
Phương pháp CAM   & Grad-CAM & Grad-CAM++ \\
Số tầng đích      & 1 & 2 \\
Tầng thô          & \texttt{base.\allowbreak features.\allowbreak 8.\allowbreak 0} & \texttt{base.\allowbreak features.\allowbreak 8.\allowbreak 0} \\
Độ phân giải thô  & $7 \times 7$ & $7 \times 7$ \\
Tầng tinh         & --- & \texttt{base.\allowbreak features.} \texttt{6.\allowbreak 3.\allowbreak block.\allowbreak 3.\allowbreak 0} \\
Độ phân giải tinh & --- & $14 \times 14$ \\
Trọng số hợp nhất & --- & $0{,}6 \times$ tinh $+\ 0{,}4 \times$ thô \\
Ngưỡng hoá        & Cố định $\geq 0{,}5$ & Phân vị theo lớp \\
Phần tử cấu trúc  & Chữ nhật $5 \times 5$ & Elip $5 \times 5$ \\
Thứ tự hình thái  & Đóng $\rightarrow$ Mở & Đóng $\rightarrow$ Mở \\
Lọc thành phần liên thông & Không & Top-3, diện tích $\geq 400$ px \\
Xử lý lớp lá khoẻ & Như lớp bệnh & Mặt nạ toàn số 0 \\
Số mặt nạ sinh ra & 3.104 & 3.104 \\
Thời gian sinh    & 14 ph 00 s & 4 ph 08 s \\
Tốc độ            & 3,69 ảnh/s & 12,51 ảnh/s \\
\bottomrule
\end{tabular}
\end{table}

\begin{table}[H]
\centering
\caption{Tỉ lệ độ phủ mục tiêu theo lớp (nhãn giả V2)}
\label{tab:app_keeppct}
\begin{tabular}{l c}
\toprule
\textbf{Lớp} & \textbf{Tỉ lệ giữ lại} \\
\midrule
ALGAL\_LEAF\_SPOT      & 0,18 \\
ALLOCARIDARA\_ATTACK   & 0,20 \\
HEALTHY\_LEAF          & 0,00 \\
LEAF\_BLIGHT           & 0,22 \\
PHOMOPSIS\_LEAF\_SPOT  & 0,18 \\
\bottomrule
\end{tabular}
\end{table}

\section{Tinh chỉnh bằng SAM}

\begin{table}[H]
\centering
\caption{Siêu tham số giai đoạn tinh chỉnh bằng SAM}
\label{tab:app_sam}
\begin{tabular}{L{5.6cm} L{7.4cm}}
\toprule
\textbf{Tham số} & \textbf{Giá trị} \\
\midrule
Mô hình              & SAM ViT-B \\
Tệp trọng số         & \texttt{sam\_\allowbreak vit\_\allowbreak b\_\allowbreak 01ec64.\allowbreak pth} (375 MB) \\
Số điểm tiền cảnh    & 5 (trọng tâm + 4 điểm ngẫu nhiên) \\
Số điểm nền          & 3 \\
Bán kính giãn nở vùng nền & 30 điểm ảnh \\
Chế độ sinh mặt nạ   & \texttt{multimask\_\allowbreak output = True}, chọn mặt nạ điểm cao nhất \\
Ngưỡng bảo vệ IoU    & 0,15 \\
Lọc thành phần liên thông & Loại nếu $< 5\%$ tổng diện tích hoặc $< 30$ px \\
Hình thái sau lọc    & Đóng, phần tử $5 \times 5$ \\
Xử lý mặt nạ rỗng    & Bỏ qua, giữ nguyên rỗng \\
Số mặt nạ xử lý      & 3.104 \\
\bottomrule
\end{tabular}
\end{table}

\section{Mô hình phân đoạn}

% Dùng longtable thay cho table+tabular: bảng này dài hơn một trang nên nếu để
% trong float [H] sẽ tràn quá đáy trang. longtable tự ngắt sang trang sau và
% lặp lại dòng tiêu đề.
\begingroup
\small
\begin{longtable}{L{3.6cm} L{3.0cm} L{3.0cm} L{3.0cm}}
\caption{Siêu tham số ba phiên bản mô hình phân đoạn}
\label{tab:app_seg} \\
\toprule
\textbf{Tham số} & \textbf{V1} & \textbf{V2} & \textbf{V3} \\
\midrule
\endfirsthead
\multicolumn{4}{l}{\small\itshape (tiếp theo trang trước)} \\
\toprule
\textbf{Tham số} & \textbf{V1} & \textbf{V2} & \textbf{V3} \\
\midrule
\endhead
\midrule
\multicolumn{4}{r}{\small\itshape (còn tiếp)} \\
\endfoot
\bottomrule
\endlastfoot
\multicolumn{4}{l}{\textit{Kiến trúc}} \\
Mô hình        & EfficientNet-UNet & UNet++ & UNet++ \\
Bộ mã hoá      & EfficientNet-B0 & EfficientNet-B0 & EfficientNet-B0 \\
Trọng số mã hoá & ImageNet & ImageNet & ImageNet \\
Tổng tham số   & 7.276.925 & 6.569.581 & 6.569.581 \\
\quad bộ mã hoá & 4.007.548 & 4.007.548 & 4.007.548 \\
\quad bộ giải mã & 3.269.312 & 2.561.888 & 2.561.888 \\
Số kênh đầu ra & 1 & 1 & 1 \\
\midrule
\multicolumn{4}{l}{\textit{Hàm mất mát và tối ưu}} \\
Hàm mất mát    & BCE + Dice & FocalDice & FocalDice \\
\quad $\alpha$ & --- & 0,25 & 0,25 \\
\quad $\gamma$ & --- & 2,0 & 2,0 \\
\quad $\epsilon$ (Dice) & 1,0 & 1,0 & 1,0 \\
Bộ tối ưu      & Adam & AdamW & AdamW \\
Tốc độ học     & $1 \times 10^{-4}$ & $1 \times 10^{-4}$ & $1 \times 10^{-4}$ \\
Suy giảm trọng số & --- & $1 \times 10^{-4}$ & $1 \times 10^{-4}$ \\
Lịch tốc độ học & Reduce\allowbreak LR\allowbreak On\allowbreak Plateau
                & Cosine\allowbreak Annealing\allowbreak LR
                & Cosine\allowbreak Annealing\allowbreak LR \\
\quad tham số  & \texttt{min}, p = 2, hệ số 0,5
               & $T_{\max} = 50$, $\eta_{\min} = 10^{-6}$
               & $T_{\max} = 50$, $\eta_{\min} = 10^{-6}$ \\
Cắt gradient   & Không & \texttt{max\_norm} = 1,0 & \texttt{max\_norm} = 1,0 \\
\midrule
\multicolumn{4}{l}{\textit{Huấn luyện}} \\
Kích thước lô  & 8 & 4 & 8 \\
Epoch tối đa   & 30 & 50 & 50 \\
Epoch thực tế  & 30 & 25 & 22 \\
Dừng sớm --- chỉ số & \texttt{val\_loss} & \texttt{val\_loss} & \texttt{val\_loss} \\
Dừng sớm --- patience & 7 & 5 & 10 \\
Epoch chọn     & 23 & 20 & 12 \\
\midrule
\multicolumn{4}{l}{\textit{Dữ liệu}} \\
Nhãn huấn luyện & Pseudo V1 & Pseudo V2 & SAM tinh chỉnh \\
Chia tập       & 70 / 15 / 15 & 70 / 15 / 15 & 70 / 15 / 15 \\
Số ảnh         & 2.172 / 465 / 467 & 2.172 / 465 / 467 & 2.172 / 465 / 467 \\
Hạt giống      & 42 & 42 & 42 \\
Bộ sinh số ngẫu nhiên & \texttt{torch} & \texttt{numpy} & \texttt{numpy} \\
Ngưỡng nhị phân hoá & 0,5 & 0,5 & 0,5 \\
\end{longtable}
\endgroup

\begin{table}[htbp]
\centering
\caption{Bộ tăng cường dữ liệu Albumentations (V2 và V3)}
\label{tab:app_augment}
\begin{tabular}{L{5.6cm} L{3.6cm} c}
\toprule
\textbf{Phép biến đổi} & \textbf{Tham số} & \textbf{Xác suất} \\
\midrule
RandomResizedCrop & \texttt{size} = (224, 224), \texttt{scale} = (0,7; 1,0) & 1,00 \\
HorizontalFlip    & --- & 0,50 \\
VerticalFlip      & --- & 0,50 \\
RandomRotate90    & --- & 0,50 \\
ShiftScaleRotate  & dịch 0,1; co giãn 0,1; xoay 30 độ & 0,50 \\
ElasticTransform  & mặc định & 0,30 \\
GridDistortion    & mặc định & 0,20 \\
ColorJitter       & sáng, tương phản, bão hoà 0,3; sắc độ 0,1 & 0,50 \\
GaussNoise        & mặc định & 0,20 \\
Normalize         & trung bình, độ lệch chuẩn ImageNet & 1,00 \\
\bottomrule
\end{tabular}
\end{table}

\input{appendix/B-training-logs}
% !TeX root = ../main.tex
\chapter{ỨNG DỤNG MINH HOẠ}
\label{app:C}

Để kiểm chứng khả năng vận hành đầu cuối của quy trình, luận văn xây dựng hai công
cụ minh hoạ: một \en{notebook} suy luận trên ảnh đơn lẻ và một ứng dụng web tương
tác.

\section{Suy luận trên ảnh đơn lẻ}

Quy trình suy luận đầy đủ gồm bốn bước nối tiếp trên một ảnh đầu vào:

\begin{enumerate}
    \item \textbf{Phân loại} --- nạp mô hình \EffNet{} đã huấn luyện, trả về lớp
          bệnh dự đoán kèm phân bố xác suất trên năm lớp.
    \item \textbf{Giải thích} --- sinh bản đồ kích hoạt \GradCAM{} cho lớp được dự
          đoán, chồng lên ảnh gốc.
    \item \textbf{Phân đoạn} --- nạp mô hình \UNetpp{} phiên bản V3, trả về bản đồ
          xác suất và mặt nạ nhị phân sau khi áp ngưỡng 0,5.
    \item \textbf{Tổng hợp} --- ghép kết quả thành một bảng năm khung và tính tỉ lệ
          diện tích tổn thương.
\end{enumerate}

\begin{figure}[H]
\centering
\includegraphics[width=\textwidth]{figures/infer_demo.png}
\caption[Kết quả suy luận đầu cuối trên một ảnh]{Kết quả suy luận đầu cuối trên
một ảnh thuộc lớp Algal Leaf Spot. Từ trái sang phải: ảnh đầu vào, phân bố xác
suất năm lớp, bản đồ kích hoạt Grad-CAM chồng lớp, mặt nạ phân đoạn, và ảnh tổng
hợp với đường bao vùng tổn thương.}
\label{fig:infer_demo}
\end{figure}

Với ảnh minh hoạ ở \cref{fig:infer_demo}, mô hình dự đoán lớp Algal Leaf Spot với
độ tin cậy \textbf{99,29\%}. Phân bố xác suất trên năm lớp cho thấy mức độ tách
bạch rất cao: lớp có xác suất lớn thứ hai là Phomopsis Leaf Spot với 0,62\%, các
lớp còn lại đều dưới 0,05\%. Điều này phù hợp với nhận định ở
\cref{sec:ketqua_cls} rằng Algal và Phomopsis là cặp lớp gần nhau nhất --- ngay cả
khi mô hình rất chắc chắn, lớp cạnh tranh duy nhất vẫn là Phomopsis.

\begin{table}[H]
\centering
\caption{Phân bố xác suất dự đoán trên ảnh minh hoạ}
\label{tab:app_infer}
\begin{tabular}{l r}
\toprule
\textbf{Lớp} & \textbf{Xác suất} \\
\midrule
\textbf{Algal Leaf Spot}   & \textbf{0,9929} \\
Phomopsis Leaf Spot        & 0,0062 \\
Leaf Blight                & 0,0005 \\
Healthy Leaf               & 0,0004 \\
Allocaridara Attack        & 0,0000 \\
\bottomrule
\end{tabular}
\end{table}

\section{Ứng dụng web tương tác}

Ứng dụng được xây dựng bằng Streamlit, cho phép người dùng tải lên ảnh lá và nhận
kết quả phân tích tức thời mà không cần biết lập trình.

\subsection{Chức năng}

\begin{itemize}
    \item Tải ảnh lên từ máy người dùng, hỗ trợ định dạng JPG và PNG.
    \item Chọn phiên bản mô hình phân đoạn trong ba lựa chọn V1, V2, V3 --- phục
          vụ so sánh trực quan giữa các phiên bản.
    \item Điều chỉnh ngưỡng nhị phân hoá bằng thanh trượt trong khoảng 0,1--0,9,
          cho phép quan sát ảnh hưởng của ngưỡng lên diện tích mặt nạ.
    \item Bật hoặc tắt lớp hiển thị bản đồ kích hoạt.
    \item Hiển thị đồng thời năm khung: ảnh gốc, biểu đồ xác suất, bản đồ kích
          hoạt, mặt nạ phân đoạn, và ảnh tổng hợp có đường bao.
    \item Hiển thị bản đồ xác suất thô và thông tin điểm kiểm tra đang dùng.
\end{itemize}

\subsection{Kiến trúc kỹ thuật}

\begin{table}[H]
\centering
\caption{Thành phần kỹ thuật của ứng dụng minh hoạ}
\label{tab:app_streamlit}
\begin{tabular}{L{4.6cm} L{8.4cm}}
\toprule
\textbf{Thành phần} & \textbf{Mô tả} \\
\midrule
Giao diện          & Streamlit, bố cục rộng, thanh điều khiển bên trái \\
Biểu đồ            & Plotly, biểu đồ cột ngang cho phân bố xác suất \\
Nạp mô hình        & Có bộ nhớ đệm, mô hình chỉ nạp một lần cho mỗi phiên làm việc \\
Thiết bị tính toán & Tự phát hiện GPU, chuyển về CPU nếu không có \\
Bảng màu           & Okabe--Ito, an toàn với người mù màu \\
Đầu ra định lượng  & Lớp dự đoán, độ tin cậy, tỉ lệ diện tích tổn thương \\
\bottomrule
\end{tabular}
\end{table}

Việc cho phép người dùng chuyển đổi giữa ba phiên bản mô hình ngay trên giao diện
có giá trị sư phạm: nó cho thấy trực quan điều mà \cref{sec:caveat} phân tích bằng
số liệu --- ba phiên bản cho ra ba kiểu mặt nạ khác nhau về hình dạng và diện tích,
và không thể kết luận phiên bản nào ``đúng'' hơn nếu không có nhãn chuẩn.

\subsection{Hạn chế của ứng dụng}

Ứng dụng chỉ nhằm mục đích minh hoạ, chưa phải sản phẩm triển khai thực tế. Bốn
hạn chế chính:

\begin{itemize}
    \item Chưa hiệu chỉnh sai lệch hệ thống của mô hình --- như đã phân tích ở
          \cref{sec:ablation}, mô hình phiên bản V3 có xu hướng ước lượng thừa
          diện tích tổn thương.
    \item Chỉ xử lý ảnh chụp lá đơn lẻ trên nền tương đối sạch, chưa xử lý được
          ảnh chụp toàn tán.
    \item Tỉ lệ diện tích tổn thương hiển thị là tỉ lệ trên \emph{toàn khung ảnh},
          không phải trên diện tích lá --- hai đại lượng chỉ trùng nhau khi lá
          chiếm trọn khung.
    \item Chưa có cơ chế cảnh báo khi ảnh đầu vào nằm ngoài phân bố dữ liệu huấn
          luyện, chẳng hạn ảnh không phải lá sầu riêng.
\end{itemize}

% !TeX root = ../main.tex
\chapter{HƯỚNG DẪN TÁI TẠO KẾT QUẢ}
\label{app:D}

\section{Cấu trúc kho mã nguồn}

\begin{table}[H]
\centering
\caption{Cấu trúc thư mục chính của kho mã nguồn}
\label{tab:app_repo}
\begin{tabular}{L{5.2cm} L{7.8cm}}
\toprule
\textbf{Đường dẫn} & \textbf{Nội dung} \\
\midrule
\texttt{notebooks/} & Mười \en{notebook} thực nghiệm, đánh số theo thứ tự chạy \\
\texttt{notebooks/\allowbreak data/\allowbreak raw/} & Ảnh gốc, chia sẵn theo tập và theo lớp \\
\texttt{notebooks/\allowbreak data/\allowbreak pseudo\_\allowbreak labels*/} & Các bộ nhãn giả sinh ra \\
\texttt{notebooks/\allowbreak models/} & Điểm kiểm tra và tệp chỉ số của các mô hình \\
\texttt{utils/} & Mã dùng chung: mô hình, tiền xử lý, chỉ số, Grad-CAM \\
\texttt{visualizations/} & Hình sinh ra từ các \en{notebook} \\
\texttt{app.py} & Ứng dụng minh hoạ Streamlit \\
\texttt{docs/thesis/} & Mã nguồn LaTeX của luận văn \\
\bottomrule
\end{tabular}
\end{table}

\section{Thứ tự thực thi}

Các \en{notebook} phải chạy theo đúng thứ tự vì mỗi bước sử dụng sản phẩm của bước
trước.

\begin{table}[H]
\centering
\caption{Thứ tự thực thi và sản phẩm của từng notebook}
\label{tab:app_order}
\small
\begin{tabular}{c L{4.6cm} L{6.4cm}}
\toprule
\textbf{STT} & \textbf{Notebook} & \textbf{Sản phẩm} \\
\midrule
1 & \texttt{01-\allowbreak data-\allowbreak exploration} & Hình EDA, thống kê bộ dữ liệu \\
2 & \texttt{02-\allowbreak classification-\allowbreak baseline} & \texttt{best\_\allowbreak model.\allowbreak pth}, chỉ số phân loại \\
3 & \texttt{03-\allowbreak xai-\allowbreak gradcam-\allowbreak analysis} & Hình bản đồ kích hoạt, thống kê độ tin cậy \\
4 & \texttt{04-\allowbreak pseudo-\allowbreak labeling} & 3.104 mặt nạ nhãn giả V1 \\
5 & \texttt{05-\allowbreak segmentation-\allowbreak model} & Mô hình phân đoạn V1, chỉ số kiểm thử \\
6 & \texttt{06-\allowbreak pseudo-\allowbreak labeling-\allowbreak v2} & 3.104 mặt nạ nhãn giả V2 \\
7 & \texttt{07-\allowbreak segmentation-\allowbreak v2} & Mô hình phân đoạn V2, chỉ số kiểm thử \\
8 & \texttt{08-\allowbreak segmentation-\allowbreak v3-\allowbreak sam} & Nhãn tinh chỉnh SAM, mô hình V3 \\
9 & \texttt{10-\allowbreak inference-\allowbreak single-\allowbreak image} & Hình minh hoạ suy luận đầu cuối \\
\bottomrule
\end{tabular}
\end{table}

Riêng \en{notebook} thứ tám có cơ chế tiếp tục: nếu thư mục nhãn tinh chỉnh đã đủ
3.104 tệp, bước sinh nhãn bằng \SAM{} sẽ được bỏ qua và mô hình \SAM{} không được
nạp lên bộ nhớ đồ hoạ. Muốn sinh lại từ đầu cần đặt biến \texttt{FORCE\_REDO} bằng
\texttt{True} hoặc xoá thư mục nhãn.

Các \en{notebook} huấn luyện đều lưu điểm kiểm tra tạm sau mỗi epoch, cho phép
tiếp tục nếu quá trình bị gián đoạn.

\section{Yêu cầu môi trường}

Môi trường đã dùng để sinh toàn bộ kết quả trong luận văn được liệt kê ở
\cref{tab:moitruong}. Các thư viện chính cần cài đặt:

\begin{itemize}
    \item \texttt{torch}, \texttt{torchvision} --- phiên bản tương thích CUDA của
          máy đích.
    \item \texttt{segmentation-\allowbreak models-\allowbreak pytorch} --- cung cấp kiến trúc \UNetpp{}.
    \item \texttt{albumentations} --- tăng cường dữ liệu đồng bộ ảnh và mặt nạ.
    \item \texttt{segment-\allowbreak anything} --- mô hình \SAM{}.
    \item \texttt{opencv-\allowbreak python}, \texttt{pillow}, \texttt{numpy},
          \texttt{pandas}, \texttt{matplotlib}, \texttt{seaborn},
          \texttt{scikit-learn}, \texttt{tqdm}.
    \item \texttt{streamlit}, \texttt{plotly} --- chỉ cần cho ứng dụng minh hoạ.
\end{itemize}

Trọng số \SAM{} ViT-B (khoảng 375 MB) được tải tự động ở lần chạy đầu tiên của
\en{notebook} thứ tám.

\section{Lưu ý về khả năng tái lập}

\subsection{Các yếu tố đã được cố định}

Hạt giống ngẫu nhiên được đặt bằng 42 cho việc chia tập dữ liệu ở mọi phiên bản
phân đoạn, bảo đảm các phiên bản V2 và V3 dùng đúng cùng một tập kiểm thử.

\subsection{Các yếu tố chưa được cố định hoàn toàn}

Cần lưu ý ba nguồn sai khác có thể khiến kết quả chạy lại không trùng khớp tuyệt
đối với số liệu báo cáo:

\begin{itemize}
    \item \textbf{Tính không tất định của cuDNN.} Chỉ \en{notebook} thứ tám đặt
          \texttt{cudnn.\allowbreak deterministic} bằng \texttt{True}; các \en{notebook} còn
          lại không đặt. Các thuật toán tích chập được chọn tự động có thể khác
          nhau giữa các lần chạy và giữa các loại GPU.
    \item \textbf{Lấy mẫu điểm gợi ý cho \SAM{}.} Bốn điểm tiền cảnh và ba điểm
          nền được lấy ngẫu nhiên mà không cố định hạt giống, nên nhãn tinh chỉnh
          sinh lại sẽ khác đôi chút so với bộ nhãn đã dùng.
    \item \textbf{Tăng cường dữ liệu ngẫu nhiên.} Chuỗi phép biến đổi Albumentations
          không cố định hạt giống trong vòng lặp huấn luyện.
    \end{itemize}

Vì vậy, khi chạy lại, các chỉ số có thể lệch ở chữ số thập phân thứ hai hoặc thứ
ba. Các kết luận định tính của luận văn --- chiều của các so sánh, hiện tượng nở
mặt nạ, sự lệch pha giữa hàm mất mát và IoU --- không phụ thuộc vào mức sai khác
này.

\subsection{Khuyến nghị cho lần chạy tiếp theo}

Để tăng khả năng tái lập, ba thay đổi nên được áp dụng: cố định hạt giống ngẫu
nhiên cho \texttt{random}, \texttt{numpy} và \texttt{torch} ở đầu mọi
\en{notebook}; đặt \texttt{cudnn.\allowbreak deterministic} bằng \texttt{True} một cách nhất
quán; và cố định hạt giống cho bước lấy mẫu điểm gợi ý của \SAM{}.

\section{Truy vết số liệu}

Toàn bộ con số xuất hiện trong luận văn được ghi lại trong tệp
\texttt{docs/\allowbreak thesis/\allowbreak INVENTORY.\allowbreak md}, kèm thông tin về \en{notebook} và ô lệnh đã
sinh ra nó. Tệp này cũng liệt kê các điểm mà số liệu chưa đủ bằng chứng, cùng
những chỗ số liệu trong công bố liên quan của nhóm không khớp với kết quả thực
thi thực tế.

Một số giá trị trong luận văn được tính lại trực tiếp từ các tệp kết quả trên đĩa
thay vì lấy từ ô kết quả của \en{notebook} --- cụ thể là bảng độ phủ của nhãn tinh
chỉnh bằng \SAM{} ở \cref{tab:sam_expansion}. Đoạn mã dùng để tính lại được ghi
đầy đủ trong tệp kiểm kê nói trên, cho phép kiểm chứng độc lập.


\end{document}
